\documentclass[a4paper,12pt]{book}
\usepackage[utf8]{inputenc}
\usepackage[T2A]{fontenc}
\usepackage[russian]{babel}
\usepackage{graphicx}
\usepackage{geometry}
\geometry{left=2cm,right=2cm,top=2cm,bottom=2cm}
\usepackage{array}
\usepackage{multirow}
\usepackage{booktabs}
\usepackage{longtable}
\usepackage{hyperref}
\usepackage{caption}
\usepackage{float}
\usepackage{textcomp}
\usepackage{amsmath}
\usepackage{amssymb}
\setkeys{Gin}{draft}
\DeclareUnicodeCharacter{2022}{\textbullet}
\DeclareUnicodeCharacter{00B0}{\textdegree}
\DeclareUnicodeCharacter{00D7}{\ensuremath{\times}}
\DeclareUnicodeCharacter{2248}{\ensuremath{\approx}}
\floatplacement{figure}{H}
\floatplacement{table}{H}

\begin{document}

% Титульный лист
\thispagestyle{empty}
\begin{center}
\vspace*{5cm}

{\Huge\bfseries Клинические случаи\par
в гематологии, биохимии\par
и анализе газов крови\par
у мелких животных\par}

\vspace{2cm}
{\Large Рафаэль А. Перес Эсиха\par}

\vfill
\includegraphics[width=0.4\textwidth]{placeholder.jpg} % Placeholder for the image
\vspace{1cm}

{\large Grupo Asis Biomedia, SL\par
2018}
\end{center}

\newpage
\thispagestyle{empty}
\begin{center}
\includegraphics[width=0.8\textwidth]{placeholder.jpg}
\end{center}

\newpage
~
\newpage

% Выходные данные
\thispagestyle{empty}
\begin{flushleft}
2018 Grupo Asis Biomedia, SL\\
Plaza Antonio Beltran Martinez n.° 1, planta 8 - letra I\\
(Centro empresarial El Trovador)\\
50002 Zaragoza -- Spain

Первое издание: 2018

Эта книга была первоначально издана на испанском языке под названием:\\
\textit{Casos clínicos de hematología, bioquímica y gasometría de pequeños animales}\\
© 2018 Grupo Asis Biomedia, SL\\
ISBN испанского издания: 978-84-17640-26-2

Перевод: Sanscrit

Дизайн, верстка и печать: Servet editorial -- Grupo Asis Biomedia, SL\\
www.grupoasis.com\\
info@grupoasis.com

\includegraphics[width=0.3\textwidth]{placeholder.jpg}

Servet -- издательство Grupo Asis

Все права защищены.

Любое воспроизведение, распространение, публикация или преобразование данной книги разрешается только с согласия правообладателей, за исключением случаев, предусмотренных законом. Для получения разрешения на фотокопирование или сканирование любой части книги обращайтесь в CEDRO (Испанский центр репродуктивных прав, www.cedro.org) или на сайт www.conlicencia.com; тел. +34 91 702 19 70 / +34 93 272 04 47).

Предупреждение:

Ветеринарная наука постоянно развивается, как и фармакология и другие науки. Поэтому ветеринарный врач обязан определять и проверять дозировку, способ применения, продолжительность лечения и любые возможные противопоказания для каждого конкретного пациента на основе своего профессионального опыта. Ни издатель, ни автор не несут ответственности за ущерб или вред, причиненный людям, животным или имуществу в результате правильного или неправильного применения информации, содержащейся в данной книге.
\end{flushleft}

\newpage
\thispagestyle{empty}
\begin{center}
{\Large Клинические случаи\par
гематологии,\par
биохимии и анализа газов крови\par
у мелких животных}

\vspace{2cm}
Рафаэль А. Перес Эсиха
\end{center}

\newpage
~
\newpage

% Благодарности
\chapter*{Благодарности}
Всему персоналу Лабораторной службы Университетской ветеринарной клиники Кордовы (Mª. Carmen González, Jennifer Parejo, Antonia Castillo и Belén Díez), а также ветеринарным врачам Службы внутренней медицины (особенно д-рам Carmen Martínez и Carles Mengual).

\newpage
~
\newpage

% Об авторе
\chapter*{Автор}
\textbf{Алехандро Перес-Эсиха}

Автор имеет степень доктора ветеринарной медицины и степень магистра в области здоровья животных, полученную в 2009 году в Университете Кордовы. В 2015 году д-р Перес-Эсиха получил Европейский диплом по ветеринарной патологии Европейского колледжа ветеринарных патологов (ECVP). В настоящее время он возглавляет Лабораторную службу Университетской ветеринарной клиники Кордовы и является клиническим ветеринаром в ее Службе гематологии и онкологии. Его основные научные интересы включают ветеринарную медицину, гематологию, иммунопатологию и клиническую патологию мелких животных и лошадей. Он является членом Европейского колледжа ветеринарной патологии, Испанского общества ветеринарной внутренней медицины, Испанского общества ветеринарной патологии и Европейского общества ветеринарной клинической патологии.

Д-р Перес-Эсиха участвовал во многих национальных и европейских исследовательских проектах, является автором ряда публикаций в рецензируемых научных журналах, посвященных ветеринарной гематологии и патологии, а также соавтором нескольких научно-популярных книг по клинической патологии и гематологии, эндокринологии, гистопатологии, цитологии и взятию проб. Он также выступал с докладами на многочисленных конференциях и курсах, связанных с ветеринарной внутренней медициной.

\newpage
%
\chapter*{Пролог}
Лабораторная диагностика является одним из самых полезных инструментов для клинических ветеринаров и одной из областей, в которых специалисты испытывают наибольшие трудности в интерпретации. Эта проблема отчасти вызвана чрезмерно жестким и статичным подходом к данной области, а также нехваткой инструментов, которые бы четко направляли эту работу в повседневной клинической практике. В связи с последним, данный сборник клинических случаев, несомненно, станет прикладным справочным источником на английском языке.

Текст «Клинические случаи мелких животных, включающие гематологию, биохимию и газометрию» Алехандро Переса-Эсихи тщательно и исчерпывающе рассматривает наиболее распространенные лабораторные изменения у собак и кошек, предоставляя превосходные иллюстрации и обеспечивая максимально возможную дидактическую и практическую ориентацию для каждой клинической ситуации. Благодаря такому подходу, это руководство углубляется в патофизиогенез, клинические симптомы и диагностику многочисленных патологических процессов в увлекательной и наглядной форме, закладывая основы для адекватной и полезной интерпретации лабораторных результатов.

Этот текст предназначен для всех клинических ветеринаров, независимо от их опыта в лабораторной диагностике, как для ознакомления с областями гематологии, биохимии и газометрии, так и для углубленного анализа сложных клинических предположений, включающих множество взаимосвязанных факторов. Благодаря выдающейся работе автора, его обширному клиническому и патологическому опыту, а также исключительной педагогической и академической истории, читатель легко и с удовольствием освоит лабораторную диагностику и узнает, как эти методы можно использовать в повседневной клинической работе.

\begin{flushright}
Д-р Франсиско Хавьер Мендоса Гарсия,\\
DVM, MSc, PhD, Diplomate ECEIM.
\end{flushright}

\newpage
~
\chapter*{Предисловие}
Студентов-ветеринаров всегда учат разделять знания по разным областям, таким как нефрология, пищеварительная система, хирургия, онкология и т.д. Только начав практиковать, мы осознаем серьезность этой ошибки: наши пациенты не приходят с табличкой на шее, указывающей, к какой специальности относится их проблема.

Каждый клинический случай — это свой собственный мир и свой хаос; прекрасный конгломерат изменений и взаимодействий, который никоим образом не напоминает жесткие правила, о которых мы читаем в теоретических руководствах. Конечно, ветеринары должны знать теорию, но гораздо важнее уметь глобально интерпретировать все данные и адекватно применять их к каждому клиническому случаю. Почему возникают эти изменения? Почему повышаются одни показатели, а другие нет? Что следует подозревать при одновременном учете симптомов и результатов различных тестов?

Это руководство по клиническим случаям, включающим гематологию, биохимию или газометрию, предоставляет клиницистам патофизиологические основы, объясняющие наиболее распространенные изменения на лабораторном уровне у мелких животных, всегда фокусируясь на реальности самым дидактичным и легким для понимания способом.

Каждый случай представлен как головоломка с собственной историей болезни и результатами физикального обследования, лабораторных и других дополнительных методов (например, рентгенографии, УЗИ или цитологии и т.д.). Задача состоит в том, чтобы понять, как каждый из этих тестов дополняет и объясняет друг друга, и получить на их основе надежный и точный диагноз.

После решения каждый случай сохраняет образовательную ценность, поскольку все они используются для предоставления диагностических алгоритмов, списков дифференциальных диагнозов, формул для расчета переменных и прикладных референсных диапазонов. Кроме того, учитывая, что лабораторная диагностика является столь наглядной областью, этот текст предоставляет огромную галерею диагностических изображений, которые послужат примером и помогут в обучении читателя для будущей лабораторной работы.

Моя конечная цель — чтобы клиницисты научились применять глобальный подход к каждому клиническому случаю, поняли, как максимально использовать свои лабораторные результаты, и осознали потенциал тестов, доступных им на кончиках пальцев. Искренне надеюсь, что я достиг своей цели.

\begin{flushright}
С уважением,\\
Алехандро Перес-Эсиха
\end{flushright}

\newpage
\setcounter{page}{1}
\tableofcontents

\chapter{Сокращения}
\label{chap:abbrev}
\begin{longtable}{p{4cm} p{10cm}}
\textbf{Сокращение} & \textbf{Значение} \\
\hline
\endhead
$\mu$g & Микрограммы \\
$\mu$L & Микролитр \\
A:G & Альбумин-глобулиновое соотношение \\
admin & Введение \\
AG & Анионный интервал \\
Al & Алюминий \\
ALB & Альбумин \\
ALT & Аланинаминотрансфераза \\
AML & Острый миелоидный лейкоз \\
AST & Аспартатаминотрансфераза \\
ATOT & Общие слабые нелетучие кислоты \\
BAN & Палочкоядерный нейтрофил \\
BASO & Базофил \\
Ca & Кальций \\
ARC & Абсолютное количество ретикулоцитов \\
MCHC & Средняя концентрация гемоглобина в эритроците \\
DIC & Диссеминированное внутрисосудистое свертывание \\
CK & Креатинфосфокиназа \\
Cl & Хлор \\
MCH & Среднее содержание гемоглобина в эритроците \\
CPK & Креатинфосфокиназа \\
cPLI & Иммунореактивность панкреатической липазы собак \\
dL & Децилитр \\
DM & Сахарный диабет \\
Dx & Медицинский диагноз \\
EDTA & Этилендиаминтетрауксусная кислота \\
ELISA & Иммуноферментный анализ \\
EOS & Эозинофил \\
AF & Щелочная фосфатаза \\
HR & Частота сердечных сокращений \\
Fe & Железо \\
FeLV & Вирус лейкемии кошек \\
FIV & Вирус иммунодефицита кошек \\
fL & Фемтолитр \\
Frag. & Фрагменты \\
g & Грамм \\
WBC & Лейкоциты \\
SGOT & Глутаминовая щавелевоуксусная трансаминаза (AST) \\
SGPT & Глутамат-пируват-трансаминаза (ALT) \\
RBC & Эритроциты \\
h & Часы \\
HCO3- & Бикарбонат \\
HCT & Гематокрит (расчетный) \\
HDL & Липопротеины высокой плотности \\
Hb & Концентрация гемоглобина \\
HF & Сердечная недостаточность \\
Ig & Иммуноглобулин \\
EPI & Экзокринная недостаточность поджелудочной железы \\
RPI & Индекс продукции ретикулоцитов \\
ESRD & Терминальная стадия почечной недостаточности \\
ARF & Острая почечная недостаточность \\
CRF & Хроническая почечная недостаточность \\
IRIS & Международное общество по изучению почек \\
UTI & Инфекция мочевыводящих путей \\
IV & Внутривенно \\
k & 1,000 \\
K & Калий \\
L & Литр \\
LDH & Лактатдегидрогеназа \\
LDL & Липопротеины низкой плотности \\
bpm & Ударов в минуту \\
LYM & Лимфоцит \\
M & 1,000,000 \\
mEq & Миллиэквиваленты \\
min. & Минуты \\
mL & Миллилитр \\
mmHg & Миллиметры ртутного столба \\
mmol & Миллимоль \\
MONO & Моноцит \\
MPV & Средний объем тромбоцита \\
N & Норма \\
Na & Натрий \\
NMB & Новый метиленовый синий \\
NEU & Нейтрофил \\
ng & Нанограмм \\
NRBC & Ядерные эритроциты \\
°C & Градусы Цельсия \\
P & Фосфор \\
PaCO2 & Артериальное парциальное давление углекислого газа \\
PaO2 & Артериальное парциальное давление кислорода \\
PCO2 & Парциальное давление углекислого газа \\
CRP & Скорректированный процент ретикулоцитов \\
PCR & Полимеразная цепная реакция \\
APP & Белок острой фазы \\
pg & Пикограмм \\
FIP & Инфекционный перитонит кошек \\
PLI & Иммунореактивность панкреатической липазы \\
PLT & Тромбоцит \\
PO2 & Парциальное давление кислорода \\
TP & Общий белок \\
RBC & Концентрация эритроцитов \\
RDW & Ширина распределения эритроцитов \\
Ref. & Референсный \\
RETIC & Ретикулоциты \\
RIA & Радиоиммунный анализ \\
bpm & Вдохов в минуту \\
SDMA & Симметричный диметиларгинин \\
sec. & Секунды \\
SIBO & Синдром избыточного бактериального роста в тонком кишечнике \\
SID & Разница сильных ионов \\
spp. & Виды \\
CT & Компьютерная томография \\
TCO2 & Общий диоксид углерода \\
TLI & Трипсиноподобная иммунореактивность \\
PT & Протромбиновое время \\
CRT & Время наполнения капилляров \\
BMBT & Время кровотечения слизистой оболочки щеки \\
aPTT & Активированное частичное тромбопластиновое время \\
IU & Международные единицы \\
UPC & Соотношение белок-креатинин в моче \\
MCV & Средний корпускулярный объем \\
VLDL & Липопротеины очень низкой плотности \\
WBC & Концентрация лейкоцитов \\
× & Умножить (×10) \\
$\gamma$GT & Гамма-глутамилтрансфераза \\
\end{longtable}

\chapter{Гематологические клинические случаи}
\label{chap:haem}

\section{Случай 1. Щенок с диареей}

\subsection*{Клиническая история и физикальное обследование}
Целый 5-недельный кобель веймаранера поступил в клинику из-за апатии, потери аппетита и перемежающейся диареи с обильной слизью (рис. 1). Животное было вакцинировано несколько дней назад (против парвовируса и чумы собак), но не было дегельминтизировано. При физикальном обследовании выявлены сухие слизистые оболочки и легкое снижение тургора кожи с клинической оценкой обезвоживания в $7\%$.

\begin{tabular}{@{}ll@{}}
\toprule
Ректальная температура & 39.4°C \\
Частота сердечных сокращений & 100 уд/мин \\
Частота дыхания & 30 дых/мин \\
Время наполнения капилляров & Менее 2 секунд \\
\bottomrule
\end{tabular}

\begin{figure}[htbp]
\centering
\includegraphics[width=0.5\textwidth]{placeholder.jpg}
\caption{Пациент при поступлении.}
\end{figure}

\subsection*{Гемограмма}
\begin{tabular}{lccccc}
\toprule
Тест & Результат & Референсный диапазон & LOW & NORMAL & HIGH \\
\midrule
HCT & 36\% & 37--55 & $\blacksquare$ \\
RBC & 4.6 × 10\textsuperscript{12}/L & 5.50--8.5 & $\blacksquare$ \\
Hb & 111 g/L & 120--180 & $\blacksquare$ \\
MCV & 79 fL & 60--77 & $\blacksquare$ \\
RDW & 16\% & 12.0--15.0 & $\blacksquare$ \\
MCHC & 310 g/L & 310--340 & $\blacksquare$ \\
MCH & 24 pg & 19.5--24.5 & $\blacksquare$ \\
WBC & 18.5 × 10\textsuperscript{9}/L & 6.0--17.0 & $\blacksquare$ \\
NEU & 11.2 × 10\textsuperscript{9}/L & 3.0--12.0 & $\blacksquare$ \\
LYM & 5.22 × 10\textsuperscript{9}/L & 1--4.80 & $\blacksquare$ \\
MONO & 1.17 × 10\textsuperscript{9}/L & 0.2--1.50 & $\blacksquare$ \\
EOS & 0.55 × 10\textsuperscript{9}/L & 0--0.8 & $\blacksquare$ \\
BASO & 0 × 10\textsuperscript{9}/L & 0--0.4 & $\blacksquare$ \\
PLT & 350 × 10\textsuperscript{9}/L & 200--500 & $\blacksquare$ \\
MPV & 9.02 fL & 3.90--11.1 & $\blacksquare$ \\
\% RETIC* & 5.5\% & 0.5--1.5 & $\blacksquare$ \\
\% BAND* & 0.47\% & 0--3.5 & $\blacksquare$ \\
\bottomrule
\multicolumn{6}{l}{*Ручной подсчет}
\end{tabular}

\subsection*{Вопрос: Какие изменения можно наблюдать в гемограмме этого животного?}

\subsection*{Мазок крови}
\begin{figure}[htbp]
\centering
\includegraphics[width=0.6\textwidth]{placeholder.jpg}
\caption{Микроскопическое изображение мазка крови животного. Окраска Diff-Quik, 400x.}
\end{figure}
\begin{figure}[htbp]
\centering
\includegraphics[width=0.6\textwidth]{placeholder.jpg}
\caption{Деталь мазка крови. Окраска Diff-Quik, 1000x.}
\end{figure}
\begin{figure}[htbp]
\centering
\includegraphics[width=0.6\textwidth]{placeholder.jpg}
\caption{Различные детали мазка крови. Окраска Diff-Quik, 1000x.}
\end{figure}
\begin{figure}[htbp]
\centering
\includegraphics[width=0.6\textwidth]{placeholder.jpg}
\caption{Деталь мазка крови. Окраска Diff-Quik, 1000x.}
\end{figure}

\subsection*{Вопрос: Какие изменения наблюдались в мазке крови этого животного?}

\subsection*{Исследование кала}
\begin{figure}[htbp]
\centering
\includegraphics[width=0.5\textwidth]{placeholder.jpg}
\caption{Изображение пробы кала после флотации с сульфатом цинка.}
\end{figure}

\subsection*{Вопрос: Учитывая данные кала и гематологические находки, каков диагноз?}

\subsection*{Разбор случая}
\subsubsection*{Какие изменения можно наблюдать в гемограмме этого животного?}
При оценке любого щенка (особенно в возрасте до 2 месяцев) необходимо всегда использовать возраст-специфичные гематологические диапазоны (рис. 7). Если используются взрослые диапазоны, как в этом клиническом случае, животное будет казаться анемичным (снижение HCT, RBC и Hb) с регенераторными признаками (повышенный процент ретикулоцитов, увеличение MCV и RDW). С учетом возраста пациента единственным изменением была легкая эритроцитоз (гемоконцентрация в результате обезвоживания). Аналогично, у здоровых щенков в возрасте до 2 месяцев процент ретикулоцитов колеблется от 6 до $3\%$ (и снижается с возрастом), при этом характерна относительная макроцитоз (высокий MCV; рис. 7) и анизоцитоз (высокий RDW), вызванный постепенным созреванием эритроцитов после рождения. В лейкограмме обнаружен легкий лейкоцитоз за счет лимфоцитоза, что может быть информативной находкой, так как лейкоцитарные и тромбоцитарные параметры у взрослых и щенков сходны.

\begin{figure}[htbp]
\centering
\includegraphics[width=0.9\textwidth]{placeholder.jpg}
\caption{Эволюция основных эритроцитарных показателей у щенков в возрасте до 2 месяцев (среднее значение у взрослых собак показано синей линией на каждом графике).}
\end{figure}

\subsubsection*{Какие изменения наблюдались в мазке крови этого животного?}
В мазке крови обнаружено множество ядерных эритроцитов (нормобластоз) (рис. 8). Они характеризуются наличием ядра с очень компактным хроматином и цитоплазмой с зубчатыми краями с вариабельной окраской, в зависимости от степени созревания (более ацидофильная у зрелых метарубрицитов и базофильная у более незрелых форм). Хотя этот признак может появляться при повреждении костного мозга или при регенераторных анемиях (наряду с другими признаками ответа костного мозга), эти клетки также очень многочисленны в период эритроидного созревания в первые месяцы жизни, при этом количество до 1000 ядерных эритроцитов/мкл считается нормальным у щенков (снижаясь от рождения до первых 2-3 месяцев). Крайне важно проверить наличие этих клеток и подсчитать их вручную, отличая от лимфоцитов, у которых хроматин более дисперсный (рис. 8).

Присутствие ядерных эритроцитов часто мешает автоматическому подсчету лейкоцитов (поскольку счетное оборудование обычно принимает их за лимфоциты), поэтому полученные значения необходимо скорректировать при их большом количестве (рис. 9).

\[
\text{Скорректированные WBC} = \frac{\text{Полученные WBC} \times 100}{\text{NRBC подсчитанные} + 100}
\]

\begin{figure}[htbp]
\centering
\includegraphics[width=0.6\textwidth]{placeholder.jpg}
\caption{Формула для коррекции лейкоцитарного счета в образцах с многочисленными ядерными эритроцитами. NRBC подсчитанные: ядерные эритроциты на 100 лейкоцитов при ручном подсчете.}
\end{figure}

Также наблюдалось множество эхиноцитов (один из наиболее распространенных артефактов в гематологических мазках) (рис. 3). Эти клетки характеризуются мелкими, равномерными выступами, равномерно распределенными по поверхности эритроцитов. Обычно эхиноцитоз затрагивает большинство эритроцитов в мазке. Однако эта находка может также встречаться у пациентов с почечной патологией, гломерулопатиями или после приема некоторых лекарств. Крайне важно не путать эти клетки с акантоцитами (которые показывают нерегулярные и асимметрично распределенные выступы на поверхности).

\paragraph{На заметку}
Наличие эхиноцитов обычно является артефактной находкой, не имеющей патологического значения, хотя важно отличать их от акантоцитов.

\paragraph{Основные причины эхиноцитоза эритроцитов}
Слишком длительное хранение образца; недостаточное прогревание образца после охлаждения; избыток EDTA (пробирка недостаточно заполнена кровью); или забор пробы из катетера при проведении высокоскоростной инфузионной терапии без предварительной промывки.

На лейкоцитарном уровне отмечалось наличие моноцитов с вакуолизированной цитоплазмой (рис. 4), изменение, обычно не имеющее патологического значения, возникающее из-за фагоцитоза EDTA и других элементов при чрезмерно длительном хранении образца. Также наблюдались лимфоциты с несколько увеличенной цитоплазмой и голубоватым окрашиванием с легкой зоной перинуклеарного просветления (рис. 5). Эти клетки соответствуют реактивным иммуноцитам или лимфоцитам и часто встречаются после вакцинации или при выраженной антигенной стимуляции.

\subsubsection*{Учитывая данные кала и гематологические находки, каков диагноз?}
После учета данных пациента с учетом возрастных референсных диапазонов мы видим, что единственным изменением была легкая гемоконцентрация из-за обезвоживания, без реальных гематологических нарушений. Легкий лейкоцитоз и лимфоцитоз, показанные автоматическим счетчиком у этого животного, были связаны с ошибкой классификации ядерных эритроцитов анализатором (реальный лейкоцитарный счет составлял $15.5 \times 10^{9} / \mathrm{L}$). Кроме того, все находки на уровне мазка были либо артефактными, либо нормальными для недавно вакцинированного пациента. Паразиты, обнаруженные в кале у этого пациента (рис. 6), были идентифицированы как \textit{Giardia} spp. и признаны ответственными за представленную клиническую картину.

\section{Случай 2. Плановый осмотр борзой}

\subsection*{Клиническая история и физикальное обследование}
Некастрированный 1-летний кобель грейхаунда поступил на плановый осмотр. Животное еще не было вакцинировано или дегельминтизировано, владельцы отметили лишь плохую консистенцию кала (рис. 1). Физикальное обследование не выявило отклонений, хотя слизистые оболочки казались слегка бледными.

\begin{tabular}{@{}ll@{}}
\toprule
Ректальная температура & 38.4°C \\
ЧСС & 90 уд/мин \\
ЧД & 20 дых/мин \\
ВНК & Менее 2 сек \\
\bottomrule
\end{tabular}

\begin{figure}[htbp]
\centering
\includegraphics[width=0.5\textwidth]{placeholder.jpg}
\caption{Пациент при поступлении.}
\end{figure}

\subsection*{Гемограмма}
\begin{tabular}{lccccc}
\toprule
Тест & Результат & Реф. диапазон & LOW & NORMAL & HIGH \\
\midrule
HCT & 41\% & 37--55 & $\blacksquare$ \\
RBC & 5.86 × 10\textsuperscript{12}/L & 5.50--8.5 & $\blacksquare$ \\
Hb & 121 g/L & 120--180 & $\blacksquare$ \\
MCV & 78 fL & 60--77 & $\blacksquare$ \\
RDW & 17\% & 12.0--15.0 & $\blacksquare$ \\
MCHC & 300 g/L & 310--340 & $\blacksquare$ \\
MCH & 21 pg & 19.5--24.5 & $\blacksquare$ \\
WBC & 5.5 × 10\textsuperscript{9}/L & 6.0--17.0 & $\blacksquare$ \\
NEU & 3.0 × 10\textsuperscript{9}/L & 3.0--12.0 & $\blacksquare$ \\
LYM & 1.02 × 10\textsuperscript{9}/L & 1--4.80 & $\blacksquare$ \\
MONO & 2.47 × 10\textsuperscript{9}/L & 0.2--1.50 & $\blacksquare$ \\
EOS & 0 × 10\textsuperscript{9}/L & 0.0--0.8 & $\blacksquare$ \\
BASO & 0 × 10\textsuperscript{9}/L & 0--0.4 & $\blacksquare$ \\
PLT & 90 × 10\textsuperscript{9}/L & 200--500 & $\blacksquare$ \\
MPV & 9.22 fL & 3.90--11.1 & $\blacksquare$ \\
\% RETIC* & 4\% & 0.5--1.5 & $\blacksquare$ \\
\% BAND* & 1.4\% & 0--3.5 & $\blacksquare$ \\
\bottomrule
\multicolumn{6}{l}{*Ручной подсчет}
\end{tabular}

\subsection*{Вопрос: Какие изменения можно наблюдать в гемограмме этого животного?}

\subsection*{Мазок крови}
\begin{figure}[htbp]
\centering
\includegraphics[width=0.6\textwidth]{placeholder.jpg}
\caption{Микроскопическое изображение мазка крови. Окраска Diff-Quik, 400x.}
\end{figure}
\begin{figure}[htbp]
\centering
\includegraphics[width=0.6\textwidth]{placeholder.jpg}
\caption{Деталь мазка крови. Окраска Diff-Quik, 1000x.}
\end{figure}
\begin{figure}[htbp]
\centering
\includegraphics[width=0.6\textwidth]{placeholder.jpg}
\caption{1000x.}
\end{figure}

\subsection*{Исследование кала}
\begin{figure}[htbp]
\centering
\includegraphics[width=0.5\textwidth]{placeholder.jpg}
\caption{Изображение пробы кала после флотации с сульфатом цинка.}
\end{figure}

\subsection*{Вопрос: Учитывая данные кала и гематологические находки, каков диагноз?}

\subsection*{Разбор случая}
\textbf{Геморрагическая анемия вследствие кишечного паразитоза}

\subsubsection*{Какие изменения можно наблюдать в гемограмме этого животного?}
Чтобы избежать путаницы при использовании общих диапазонов нормы для большинства собак (табл. 1), важно учитывать значительные гематологические особенности грейхаундов. Значения HCT, RBC и гемоглобина у этого пациента были снижены для грейхаунда, что указывает на анемию. Высокий процент ретикулоцитов вместе с увеличением MCV, RDW и снижением MCHC были маркерами регенераторной анемии (табл. 2). Общее количество лейкоцитов и тромбоцитов снова находилось в пределах нормы для этой породы (табл. 1), хотя, по данным анализатора, наблюдался моноцитоз и эозинопения.

\begin{table}[htbp]
\centering
\caption{Гематологические референсные диапазоны, специфичные для грейхаундов.}
\begin{tabular}{lcc}
\toprule
Параметр (единицы) & Референсный диапазон для грейхаундов \\
\midrule
HCT (\%) & 55--65 \\
RBC (×10\textsuperscript{12}/L) & 7.4--9 \\
Hb (g/L) & 190--215 \\
WBC (×10\textsuperscript{9}/L) & 3.5--6.5 \\
PLT (×10\textsuperscript{9}/L) & 80--200 \\
\bottomrule
\end{tabular}
\end{table}

\begin{table}[htbp]
\centering
\caption{Признаки, указывающие на регенерацию у анемичных пациентов.}
\begin{tabular}{p{3cm}p{4cm}p{5cm}}
\toprule
Техника & Признак регенерации & Другие причины/соображения \\
\midrule
\multirow{4}{*}{Автоматические анализаторы} & ↑ количество (или \%) ретикулоцитов & Общее количество ретикулоцитов более информативно, чем процент (референс: < 70,000/мкл или > 150,000/мкл для регенераторных клинических картин). \\
& ↑ MCV & Миелодисплазия, инфекция FeLV или врожденные состояния у чихуахуа и миниатюрных пород. \\
& ↑ RDW & Дефицит железа (исключить путем оценки MCV). \\
& ↓ MCHC & Дефицит железа (исключить путем оценки MCV). \\
\multirow{5}{*}{Мазки крови} & Полихроматофилы & --- \\
& Тельца Хауэлла-Жолли & Вовлечение селезенки или спленэктомия. \\
& Ядерные эритроциты & Если не сопровождаются другими регенераторными признаками, указывают на повреждение костного мозга. Относительно часто у новорожденных (возраст менее 1 месяца). \\
& Макроцитоз & Миелодисплазия, инфекция FeLV или врожденные состояния у чихуахуа и миниатюрных пород. \\
& Анизоцитоз и полихромазия & Дефицит железа или морфологические изменения эритроидных клеток. \\
\bottomrule
\end{tabular}
\end{table}

\subsubsection*{Какие изменения наблюдались в мазке крови этого животного?}
В мазке крови обнаружены многочисленные полихроматофилы (рис. 6), которые обычно немного крупнее других эритроцитов (макроцитоз или анизоцитоз), имеют более голубоватое окрашивание (полихромазия) и отсутствие узнаваемого ореола. Помимо полихроматофилов, на детали мазка (рис. 3) видны три тельца Хауэлла-Жолли и один ядерный эритроцит. Все эти находки характерны для регенераторных анемий (табл. 2). На уровне лейкоцитов (рис. 4) обнаружена выраженная цитоплазматическая вакуолизация эозинофилов с потерей типичных гранул (сравните цитоплазму двулопастной клетки справа с цитоплазмой нейтрофила слева).

\paragraph{На заметку}
Эозинопения обычно появляется при использовании избытка кортикостероидов (гиперадренокортицизм или другие виды лечения и т.д.) или вследствие острого стресса (катехоламины). Эозинофилия (и базофилия) является гораздо более частой находкой.

Эта находка очень распространена у грейхаундов и часто мешает дифференцировке лейкоцитов автоматическими анализаторами (которые иногда ошибочно классифицируют эти эозинофилы как моноциты или нейтрофилы). Типичная морфология эозинофилов и базофилов у мелких животных показана на рис. 7.

\begin{figure}[htbp]
\centering
\includegraphics[width=0.6\textwidth]{placeholder.jpg}
\caption{Типичный мазок крови при регенераторной анемии с наличием многочисленных полихроматофилов (стрелки).}
\end{figure}

\begin{table}[htbp]
\centering
\caption{Основные причины эозинофилии и базофилии.}
\begin{tabular}{p{14cm}}
\toprule
Основные причины эозинофилии и базофилии \\
\midrule
Паразитозы, гиперчувствительность, эозинофильный синдром, мастоцитомы или эозинофильные лейкемии и т.д. \\
\bottomrule
\end{tabular}
\end{table}

\begin{figure}[htbp]
\centering
\includegraphics[width=0.8\textwidth]{placeholder.jpg}
\caption{Эозинофилы и базофилы у мелких животных, (a) эозинофил собаки (несколько возможных морфологий гранул); (b) эозинофил кошки; (c) базофил собаки; (d) базофил кошки.}
\end{figure}

\subsubsection*{Учитывая данные кала и гематологические находки, каков диагноз?}
Поскольку у животного не было признаков гемолиза или повышенного гемоглобина в мазке, клиническая картина, скорее всего, представляла собой геморрагическую анемию с регенерацией (подострое течение). Основные причины и эволюция регенерации при геморрагической анемии показаны в табл. 4.

Паразиты, обнаруженные в кале у этого пациента (рис. 5), были идентифицированы как \textit{Strongyloides stercoralis}, и, учитывая высокую паразитарную нагрузку, был сделан вывод, что наиболее вероятной причиной анемической клинической картины является кишечная паразитарная инвазия.

\begin{table}[htbp]
\centering
\caption{Ключевые моменты, которые следует учитывать при геморрагической анемии.}
\begin{tabular}{p{7cm}p{7cm}}
\toprule
\multicolumn{2}{c}{Наиболее частые причины геморрагической анемии у мелких животных} \\
\midrule
• Травмы (несчастные случаи, ранения или операции и т.д.) & • Коагулопатии (отравление родентицидами, тромбоцитопения или тромбоцитопатии и т.д.) \\
• Желудочно-кишечные кровотечения (язвы, новообразования, паразитозы или парвовирус и т.д.) & • Диссеминированное внутрисосудистое свертывание \\
• Наружные паразиты (блохи и клещи и т.д.) & \\
\midrule
\multicolumn{2}{c}{Эволюция регенерации при геморрагических клинических картинах} \\
Острое & Нет признаков регенерации в первые 2--3 дня. В первые часы HCT может не падать (из-за сокращения селезенки), поэтому диагноз ставится на основании анамнеза и снижения общего белка. \\
Подострое & Признаки регенерации (пропорционально текущему кровотечению). Типичный реактивный тромбоцитоз. \\
Хроническое & Нет признаков регенерации. Начинается микроцитарная и гипохромная железодефицитная анемия. Наступление этой последней фазы зависит от запасов железа у животного. \\
\bottomrule
\end{tabular}
\end{table}

\section{Случай 3. Кот с общей слабостью}
\subsection*{Клиническая история и физикальное обследование}
Некастрированный 8-летний кот, владельцы которого описали недавний эпизод множественной травмы (был сбит машиной 20 дней назад; рис. 1). С тех пор животное проявляет апатию и выраженную потерю аппетита с периодическими вокализациями. Животное не было вакцинировано, но было дегельминтизировано. Физикальное обследование выявило легкую депрессию, заметную одышку и бледные слизистые оболочки, но адекватную гидратацию.

\begin{tabular}{@{}ll@{}}
\toprule
Ректальная температура & 37.4°C \\
ЧСС & 150 уд/мин \\
ЧД & Не оценивается \\
ВНК & Менее 2 сек \\
\bottomrule
\end{tabular}

\begin{figure}[htbp]
\centering
\includegraphics[width=0.5\textwidth]{placeholder.jpg}
\caption{Пациент при поступлении.}
\end{figure}

\subsection*{Гемограмма}
\begin{tabular}{lccccc}
\toprule
Тест & Результат & Реф. диапазон & LOW & NORMAL & HIGH \\
\midrule
HCT & 18\% & 24--45 & $\blacksquare$ \\
RBC & 3.79 × 10\textsuperscript{12}/L & 5.0--10.0 & $\blacksquare$ \\
Hb & 62 g/L & 80--150 & $\blacksquare$ \\
MCV & 52 fL & 41--58 & $\blacksquare$ \\
RDW & 24\% & 17.3--22.0 & $\blacksquare$ \\
MCHC & 340 g/L & 300--360 & $\blacksquare$ \\
MCH & 16 pg & 12.5--17.5 & $\blacksquare$ \\
WBC & 26.3 k/$\mu$L & 5.5--19.5 & $\blacksquare$ \\
NEU & 14.5 k/$\mu$L & 2.5--14.0 & $\blacksquare$ \\
LYM & 9.82 × 10\textsuperscript{9}/L & 1.5--7.0 & $\blacksquare$ \\
MONO & 1.02 × 10\textsuperscript{9}/L & 0--1.50 & $\blacksquare$ \\
EOS & 0.75 × 10\textsuperscript{9}/L & 0--1.00 & $\blacksquare$ \\
BASO & 0.1 × 10\textsuperscript{9}/L & 0--0.20 & $\blacksquare$ \\
PLT & 1000 × 10\textsuperscript{9}/L & 300--800 & $\blacksquare$ \\
MPV & 12 fL & 12--17 & $\blacksquare$ \\
\% RETIC* & 0.5\% & 0.2--1.0 & $\blacksquare$ \\
\% BAND* & 0.3\% & 0--3.5 & $\blacksquare$ \\
\bottomrule
\multicolumn{6}{l}{*Ручной подсчет}
\end{tabular}

\subsection*{Вопрос: Какие изменения можно наблюдать в гемограмме этого животного?}
\subsection*{Мазок крови и окраска ретикулоцитов}
\begin{figure}[htbp]
\centering
\includegraphics[width=0.6\textwidth]{placeholder.jpg}
\caption{Микроскопическое изображение мазка крови. Окраска Diff-Quik, 400x.}
\end{figure}
\begin{figure}[htbp]
\centering
\includegraphics[width=0.6\textwidth]{placeholder.jpg}
\caption{Мазок крови пациента. Окраска Diff-Quik, 400x.}
\end{figure}
\begin{figure}[htbp]
\centering
\includegraphics[width=0.6\textwidth]{placeholder.jpg}
\caption{Окраска ретикулоцитов (новый метиленовый синий). 400x. Процент ретикулоцитов (ручной): 0.5\% агрегатных; 20\% пунктатных.}
\end{figure}

\subsection*{Вопрос: Какие изменения наблюдались в мазке крови этого животного? Какова степень медуллярной регенерации у пациента?}

\subsection*{Рентгенограмма и дальнейшие тесты}
\begin{figure}[htbp]
\centering
\includegraphics[width=0.6\textwidth]{placeholder.jpg}
\caption{Латеро-латеральная рентгенограмма животного.}
\end{figure}

\subsection*{Вопрос: Учитывая рентгенологические и гематологические находки, каков диагноз? Какие еще тесты следует провести этому пациенту?}

\subsection*{Разбор случая}
\subsubsection*{Какие изменения можно наблюдать в гемограмме этого животного?}
Гемограмма у этого кота указывала на анемию (снижение HCT, RBC и Hb) с явно регенераторным характером (не было увеличения процента ретикулоцитов или MCV, только небольшое увеличение RDW). Однако следует учитывать, что у кошек есть два типа ретикулоцитов. Поэтому для оценки степени медуллярной регенерации у кошек всегда показано использование специальной окраски для подсчета процента агрегатных и пунктатных ретикулоцитов (табл. 1).

На уровне лейкограммы у пациента наблюдался лейкоцитоз со зрелой нейтрофилией (без левого сдвига) и лимфоцитоз. Отсутствие увеличения палочкоядерных нейтрофилов исключало активный миелоидный медуллярный ответ (воспалительное и/или инфекционное состояние маловероятно), хотя эта находка всегда должна быть проверена с помощью мазка крови. У кошек это изменение в лейкограмме часто связано со стрессом из-за эффекта катехоламинов (вследствие неправильного обращения при венепункции или более устойчиво у госпитализированных животных или животных, испытывающих боль).

Наконец, в гемограмме наблюдался выраженный тромбоцитоз. Основными причинами этой находки являются стресс (эффект катехоламинов и сокращение селезенки), ответ на потерю тромбоцитов или анемию (реактивный тромбоцитоз), эссенциальная тромбоцитемия (редкое лейкемическое состояние) и ошибки распознавания тромбоцитов анализатором (из-за присутствия телец Гейнца или выраженного гемолиза).

\paragraph{На заметку}
Стресс может вызывать лейкоцитоз ($< 30,000 / \mu\mathrm{L}$) у кошек за счет зрелой нейтрофилии в сочетании с лимфоцитозом (обычно $< 10,000 / \mu\mathrm{L}$).

\paragraph{Рекомендуемое действие}
Проверьте отсутствие палочкоядерных нейтрофилов и аномальных лимфоцитов. Повторите забор пробы через 30-60 минут (если причина стресса больше не присутствует).

\begin{table}[htbp]
\centering
\caption{Характеристики различных типов ретикулоцитов, которые могут появляться у кошек.}
\begin{tabular}{p{2.5cm}p{3cm}p{2.5cm}p{2.5cm}p{2cm}}
\toprule
Тип & Изображение & Окраска новым метиленовым синим & Окраска Diff-Quik & Интерпретация \& Обнаружение автоматическим оборудованием \\
\midrule
Агрегатные & 1--2 агрегата или голубоватые кластеры среднего--крупного размера. & Полихроматофилы и макроциты. & Медуллярная регенерация за последние 3--6 дней. & ДА \\
Пунктатные & Множественные мелкие голубоватые пятна. & Нет полихроматофилов; некоторые макроциты. & Медуллярная регенерация за последние 7--21 день. & НЕТ \\
\bottomrule
\end{tabular}
\end{table}

\subsubsection*{Какие изменения наблюдались в мазке крови? Степень регенерации?}
В отношении лейкоцитов значительных изменений не было (рис. 2 и 3), что подтвердило, что лимфоциты имели типичную морфологию и палочкоядерных нейтрофилов не было. Эти находки подтвердили, что изменения в лейкограмме были вторичны по отношению к стрессу. Однако наблюдалось образование «монетных столбиков» (руло) эритроцитов (рис. 3). Этот паттерн часто наблюдается у кошек и еще более част при гиперглобулинемии (и, следовательно, при какой-то основной патологии). Хотя наличие руло неспецифично, важно отличать это изменение от аутоагглютинации (рис. 6). Лучший способ различить два процесса — выполнить разведение образца физиологическим раствором (1:4, хотя иногда необходимо до 1:10), чтобы наблюдать неокрашенную аликвоту под микроскопом; руло исчезают в разведенных образцах, в то время как аутоагглютинация обычно сохраняется.

Наконец, при окраске новым метиленовым синим (рис. 4) в основном наблюдались пунктатные и очень мало агрегатных ретикулоцитов (что подтверждало данные анализатора). При оценке степени медуллярной регенерации использование абсолютных значений или ретикулоцитарных индексов всегда гораздо информативнее, чем использование процента ретикулоцитов (табл. 2).

У этого пациента оценка обеих популяций ретикулоцитов показала, что регенераторный ответ изначально был умеренным или слабым (7--21 день назад, о чем свидетельствует количество пунктатных эритроцитов), но недавно изменился со слабого на очень слабый (3--6 дней назад, согласно количеству агрегатных эритроцитов). Этот паттерн обычно наблюдается, когда причина анемии разрешилась (и, следовательно, регенерация становится все менее необходимой) или у животных, у которых медуллярный ответ ослабевает (особенно если причина анемии все еще присутствует).

\begin{figure}[htbp]
\centering
\includegraphics[width=0.8\textwidth]{placeholder.jpg}
\caption{Сравнение образцов с руло (слева), где эритроциты обычно располагаются рядами или стопками, и аутоагглютинацией (справа), где эритроциты образуют большие трехмерные кластеры. Наиболее точная дифференцировка требует определения сохранения этих явлений после разведения физиологическим раствором.}
\end{figure}

\begin{table}[htbp]
\centering
\caption{Оценка степени медуллярной регенерации у мелких животных с использованием различных ретикулоцитарных индексов.}
\begin{tabular}{p{3cm}p{2cm}p{2cm}p{2cm}p{2cm}p{2cm}p{2cm}}
\toprule
\multirow{2}{*}{Степень регенерации} & \multicolumn{3}{c}{Собака} & \multicolumn{3}{c}{Кошка} \\
\cmidrule(lr){2-4} \cmidrule(lr){5-7}
& \% RETIC & ARC & \% Агрегатные RETIC & \% Пунктатные RETIC & Агрегатные RETIC ARC & Пунктатные RETIC ARC \\
\midrule
Отсутствует & 175,000 & 0--0.4 & 1--10 & <15,000 & <200,000 \\
Слабая & 1--4 & 150,000 & 0.5--2 & 10--20 & 50,000 & 200,000--500,000 \\
Умеренная & 5--20 & 300,000 & 3--5 & 20--50 & 100,000 & 500,000--1,000,000 \\
Выраженная & 21--50 & >500,000 & >5 & >50 & >200,000 & >1,500,000 \\
\bottomrule
\multicolumn{7}{l}{Абсолютное количество ретикулоцитов (ARC) = \% RETIC × количество эритроцитов} \\
\multicolumn{7}{l}{Другие индексы (наиболее используемые у собак):} \\
\multicolumn{7}{l}{Скорректированный процент ретикулоцитов (CRP) = \% RETIC × (HCT пациента / Нормальный HCT)} \\
\multicolumn{7}{l}{Принимая HCT 45\% у собак и 35\% у кошек за норму.} \\
\multicolumn{7}{l}{Регенерация, если CRP > 1\% у собак или > 0.4\% у кошек.} \\
\multicolumn{7}{l}{Индекс продукции ретикулоцитов (RPI) = (HCT пациента / Нормальный HCT) × \% RETIC} \\
\multicolumn{7}{l}{Те же нормальные HCT, что и для CRP. Время созревания в зависимости от HCT: 45 = 1; 35 = 1.5; 25 = 2; 15 = 2.5.} \\
\multicolumn{7}{l}{RPI < 1 — нет регенерации. RPI = 2 — регенерация со скоростью, вдвое превышающей норму.} \\
\end{tabular}
\end{table}

\subsubsection*{Учитывая рентгенологические и гематологические находки, каков диагноз? Какие еще тесты?}
Рентгенологически у пациента выявлен обильный плевральный выпот (обведен желтым на рис. 7); его содержимое было совместимо с активным кровотечением (вставка на рис. 7).

Вместе с гематологическими изменениями эта находка указывала на клиническую картину подострой геморрагической анемии вследствие травмы (совпадающей с анамнезом). В ходе операции были обнаружены сломанное ребро и разрыв легкого, которые также вызывали анемию. Геморрагическая анемия с такой эволюцией во времени обычно регенераторная (случай 2). Однако, учитывая слабый медуллярный регенераторный ответ, наблюдаемый у этого кота (который ухудшился, когда кровотечение все еще было активным), необходимо провести более детальное исследование для поиска причин недостаточности медуллярного эритропоэза (цитология или биопсия костного мозга; диагностика вирусных клинических картин, которые могут вызывать медуллярную недостаточность; исключить контакт с лекарствами или миелосупрессивными токсинами; изучить функцию почек пациента и т.д.). Наличие руло может указывать на существование какого-либо хронического основного состояния, которое может быть связано с медуллярной недостаточностью у этого пациента.

Наконец, остальные гематологические изменения (в лейкограмме и тромбоцитоз) легко объясняются уровнем стресса у пациента, а также его ответом на анемию.

\begin{figure}[htbp]
\centering
\includegraphics[width=0.5\textwidth]{placeholder.jpg}
\caption{Травматический гемоторакс.}
\end{figure}

\section{Случай 4. Плановый осмотр кота}
\subsection*{Клиническая история и физикальное обследование}
10-летний некастрированный кот поступил на плановый осмотр. Пациент был правильно вакцинирован и дегельминтизирован (рис. 1). Владельцы сообщили о прогрессирующей потере аппетита, небольшой потере веса и мышечной массы, а также о некоторой вялости. При физикальном обследовании выявлены слегка бледные и желтушные слизистые оболочки, но адекватное состояние гидратации.

\begin{tabular}{@{}ll@{}}
\toprule
Ректальная температура & 38.4°C \\
ЧСС & 140 уд/мин \\
ЧД & 20 дых/мин \\
ВНК & Менее 2 сек \\
\bottomrule
\end{tabular}

\begin{figure}[htbp]
\centering
\includegraphics[width=0.5\textwidth]{placeholder.jpg}
\caption{Пациент при поступлении.}
\end{figure}

\subsection*{Гемограмма}
\begin{tabular}{lccccc}
\toprule
Тест & Результат & Реф. диапазон & LOW & NORMAL & HIGH \\
\midrule
HCT & 19\% & 24--45 & $\blacksquare$ \\
RBC & 4.19 × 10\textsuperscript{12}/L & 5.0--10.0 & $\blacksquare$ \\
Hb & 67 g/L & 80--150 & $\blacksquare$ \\
MCV & 37 fL & 39--55 & $\blacksquare$ \\
RDW & 26\% & 17.3--22.0 & $\blacksquare$ \\
MCHC & 350 g/L & 300--360 & $\blacksquare$ \\
MCH & 16 pg & 12.5--17.5 & $\blacksquare$ \\
WBC & 16.3 × 10\textsuperscript{9}/L & 5.5--19.5 & $\blacksquare$ \\
NEU & 10.5 × 10\textsuperscript{9}/L & 2.5--14.0 & $\blacksquare$ \\
LYM & 3.86 × 10\textsuperscript{9}/L & 1.5--7.0 & $\blacksquare$ \\
MONO & 1.55 × 10\textsuperscript{9}/L & 0--1.5 & $\blacksquare$ \\
EOS & 0.24 × 10\textsuperscript{9}/L & 0--1.0 & $\blacksquare$ \\
BASO & 0.05 × 10\textsuperscript{9}/L & 0--0.20 & $\blacksquare$ \\
PLT & 550 × 10\textsuperscript{9}/L & 300--800 & $\blacksquare$ \\
MPV & 14 fL & 12--17 & $\blacksquare$ \\
\% RETIC* & 0.3\% & 0.2--1.0 & $\blacksquare$ \\
\% BAND* & 0.6\% & 0--3.5 & $\blacksquare$ \\
\bottomrule
\multicolumn{6}{l}{*Ручной подсчет}
\end{tabular}

\subsection*{Вопрос: Какие изменения можно наблюдать в гемограмме этого животного?}

\subsection*{Мазки крови}
\begin{figure}[htbp]
\centering
\includegraphics[width=0.6\textwidth]{placeholder.jpg}
\caption{Микроскопическое изображение мазка крови. Окраска Diff-Quik, 400x.}
\end{figure}
\begin{figure}[htbp]
\centering
\includegraphics[width=0.6\textwidth]{placeholder.jpg}
\caption{Мазок крови пациента. Окраска Diff-Quik, 1000x.}
\end{figure}

\subsection*{Вопрос: Какие изменения наблюдались в мазке крови?}

\subsection*{УЗИ и цитологическое исследование}
\begin{figure}[htbp]
\centering
\includegraphics[width=0.6\textwidth]{placeholder.jpg}
\caption{Ультразвуковой вид левой почки (изображение правой почки было аналогичным). Обе почки были слегка уменьшены в размере.}
\end{figure}
\begin{figure}[htbp]
\centering
\includegraphics[width=0.6\textwidth]{placeholder.jpg}
\caption{Цитология, полученная тонкоигольной пункцией печени без аспирации. Окраска Diff-Quik, 1000x.}
\end{figure}

\subsection*{Вопрос: Учитывая ультразвуковые, цитологические и гематологические данные, каков диагноз?}

\subsection*{Разбор случая}
\subsubsection*{Нерегенераторная анемия вследствие хронического заболевания}

\subsubsection*{Какие изменения можно наблюдать в гемограмме этого животного?}
Была обнаружена умеренная анемия (снижение HCT, RBC и Hb), которая казалась регенераторной (нет повышения процента ретикулоцитов или MCV, только небольшое повышение RDW). Аналогично, у этого пациента не было увеличения пунктатных ретикулоцитов, поэтому более хронический регенераторный ответ можно было исключить. Основные причины нерегенераторной анемии перечислены в таблице 1. Лейкограмма анализатора и анализ тромбоцитов не выявили отклонений.

\begin{table}[htbp]
\centering
\caption{Наиболее частые причины регенераторной анемии у мелких животных.}
\begin{tabular}{p{14cm}}
\toprule
Наиболее частые причины НЕрегенераторной анемии у мелких животных \\
\midrule
Острое (2--3 дня) или хроническое кровотечение (дефицит железа). \\
Подавление эритропоэза (анемия, вызванная хроническим заболеванием или воспалением, нефропатией, токсинами или лекарствами и т.д.). \\
Инфекции (парвовирус и др.). \\
Миелопатии (миелофиброз, изолированная аплазия красного ростка или неоплазия и т.д.). \\
Лучевая терапия или химиотерапия. \\
\bottomrule
\end{tabular}
\end{table}

\subsubsection*{Какие изменения наблюдались в мазке крови?}
Наиболее поразительным изменением была выраженная пойкилоцитоз (вариабельность формы эритроцитов), как показано на рис. 2 и 3. Пойкилоцитоз вызван наличием многочисленных акантоцитов и эллиптоцитов. Акантоциты необходимо дифференцировать от эхиноцитов (случай 1), которые характеризуются появлением нерегулярных тупых выступов, асимметрично распределенных по их поверхности (рис. 6).

\begin{figure}[htbp]
\centering
\includegraphics[width=0.8\textwidth]{placeholder.jpg}
\caption{Сравнение эхиноцитов (a) и акантоцитов (b). Обратите внимание на отсутствие однородности выступов у акантоцита по сравнению с эхиноцитом. Кроме того, эхиноцитоз обычно затрагивает большинство эритроцитов в мазке.}
\end{figure}

\paragraph{На заметку}
Акантоциты могут появляться при различных патологиях у мелких животных, хотя они обычно связаны с фрагментационным повреждением (вместе с наличием шистоцитов) или наличием изменений в липидном составе плазмы (главным образом холестерина и липопротеинов).

\paragraph{Патологии, при которых появляются акантоциты}
ДВС-синдром, заболевания печени, дефицит железа, опухоли и др.

Что касается эллиптоцитов, это морфологическое изменение вызвано подтипом акантоцитов, который может появляться при различных патологиях, хотя у кошек он довольно специфично связан с клинической картиной печеночного липидоза и другими заболеваниями печени (рис. 7). Когда эти эллиптоциты имеют зазубренную поверхность, их также называют «клетками кленового листа» или «листьями остролиста».

\begin{figure}[htbp]
\centering
\includegraphics[width=0.4\textwidth]{placeholder.jpg}
\caption{Эллиптоциты (клетки кленового листа).}
\end{figure}

\subsubsection*{Учитывая данные УЗИ, цитологии и гематологии, каков диагноз?}
На УЗИ почек (рис. 4) обнаружена измененная эхогенность как коры почек (гиперэхогенная по сравнению с окружающими тканями), так и мозгового вещества; обе части также имели плохую дифференцировку и были неоднородными. Эти находки заставили заподозрить нефропатию (подтвержденную у этого пациента лабораторными маркерами повреждения почек). Что касается печени (рис. 5), цитологический образец показал данные, типичные для печеночного липидоза (тяжелая вакуолярная дегенерация гепатоцитов с наличием вакуолей с четкими границами и прозрачным содержимым, которые часто оттесняли ядро к периферии, и которые также были видны в фоне препарата).

Эти находки, вместе с гематологическими изменениями, указывают на клиническую картину нерегенераторной анемии, вызванной хроническим заболеванием печени и почек. Повышение RDW и снижение MCV у этого пациента было результатом выраженного пойкилоцитоза. Маркеров, указывающих на регенерацию (ни в анализаторе, ни в мазке крови), не наблюдалось. Учитывая, что это был пациент кошачьего вида, отсутствие явных ореолов исключало первичную клиническую картину дефицита железа у этого животного. Хотя для исключения миелосупрессивных инфекций или контакта с лекарствами или миелосупрессивными токсинами необходимо было бы провести детальное исследование костного мозга, совокупность проведенных тестов указывала на то, что наиболее вероятными причинами анемии были заболевание печени и нефропатия. В таблице 2 показано патофизиогенетическое объяснение этих типов нерегенераторных анемий.

\begin{table}[htbp]
\centering
\caption{Патофизиогенетические механизмы анемии при хронических заболеваниях и при заболеваниях почек (обе НЕрегенераторные).}
\begin{tabular}{p{7cm}p{7cm}}
\toprule
Патофизиогенетические механизмы анемии при хроническом заболевании & Патофизиогенетические механизмы анемии, вызванной заболеванием почек \\
\midrule
• Делеционные эффекты воспалительных цитокинов на эритропоэз. & • Снижение продукции эритропоэтина. \\
• Снижение продукции или эффективности эритропоэтина. & • Уремия (подавляет эритропоэз и может вызывать гемолиз). \\
• Снижение использования запасов железа. & • Развитие гастроэнтерических язв (хроническое кровотечение). \\
• Заболевания, наиболее часто связанные с этим процессом: эндокринопатии (гипотиреоз, гипоадренокортицизм или гиперэстрогенизм), инфекционные клинические картины, воспаление или новообразования. & • Тромбоцитопатический эффект уремических токсинов (хроническое кровотечение). \\
\bottomrule
\end{tabular}
\end{table}

\section{Случай 5. Собака с общей слабостью и кашлем}
\subsection*{Клиническая история и физикальное обследование}
15-летняя нестерилизованная сука смешанной породы поступила в клинику с постоянным кашлем в течение двух месяцев и клинической картиной общей слабости и апатии. При аускультации патологических шумов не выявлено. Слизистые оболочки были слегка бледными.

\begin{tabular}{@{}ll@{}}
\toprule
Ректальная температура & 38.3°C \\
ЧСС & 160 уд/мин \\
ЧД & 50 дых/мин \\
ВНК & Менее 2 сек \\
\bottomrule
\end{tabular}

\subsection*{Гемограмма}
\begin{tabular}{lccccc}
\toprule
Тест & Результат & Реф. диапазон & LOW & NORMAL & HIGH \\
\midrule
HCT & 26\% & 37--55 & $\blacksquare$ \\
RBC & 4.06 × 10\textsuperscript{12}/L & 5.50--8.50 & $\blacksquare$ \\
Hb & 85 g/L & 120--180 & $\blacksquare$ \\
MCV & 64 fL & 60--77 & $\blacksquare$ \\
RDW & 13\% & 12.0--15.0 & $\blacksquare$ \\
MCHC & 330 g/L & 310--340 & $\blacksquare$ \\
MCH & 26 pg & 19.5--24.5 & $\blacksquare$ \\
WBC & 12.5 × 10\textsuperscript{9}/L & 6.0--17.0 & $\blacksquare$ \\
NEU & 10 × 10\textsuperscript{9}/L & 3.0--12.0 & $\blacksquare$ \\
LYM & 1.41 × 10\textsuperscript{9}/L & 1--4.80 & $\blacksquare$ \\
MONO & 0.77 × 10\textsuperscript{9}/L & 0.20--1.50 & $\blacksquare$ \\
EOS & 0.38 × 10\textsuperscript{9}/L & 0--0.80 & $\blacksquare$ \\
BASO & 0 × 10\textsuperscript{9}/L & 0--0.40 & $\blacksquare$ \\
PLT & 54 × 10\textsuperscript{9}/L & 200--500 & $\blacksquare$ \\
MPV & 7.46 fL & 3.9--11.1 & $\blacksquare$ \\
\% RETIC* & 2\% & 0.5--1.5 & $\blacksquare$ \\
\% BAND* & 3\% & 0--3.5 & $\blacksquare$ \\
\bottomrule
\multicolumn{6}{l}{*Ручной подсчет}
\end{tabular}

\subsection*{Вопрос: Какие изменения можно наблюдать в гемограмме этого животного?}
\begin{figure}[htbp]
\centering
\includegraphics[width=\textwidth]{placeholder.jpg}
\caption{Гистограммы крови пациента, классифицирующие элементы по размеру (горизонтальная ось) и количеству (вертикальная ось); (a) графическое представление лейкоцитов; (b) эритроцитов; и (c) тромбоцитов.}
\end{figure}

\subsection*{Мазок крови}
\begin{figure}[htbp]
\centering
\includegraphics[width=0.6\textwidth]{placeholder.jpg}
\caption{Микроскопическое изображение мазка крови. Окраска Diff-Quick. 400x.}
\end{figure}

\subsection*{Вопрос: Какие изменения можно наблюдать в мазке крови?}
\begin{figure}[htbp]
\centering
\includegraphics[width=0.6\textwidth]{placeholder.jpg}
\caption{Латеро-латеральная рентгенограмма грудной клетки пациента.}
\end{figure}

\subsection*{Вопрос: Учитывая эту рентгенограмму, каков диагноз?}

\subsection*{Разбор случая}
\textbf{Анемия с микроангиопатической фрагментацией, вызванная метастазами в легкие}

\subsubsection*{Какие изменения можно наблюдать в гемограмме?}
У животного наблюдается регенераторная нормохромная нормоцитарная анемия с тромбоцитопенией. Гистограммы животного (рис. 1) показывают регенераторную анемию (асимметрия вправо эритроидной кривой из-за наличия многочисленных ретикулоцитов; рис. 1b) и истинную тромбоцитопению (типичная форма кривой, отсутствие тромбоцитарных агрегатов, образующих «лестничную форму»; рис. 1c). Лейкоцитарная гистограмма (рис. 1a) нормальна для собак.

\subsubsection*{Какие изменения наблюдались в мазке крови?}
В мазке крови (рис. 2) видны анизоцитоз, пойкилоцитоз и шистоциты.

\paragraph{На заметку}
Наличие шистоцитов в мазке крови обычно указывает на фрагментацию эритроцитов, которая может быть микро- или макроангиопатической. Шистоциты (с аномальной удлиненной формой) необходимо дифференцировать от акантоцитов (с асимметричными и неравномерно распределенными выступами на поверхности).

\paragraph{Основные причины фрагментации эритроцитов:} диссеминированное внутрисосудистое свертывание (ДВС), пороки клапанов сердца, гломерулярные заболевания, высоковаскуляризированные опухоли (гемангиосаркомы), \textit{Dirofilaria immitis} и др.

\begin{figure}[htbp]
\centering
\includegraphics[width=0.8\textwidth]{placeholder.jpg}
\caption{Разница между акантоцитами (a, b) и шистоцитами (c).}
\end{figure}

\subsubsection*{Учитывая этот мазок, каков диагноз?}
Учитывая отличительные признаки эритроидной фрагментации, наблюдаемые в гемограмме (рис. 4), вместе с наличием тромбоцитопении и изображением, совместимым с метастазами в легкие типа «пушечного ядра», наиболее вероятным диагнозом будет высоковаскуляризированная злокачественная опухоль с метастазами в легкие (рис. 5), возможно, гемангиосаркома.

\begin{figure}[htbp]
\centering
\includegraphics[width=0.6\textwidth]{placeholder.jpg}
\caption{Латеро-латеральная рентгенограмма грудной клетки пациента, показывающая многочисленные округло-овальные участки с интерстициальным рисунком в легочном поле, совместимые с наличием метастазов в легочной паренхиме.}
\end{figure}

\section{Случай 6. Атаксичная и спутанная сука}
\subsection*{Клиническая история и физикальное обследование}
Нестерилизованная 4-летняя сука смешанной породы поступила в клинику с внезапным неврологическим состоянием после пребывания в автомобиле в течение 30 минут. Пациент был правильно вакцинирован и дегельминтизирован, без предшествующей истории других процессов. Физикальное обследование выявило выраженную спутанность сознания, атаксию, гиперсаливацию, слегка бледные слизистые оболочки и наличие петехий и экхимозов (рис. 1), а также снижение тургора кожи с клинической оценкой обезвоживания в $7\%$. У животного также была рвота и обильная гематохезия.

\begin{tabular}{@{}ll@{}}
\toprule
Ректальная температура & 42.0°C \\
ЧСС & 180 уд/мин \\
ЧД & 40 дых/мин \\
ВНК & Менее 2 сек \\
\bottomrule
\end{tabular}

\begin{figure}[htbp]
\centering
\includegraphics[width=0.5\textwidth]{placeholder.jpg}
\caption{При поступлении у пациента наблюдалось конъюнктивальное кровоизлияние.}
\end{figure}

\subsection*{Гемограмма и биохимия}
\begin{tabular}{lccccc}
\toprule
Тест & Результат & Реф. диапазон & LOW & NORMAL & HIGH \\
\midrule
HCT & 39\% & 37--55 & $\blacksquare$ \\
RBC & 5.91 × 10\textsuperscript{12}/L & 5.50--8.50 & $\blacksquare$ \\
Hb & 131 g/L & 120--180 & $\blacksquare$ \\
MCV & 68 fL & 60--77 & $\blacksquare$ \\
RDW & 14\% & 12.0--15.0 & $\blacksquare$ \\
MCHC & 330 g/L & 310--340 & $\blacksquare$ \\
MCH & 22 pg & 19.5--24.5 & $\blacksquare$ \\
WBC & 19.1 × 10\textsuperscript{9}/L & 6.0--17.0 & $\blacksquare$ \\
NEU & 10.9 × 10\textsuperscript{9}/L & 3.0--12.0 & $\blacksquare$ \\
LYM & 6.39 × 10\textsuperscript{9}/L & 1--4.80 & $\blacksquare$ \\
MONO & 1.01 × 10\textsuperscript{9}/L & 0.20--1.50 & $\blacksquare$ \\
EOS & 0.75 × 10\textsuperscript{9}/L & 0--0.80 & $\blacksquare$ \\
BASO & 0 × 10\textsuperscript{9}/L & 0--0.40 & $\blacksquare$ \\
PLT & 100 × 10\textsuperscript{9}/L & 200--500 & $\blacksquare$ \\
MPV & 9.02 fL & 3.9--11.1 & $\blacksquare$ \\
Общий белок & 4.80 g/dL & 6--7.5 & $\blacksquare$ \\
Альбумин-глобулиновое соотношение & 0.9 & 0.52--1.15 & $\blacksquare$ \\
\% RETIC* & 0.5\% & 0.5--1.5 & $\blacksquare$ \\
\% BAND* & 1.8\% & 0--3.5 & $\blacksquare$ \\
\bottomrule
\multicolumn{6}{l}{*Ручной подсчет}
\end{tabular}

\subsection*{Вопрос: Какие изменения можно наблюдать в гемограмме этого животного?}

\subsection*{Мазки крови}
\begin{figure}[htbp]
\centering
\includegraphics[width=0.6\textwidth]{placeholder.jpg}
\caption{Мазок крови пациента. Окраска Diff-Quik, 1000x.}
\end{figure}
\begin{figure}[htbp]
\centering
\includegraphics[width=0.6\textwidth]{placeholder.jpg}
\caption{Различные детали мазка крови. Окраска Diff-Quik, 1000x.}
\end{figure}

\subsection*{Вопрос: Какие изменения наблюдались в мазке крови?}

\subsection*{Гемостатическое исследование}
\begin{tabular}{lcc}
\toprule
Тест & Результат & Референсный диапазон \\
\midrule
Прогромбиновое время & 29 секунд & 6--12 \\
Активированное частичное тромбопластиновое время & 42 секунды & 12--22 \\
Время кровотечения слизистой оболочки щеки & 6.5 минут & < 4 \\
D-димер & 1,500 ng/mL & < 250 \\
\bottomrule
\end{tabular}

\subsection*{Вопрос: Учитывая гемостатические и гематологические находки, каков диагноз?}

\subsection*{Разбор случая}
\textbf{Диссеминированное внутрисосудистое свертывание (ДВС) вследствие теплового удара}

\subsubsection*{Какие изменения можно наблюдать в гемограмме?}
У животного не было заметных изменений в эритрограмме, хотя мы должны учитывать, что обезвоживание ложно повышает эти показатели. Клинически наличие кровоизлияний, гипопротеинемия (даже при обезвоживании) и сохранение альбумин-глобулинового соотношения делали наличие подострой геморрагической анемии очень вероятным. Эта находка была подтверждена после регидратации пациента (HCT: 34\%; общий белок: 4.4 g/dL).

В лейкограмме обнаружен легкий лейкоцитоз с лимфоцитозом.

Наконец, анализ тромбоцитов показал умеренную или выраженную тромбоцитопению. Эта находка была верифицирована ручным подсчетом в цитратных пробах, что привело к классификации истинной тромбоцитопении (табл. 1).

\paragraph{На заметку}
Основными причинами лимфоцитоза являются инфекции с высокой антигенной нагрузкой, вакцинации, выброс адреналина вследствие обращения (у кошек), гипоадренокортицизм, лимфоидные лейкемии или ошибки классификации лейкоцитов автоматическими анализаторами (например, из-за присутствия ядерных эритроцитов).

\begin{table}[htbp]
\centering
\caption{Основные причины тромбоцитопении у мелких животных и приблизительная формула для коррекции ошибочных подсчетов тромбоцитов.}
\begin{tabular}{p{14cm}}
\toprule
Наиболее частые причины тромбоцитопении у мелких животных \\
\midrule
• Ложная тромбоцитопения (псевдотромбоцитопения): наличие тромбоцитарных агрегатов или сбои автоматического анализатора и т.д. \\
• Истинная тромбоцитопения: снижение тромбопоэза или медуллярная недостаточность (идиопатическая, ретровирусная инфекция или \textit{Anaplasma platys} и т.д.). Повышенное использование или разрушение: ДВС, эрлихиоз, бабезиоз, спленомегалия или иммуноопосредованная тромбоцитопения и т.д. \\
\midrule
Формула для расчета (приблизительного) количества тромбоцитов в мазках крови. \\
N° тромбоцитов/мкл = (наблюдаемые WBC × WBC/мкл) \\
Подсчитайте эритроциты и тромбоциты как минимум в 30 полях (40×) мазка. \\
\bottomrule
\end{tabular}
\end{table}

\subsubsection*{Какие изменения наблюдались в мазке крови?}
В мазке крови обнаружены многочисленные ядерные эритроциты (нормобластоз) (рис. 2), но без других признаков заметной регенерации (и этот пациент не был новорожденным), следовательно, имела место неадекватная метарубрицитемия, указывающая на выраженное повреждение костного мозга. Как упоминалось в случае 1, эти клетки могут мешать подсчету лейкоцитов (лимфоцитов), поэтому мы должны скорректировать полученные данные (рис. 9 из случая 1). После этой коррекции показатели были в пределах нормы.

Также наблюдались многочисленные признаки фрагментации эритроцитов, такие как наличие шистоцитов, акантоцитов и кератоцитов (рис. 3). Кератоциты (серповидной формы) могут появляться как вследствие окислительного повреждения, так и микроангиопатической фрагментации.

\subsubsection*{Учитывая гемостатические и гематологические находки, каков диагноз?}
На гемостатическом уровне наблюдалась коагулопатия потребления с первичной (удлиненное время кровотечения из слизистой оболочки щеки) и вторичной (активированное частичное тромбопластиновое и протромбиновое время) недостаточностью гемостаза и гиперфибринолизом (повышенные D-димеры). Эти находки вместе с тромбоцитопенией подтвердили клиническую картину ДВС (табл. 2). Появление признаков фрагментации эритроцитов в мазках крови типично для этого синдрома. Учитывая анамнез и клинические данные у этого пациента, наиболее вероятной причиной ДВС был тепловой удар, что также объясняло бы наличие ядерных эритроцитов (вследствие повреждения костного мозга).

\paragraph{На заметку}
Диссеминированное внутрисосудистое свертывание (ДВС) всегда является вторичным состоянием и характеризуется коагулопатией потребления с гиперкоагуляцией и гиперфибринолизом.

\paragraph{Основные причины ДВС}
Травма, аутоиммунная гемолитическая анемия, панкреатит, тепловой удар, дилатация или заворот желудка, метастатические новообразования, сепсис и т.д.

\begin{table}[htbp]
\centering
\caption{Результаты различных гемостатических тестов при основных коагуляционных патологиях у мелких животных.}
\begin{tabular}{lccccc}
\toprule
Патология & Количество тромбоцитов & BMBT & PT & aPTT & D-димеры \\
\midrule
Тромбоцитопения & ↓ & ↑ & N & N & N \\
Тромбоцитопатии & N & ↑ & N & N & N \\
Болезнь фон Виллебранда & N & ↑ & N & N & N \\
Отравление варфарином & N & N & ↑ & ↑ & N \\
Печеночная недостаточность & N & N & ↑ & ↑ & N \\
ДВС & ↓ & ↑ & ↑ & ↑ & ↑ \\
\bottomrule
\end{tabular}
\end{table}

\section{Случай 7. Собака с желтушными слизистыми оболочками}
\subsection*{Клиническая история и физикальное обследование}
5-летний далматинец поступил на прием с выраженной желтухой всех слизистых оболочек (в течение как минимум недели), а также с картиной апатии и потери аппетита (рис. 1). Животное было правильно вакцинировано и дегельминтизировано, без предшествующей истории других процессов. Физикальное обследование выявило депрессию, легкую спутанность сознания, слегка бледные и желтушные слизистые оболочки, но адекватную гидратацию.

\begin{tabular}{@{}ll@{}}
\toprule
Ректальная температура & 40.2°C \\
ЧСС & 120 уд/мин \\
ЧД & 36 дых/мин \\
ВНК & Менее 2 сек \\
\bottomrule
\end{tabular}

\begin{figure}[htbp]
\centering
\includegraphics[width=0.5\textwidth]{placeholder.jpg}
\caption{Слизистая оболочка рта пациента при поступлении.}
\end{figure}

\subsection*{Гемограмма}
\begin{tabular}{lccccc}
\toprule
Тест & Результат & Реф. диапазон & LOW & NORMAL & HIGH \\
\midrule
HCT & 24\% & 37--55 & $\blacksquare$ \\
RBC & 3.13 × 10\textsuperscript{12}/L & 5.50--8.50 & $\blacksquare$ \\
Hb & 140 g/L & 120--180 & $\blacksquare$ \\
MCV & 80 fL & 60--77 & $\blacksquare$ \\
RDW & 28\% & 12.0--15.0 & $\blacksquare$ \\
MCHC & 580 g/L & 310--340 & $\blacksquare$ \\
MCH & 44 pg & 19.5--24.5 & $\blacksquare$ \\
WBC & 28.1 × 10\textsuperscript{9}/L & 6.0--16.0 & $\blacksquare$ \\
NEU & 20.1 × 10\textsuperscript{9}/L & 3.0--12.0 & $\blacksquare$ \\
LYM & 4.39 × 10\textsuperscript{9}/L & 1--4.80 & $\blacksquare$ \\
MONO & 3.42 × 10\textsuperscript{9}/L & 0.20--1.50 & $\blacksquare$ \\
EOS & 0.15 × 10\textsuperscript{9}/L & 0--0.80 & $\blacksquare$ \\
BASO & 0 × 10\textsuperscript{9}/L & 0--0.40 & $\blacksquare$ \\
PLT & 450× 10\textsuperscript{9}/L & 200--500 & $\blacksquare$ \\
MPV & 9.91 fL & 3.9--11.1 & $\blacksquare$ \\
\% RETIC* & 5.5\% & 0.5--1.5 & $\blacksquare$ \\
\% BAND* & 7.9\% & 0--3.5 & $\blacksquare$ \\
\bottomrule
\multicolumn{6}{l}{*Ручной подсчет}
\end{tabular}

\subsection*{Вопрос: Какие изменения можно наблюдать в гемограмме?}

\subsection*{Мазки крови}
\begin{figure}[htbp]
\centering
\includegraphics[width=0.6\textwidth]{placeholder.jpg}
\caption{Мазок крови пациента. Окраска Diff-Quik, 400x.}
\end{figure}
\begin{figure}[htbp]
\centering
\includegraphics[width=0.6\textwidth]{placeholder.jpg}
\caption{Деталь мазка крови. Окраска Diff-Quik, 1000x.}
\end{figure}
\begin{figure}[htbp]
\centering
\includegraphics[width=0.6\textwidth]{placeholder.jpg}
\caption{Мазок крови пациента. Окраска Diff-Quik, 400x.}
\end{figure}

\subsection*{Тест на аутоагглютинацию на предметном стекле}
\begin{figure}[htbp]
\centering
\includegraphics[width=0.6\textwidth]{placeholder.jpg}
\caption{Тест на аутоагглютинацию крови пациента (сверху) по сравнению со здоровым контролем (снизу).}
\end{figure}

\subsection*{Вопрос: Какие изменения наблюдались в мазке крови?}
\subsection*{Вопрос: Учитывая гематологические находки, каков диагноз?}

\subsection*{Разбор случая}
\subsubsection*{Какие изменения в гемограмме?}
У пациента наблюдалась заметная регенераторная анемия (низкие HCT и RBC, высокий процент ретикулоцитов, MCV и RDW) и гемолитическая (высокие Hb, MCHC и MCH). Наличие желтухи клинически исключало возможность артефактного гемолиза.

Также наблюдалась типичная воспалительная лейкограмма с лейкоцитозом, нейтрофилией, левым сдвигом (увеличение палочкоядерных нейтрофилов) и моноцитозом (табл. 1). В данном случае сдвиг влево был классифицирован как регенераторный, поскольку также присутствовала зрелая нейтрофилия.

\paragraph{На заметку}
Основными причинами нейтрофилии являются инфекционные или воспалительные состояния, избыток глюкокортикоидов (эндогенных или экзогенных), выброс катехоламинов (обращение и страх, особенно у кошек) и миелоидные лейкемии.

\begin{table}[htbp]
\centering
\caption{Паттерны лейкограммы с лейкоцитозом и нейтрофилией у мелких животных.}
\begin{tabular}{lcccccc}
\toprule
& Нейтрофилы & Палочкоядерные & Лимфоциты & Моноциты & Эозинофилы & Примеры \\
\midrule
Воспалительная лейкограмма & ↑ & ↑ & ≈ & ↑ & ≈ & Инфекции, панкреатит или аутоиммунная гемолитическая анемия и т.д. \\
Стрессовая лейкограмма & ↑ & - & ↓ & ↑ & ↓ & Гиперадренокортицизм, хронический стресс и т.д. \\
Лейкограмма обращения/страха & ≈ & - & ↑ & = & = & Реакция бегства у кошек. \\
\bottomrule
\multicolumn{7}{l}{=: вариабельно}
\end{tabular}
\end{table}

\subsubsection*{Какие изменения наблюдались в мазке крови?}
В мазке крови при малом увеличении видны образование «монетных столбиков» и наличие многочисленных полихроматофилов (рис. 2 и 3), связанных с выраженной регенерацией. При большем увеличении наблюдались многочисленные сфероциты (рис. 6), характерные для аутоиммунной гемолитической анемии.

\begin{figure}[htbp]
\centering
\includegraphics[width=0.6\textwidth]{placeholder.jpg}
\caption{Сравнение нормальных эритроцитов, сфероцитов (синие стрелки; круглой формы, немного меньше и без видимого ореола) и полихроматофилов (красные стрелки). Окраска Diff-Quik, 1000x.}
\end{figure}

\subsubsection*{Учитывая гематологические находки, каков диагноз?}
В тесте на предметном стекле макроскопически обнаружена аутоагглютинация (рис. 5). Эта находка также была подтверждена микроскопически, даже после разведения 1:20 физиологическим раствором (рис. 8), что с высокой вероятностью указывает на наличие аутоантител против эритроцитов (окончательное подтверждение потребовало бы проведения теста Кумбса). Учитывая эту находку и наличие сфероцитов в мазке, диагноз — аутоиммунная гемолитическая анемия.

На уровне лейкоцитов наблюдались многочисленные зрелые нейтрофилы, а также моноциты и обильные палочкоядерные нейтрофилы (рис. 7).

\begin{figure}[htbp]
\centering
\includegraphics[width=0.6\textwidth]{placeholder.jpg}
\caption{Дифференциация палочкоядерных нейтрофилов (синие стрелки; компактный и не «пережатый» хроматин) и моноцитов (красные стрелки, «переплетенный» хроматин). Окраска Diff-Quik, 400x.}
\end{figure}
\begin{figure}[htbp]
\centering
\includegraphics[width=0.6\textwidth]{placeholder.jpg}
\caption{Дифференциация аутоагглютинации (слева) и «монетных столбиков» (справа) на микроскопическом уровне после теста на стекле. Обратите внимание на потерю разграничения между эритроцитами в скоплениях, которые появляются при аутоагглютинации. 400x.}
\end{figure}

\section{Случай 8. Очень слабая гериатрическая собака}
\subsection*{Клиническая история и физикальное обследование}
12-летняя собака смешанной породы поступила в клинику после двух месяцев снижения аппетита и физической активности (рис. 1). Животное было правильно вакцинировано и дегельминтизировано, без предшествующей истории других процессов. При физикальном обследовании обнаружен лишь слегка бледный оттенок слизистых оболочек.

\begin{tabular}{@{}ll@{}}
\toprule
Ректальная температура & 38.4°C \\
ЧСС & 100 уд/мин \\
ЧД & 20 дых/мин \\
ВНК & Менее 2 сек \\
\bottomrule
\end{tabular}

\begin{figure}[htbp]
\centering
\includegraphics[width=0.5\textwidth]{placeholder.jpg}
\caption{Пациент при поступлении.}
\end{figure}

\subsection*{Гемограмма}
\begin{tabular}{lccccc}
\toprule
Тест & Результат & Реф. диапазон & LOW & NORMAL & HIGH \\
\midrule
HCT & 33\% & 37--55 & $\blacksquare$ \\
RBC & 4.78 × 10\textsuperscript{12}/L & 5.50--8.50 & $\blacksquare$ \\
Hb & 108 g/L & 120--180 & $\blacksquare$ \\
MCV & 69 fL & 60--77 & $\blacksquare$ \\
RDW & 15\% & 12.0--15.0 & $\blacksquare$ \\
MCHC & 330 g/L & 310--340 & $\blacksquare$ \\
MCH & 22 pg & 19.5--24.5 & $\blacksquare$ \\
WBC & 66.9 × 10\textsuperscript{9}/L & 6.0--17.0 & $\blacksquare$ \\
NEU & 28 × 10\textsuperscript{9}/L & 3.0--12.0 & $\blacksquare$ \\
LYM & 11.9 × 10\textsuperscript{9}/L & 1--4.80 & $\blacksquare$ \\
MONO & 22.1 × 10\textsuperscript{9}/L & 0.20--1.50 & $\blacksquare$ \\
EOS & 1.9 × 10\textsuperscript{9}/L & 0--0.80 & $\blacksquare$ \\
BASO & 2.9 × 10\textsuperscript{9}/L & 0--0.40 & $\blacksquare$ \\
PLT & 88 × 10\textsuperscript{9}/L & 200--500 & $\blacksquare$ \\
MPV & 9.1 fL & 3.9--11.1 & $\blacksquare$ \\
\% RETIC* & 1.5\% & 0.5--1.5 & $\blacksquare$ \\
\% BAND* & 12.7\% & 0--3.5 & $\blacksquare$ \\
\bottomrule
\multicolumn{6}{l}{*Ручной подсчет}
\end{tabular}

\begin{tabular}{l}
Внимание: аномальное распределение лейкоцитов. \\
\end{tabular}

\subsection*{Вопрос: Какие изменения можно наблюдать в гемограмме?}

\subsection*{Мазки крови}
\begin{figure}[htbp]
\centering
\includegraphics[width=0.6\textwidth]{placeholder.jpg}
\caption{Мазок крови пациента. Окраска Diff-Quick. 400x.}
\end{figure}
\begin{figure}[htbp]
\centering
\includegraphics[width=0.6\textwidth]{placeholder.jpg}
\caption{Деталь мазка крови пациента. Окраска Diff-Quick. 400x.}
\end{figure}

\subsection*{Вопрос: Какие изменения наблюдались в мазке крови?}

\subsection*{Цитологическое исследование костного мозга}
\begin{figure}[htbp]
\centering
\includegraphics[width=0.6\textwidth]{placeholder.jpg}
\caption{Цитология костного мозга. Окраска Diff-Quick. 400x.}
\end{figure}

\begin{tabular}{lcc}
\toprule
Тест & Результат & Референсный диапазон/Ожидаемое значение \\
\midrule
Клеточность & 90 & ≈ 40--50\% \\
Миелоидно-эритроидное соотношение & 2.25 & 0.75--2.5 \\
\% Эритроидных бластов & 1.5 & < 3\% \\
\% Миелоидных бластов & 30.5 & < 3\% \\
\bottomrule
\end{tabular}

\subsection*{Вопрос: Учитывая гематологические и медуллярные находки, каков диагноз?}

\subsection*{Разбор случая}
\subsubsection*{Какие изменения можно наблюдать в гемограмме?}
У животного наблюдалась умеренная анемия с незначительной видимой регенерацией, а также выраженная тромбоцитопения. Лейкограмма показала тяжелый лейкоцитоз с обостренным увеличением всех субпопуляций. Эта находка требовала детального анализа мазка крови, поскольку сбой в классификации лейкоцитов автоматическим анализатором был весьма вероятен (табл. 1).

\subsubsection*{Какие изменения наблюдались в мазке крови?}
В мазке крови обнаружены многочисленные ядерные эритроциты, которые не сопровождались другими признаками регенерации (такими как тельца Хауэлла-Жолли). Полихроматофилов также было мало, что указывало на неадекватную метарубрицитемию вследствие заметного повреждения костного мозга.

На уровне лейкоцитов обнаружены циркулирующие миелоидные бласты, что всегда является атипичной находкой (рис. 5). На изображениях конкретно видны метамиелоциты и миелобласты-промиелоциты. Наряду с этими клетками также наблюдались многочисленные сегментоядерные нейтрофилы, палочкоядерные нейтрофилы и моноциты с типичной морфологией. Учитывая большой размер и атипичную морфологию миелоидных бластов, высокие показатели лимфоцитов, моноцитов и даже эозинофилов и базофилов можно объяснить неправильной классификацией этих атипичных клеток автоматическим анализатором.

Важно отметить, что присутствие циркулирующих миелоидных клеток-предшественников не является исключительным для лейкемических состояний, поскольку этот профиль также может наблюдаться при так называемых лейкемоидных реакциях (которые, по-видимому, вызваны очень тяжелыми воспалительными лейкограммами). Чтобы отличить лейкемию от лейкемоидной реакции, необходимо провести исследование костного мозга и исключить инфекционные или воспалительные причины; также полезно использование серийных гематологических тестов.

Наконец, поскольку тромбоцитарных агрегатов обнаружено не было, тромбоцитопения была классифицирована как наиболее вероятно истинная.

\begin{table}[htbp]
\centering
\caption{Основные причины сбоев в подсчете и классификации гематических клеток автоматическими анализаторами.}
\begin{tabular}{p{14cm}}
\toprule
Наиболее частые причины сбоев автоматического подсчета или классификации WBC \\
\midrule
• Наличие атипичных лейкоцитарных клеток (бласты или тяжелый левый сдвиг и т.д.). \\
• Наличие ядерных эритроцитов (часто ошибочно принимаются за лимфоциты). \\
• Старение образца (лизис лейкоцитов или «набухание» клеток и т.д.). \\
\midrule
Наиболее частые причины сбоев автоматического подсчета RBC и тромбоцитов \\
\midrule
• Ложные подсчеты тромбоцитов из-за тромбоцитарных агрегатов, наличия телец Гейнца или тяжелого гемолиза. \\
• Наличие макротромбоцитов. \\
• Старение образца («набухание» клеток или лизис лейкоцитов и т.д.). \\
\bottomrule
\end{tabular}
\end{table}

\begin{figure}[htbp]
\centering
\includegraphics[width=0.6\textwidth]{placeholder.jpg}
\caption{Пункция костного мозга показала наличие многочисленных миелоидных бластов (черные стрелки). Окраска Diff-Quick. 400x.}
\end{figure}
\begin{figure}[htbp]
\centering
\includegraphics[width=0.8\textwidth]{placeholder.jpg}
\caption{Различные тромбоцитарные, эритроидные и миелоидные популяции-предшественники, присутствующие в цитологическом анализе костного мозга. Верхний ряд, тромбоцитарные и эритроидные линии, (a) мегакариоцит; (b) рубрибласт; (c) прорубрицит; (d) рубрицит; и (e) метарубрицит. Нижний ряд, миелоидные линии, (f) промиелоцит; (g) миелоцит; (h) метамиелоцит; (i) палочкоядерный нейтрофил; и (j) сегментоядерный нейтрофил. Окраска Diff-Quick. 1000x.}
\end{figure}

\section{Случай 9. Угнетенная сука с желтоватыми слизистыми оболочками}
\subsection*{Клиническая история и физикальное обследование}
7-летняя нестерилизованная сука золотистого ретривера поступила с рвотой, апатией, непереносимостью физических нагрузок и полиурией (рис. 1). Пациент была правильно вакцинирована и дегельминтизирована, без предшествующей истории других процессов. Владельцы лечили собаку в течение недели парацетамолом (с переменной и плохо определенной дозой), полагая, что у животного «простуда». У животного были слегка бледные и желтушные слизистые оболочки без снижения тургора кожи.

\begin{tabular}{@{}ll@{}}
\toprule
Ректальная температура & 39.5°C \\
ЧСС & 170 уд/мин \\
ЧД & 50 дых/мин \\
ВНК & Менее 2 сек \\
\bottomrule
\end{tabular}

\begin{figure}[htbp]
\centering
\includegraphics[width=0.5\textwidth]{placeholder.jpg}
\caption{Пациент при поступлении.}
\end{figure}

\subsection*{Гемограмма}
\begin{tabular}{lccccc}
\toprule
Тест & Результат & Реф. диапазон & LOW & NORMAL & HIGH \\
\midrule
HCT & 33\% & 37--55 & $\blacksquare$ \\
RBC & 4.71 × 10\textsuperscript{12}/L & 5.50--8.50 & $\blacksquare$ \\
Hb & 181 g/L & 120--180 & $\blacksquare$ \\
MCV & 78 fL & 60--77 & $\blacksquare$ \\
RDW & 21\% & 12.0--15.0 & $\blacksquare$ \\
MCHC & 540 g/L & 310--340 & $\blacksquare$ \\
MCH & 38 pg & 19.5--24.5 & $\blacksquare$ \\
WBC & 5.2 × 10\textsuperscript{9}/L & 6.0--17.0 & $\blacksquare$ \\
NEU & 2.3 × 10\textsuperscript{9}/L & 3.0--12.0 & $\blacksquare$ \\
LYM & 2.61 × 10\textsuperscript{9}/L & 1--4.80 & $\blacksquare$ \\
MONO & 0.20 × 10\textsuperscript{9}/L & 0.20--1.50 & $\blacksquare$ \\
EOS & 0.14 × 10\textsuperscript{9}/L & 0--0.80 & $\blacksquare$ \\
BASO & 0 × 10\textsuperscript{9}/L & 0--0.40 & $\blacksquare$ \\
PLT & 1,000 × 10\textsuperscript{9}/L & 200--500 & $\blacksquare$ \\
MPV & 2.56 fL & 3.9--11.1 & $\blacksquare$ \\
\% RETIC* & 1.9\% & 0.5--1.5 & $\blacksquare$ \\
\% BAND* & 28\% & 0--3.5 & $\blacksquare$ \\
\bottomrule
\multicolumn{6}{l}{*Ручной подсчет}
\end{tabular}

\subsection*{Вопрос: Какие изменения можно наблюдать в гемограмме?}

\subsection*{Мазки крови}
\begin{figure}[htbp]
\centering
\includegraphics[width=0.6\textwidth]{placeholder.jpg}
\caption{Деталь мазка крови пациента. Окраска Diff-Quik, 1000x.}
\end{figure}
\begin{figure}[htbp]
\centering
\includegraphics[width=0.6\textwidth]{placeholder.jpg}
\caption{Деталь мазка крови пациента. Окраска Diff-Quik, 1000x.}
\end{figure}
\begin{figure}[htbp]
\centering
\includegraphics[width=0.6\textwidth]{placeholder.jpg}
\caption{Мазок крови животного. Окраска Diff-Quik, 400x.}
\end{figure}
\begin{figure}[htbp]
\centering
\includegraphics[width=0.6\textwidth]{placeholder.jpg}
\caption{Деталь мазка крови животного. Окраска Diff-Quik, 400x.}
\end{figure}

\subsection*{Вопрос: Какие изменения наблюдались в мазке крови?}
\subsection*{Вопрос: Учитывая анамнез и гематологические находки, каков предположительный диагноз?}

\subsection*{Разбор случая}
\textbf{Пиометра с признаками токсичности и гемолизом с тельцами Гейнца}

\subsubsection*{Какие изменения в гемограмме?}
У пациентки умеренная регенераторная гемолитическая анемия. Лейкограмма показала лейкоцитопению со зрелой нейтропенией и увеличенным количеством палочкоядерных нейтрофилов (дегенеративный левый сдвиг). Этот тип ответа типичен для тяжелых воспалительных или инфекционных процессов (табл. 1). Также наблюдался тромбоцитоз, хотя низкий объем тромбоцитов, казалось, указывал на то, что эта находка была артефактом.

\subsubsection*{Какие изменения в мазке крови?}
В мазке крови обнаружены многочисленные тельца Гейнца (эритроцитарные цилиндрические выпячивания; рис. 6), эксцентроциты (с прозрачной областью гемоглобина на одной стороне цитоплазмы) и шистоциты (рис. 3). Все эти находки указывают на наличие окислительного повреждения эритроцитов, вызывающего фрагментацию и гемолитическую анемию.

На лейкоцитарном уровне наблюдалось большое количество палочкоядерных нейтрофилов, которые в некоторых полях преобладали над сегментоядерными (рис. 4). В нейтрофилах наблюдались выраженные признаки токсичности, такие как цитоплазматическая вакуолизация, тельца Дёле, цитоплазматическая базофилия и токсическая зернистость (рис. 5 и 7).

\paragraph{На заметку}
Основными причинами лейкоцитопении (нейтропении) являются тяжелые бактериальные инфекции, медуллярная токсичность (химиотерапия, эстрогены и т.д.), вирусные инфекции (FeLV, FIV, парвовирус, эрлихиоз и т.д.) или гематические новообразования.

\begin{figure}[htbp]
\centering
\includegraphics[width=0.6\textwidth]{placeholder.jpg}
\caption{Тельца Гейнца в эритроцитах собаки. Окраска Diff-Quik, 400x.}
\end{figure}
\begin{figure}[htbp]
\centering
\includegraphics[width=0.8\textwidth]{placeholder.jpg}
\caption{Примеры токсичности в сегментоядерных и палочкоядерных нейтрофилах: базофилия и цитоплазматическая вакуолизация в палочкоядерных нейтрофилах (a и b); тельца Дёле (c); и токсическая зернистость (d).}
\end{figure}

\begin{table}[htbp]
\centering
\caption{Тяжесть миелоидного ответа при воспалительных лейкограммах.}
\begin{tabular}{lccccc}
\toprule
Тяжесть & Паттерн & Общие лейкоциты & Сегментоядерные нейтрофилы & Палочкоядерные нейтрофилы \\
\midrule
+ & Регенераторный левый сдвиг & ↑ & ↑ & ↑ \\
++ & Дегенеративный левый сдвиг & ↑ & ↓ & ↑ \\
+++ & Лейкоцитопения & ↓ & ↓ & ↓ \\
\bottomrule
\multicolumn{5}{l}{Все паттерны токсических признаков серьезны, но еще более серьезны для палочкоядерных нейтрофилов.} \\
\multicolumn{5}{l}{Тяжесть миелоидного ответа может не быть синонимом тяжести состояния.} \\
\end{tabular}
\end{table}

\subsubsection*{Учитывая анамнез и гематологические данные, каков предположительный диагноз?}
Наличие воспалительной лейкограммы с дегенеративным левым сдвигом и признаками токсичности указывает на тяжелую воспалительную/инфекционную клиническую картину (табл. 1). В случае этой суки эта находка была вызвана запущенной пиометрой (рис. 8). Признаки окислительного повреждения и фрагментации эритроцитов могли быть связаны с токсичностью парацетамола/ацетаминофена. Наличие многочисленных телец Гейнца также часто ассоциируется с появлением тромбоцитоза (артефактного) из-за ошибки распознавания автоматическим анализатором.

\begin{figure}[htbp]
\centering
\includegraphics[width=0.5\textwidth]{placeholder.jpg}
\caption{Матка пациентки после экстренного вмешательства.}
\end{figure}

\section{Случай 10. Собака с лихорадочными пиками}
\subsection*{Клиническая история и физикальное обследование}
5-летний кобель из сельской местности поступил в клинику с одышкой (смешанного типа), потерей веса и аппетита, вялостью (рис. 1). После выявления значительного обезвоживания пациент был госпитализирован. Во время пребывания отмечались лихорадочные пики (два-три в день), генерализованная лимфаденомегалия и бледные, слегка желтушные слизистые оболочки. Животное не было вакцинировано или дегельминтизировано в течение нескольких лет.

\begin{tabular}{@{}ll@{}}
\toprule
Ректальная температура & 41.2°C \\
ЧСС & 170 уд/мин \\
ЧД & Выраженная одышка \\
ВНК & Менее 2 сек \\
\bottomrule
\end{tabular}

\begin{figure}[htbp]
\centering
\includegraphics[width=0.5\textwidth]{placeholder.jpg}
\caption{Пациент в домашней обстановке.}
\end{figure}

\subsection*{Гемограмма}
\begin{tabular}{lccccc}
\toprule
Тест & Результат & Реф. диапазон & LOW & NORMAL & HIGH \\
\midrule
HCT & 30\% & 37--55 & $\blacksquare$ \\
RBC & 3.71 × 10\textsuperscript{12}/L & 5.50--8.50 & $\blacksquare$ \\
Hb & 195 g/L & 120--180 & $\blacksquare$ \\
MCV & 80 fL & 60--77 & $\blacksquare$ \\
RDW & 16\% & 12.0--15.0 & $\blacksquare$ \\
MCHC & 640 g/L & 310--340 & $\blacksquare$ \\
MCH & 52 pg & 19.5--24.5 & $\blacksquare$ \\
WBC & 9.1 × 10\textsuperscript{9}/L & 6.0--17.0 & $\blacksquare$ \\
NEU & 2.2 × 10\textsuperscript{9}/L & 3.0--12.0 & $\blacksquare$ \\
LYM & 5.1 k/L & 1--4.80 & $\blacksquare$ \\
MONO & 0.91 × 10\textsuperscript{9}/L & 0.20--1.50 & $\blacksquare$ \\
EOS & 0.7 × 10\textsuperscript{9}/L & 0--0.80 & $\blacksquare$ \\
BASO & 0 × 10\textsuperscript{9}/L & 0--0.40 & $\blacksquare$ \\
PLT & 30 × 10\textsuperscript{9}/L & 200--500 & $\blacksquare$ \\
MPV & 9.5 fL & 3.9--11.1 & $\blacksquare$ \\
\% RETIC* & 4\% & 0.5--1.5 & $\blacksquare$ \\
\% BAND* & 1.6\% & 0--3.5 & $\blacksquare$ \\
\bottomrule
\multicolumn{6}{l}{*Ручной подсчет}
\end{tabular}

\subsection*{Вопрос: Какие изменения в гемограмме?}

\subsection*{Мазки крови}
\begin{figure}[htbp]
\centering
\includegraphics[width=0.6\textwidth]{placeholder.jpg}
\caption{Мазок крови пациента. Окраска Diff-Quik, 400x.}
\end{figure}

\subsection*{Специальная окраска крови (Гимза)}
\begin{figure}[htbp]
\centering
\includegraphics[width=0.6\textwidth]{placeholder.jpg}
\caption{Мазок капиллярной крови (взятой из ушной раковины) животного. Окраска Гимза, 1000x.}
\end{figure}

\subsection*{Вопрос: Какие структуры можно увидеть с помощью этой окраски?}
\subsection*{Вопрос: Учитывая клиническую картину и гематологические находки, каков диагноз?}

\subsection*{Разбор случая}
\subsubsection*{Какие изменения в гемограмме?}
У пациента гемолитическая анемия (снижение HCT и RBC при повышенном Hb), что также объясняло наличие бледных и иктеричных слизистых оболочек. Этот гемолиз также вызвал повышение параметров MCHC и MCH. Регенерация у этого животного была высокой (как обычно бывает при гемолитических анемиях).

Что касается лейкоцитов, наблюдалась умеренная нейтропения с легким лимфоцитозом без увеличения количества палочкоядерных нейтрофилов. Наконец, у животного была тяжелая тромбоцитопения.

\subsubsection*{Какие изменения в мазке крови?}
В мазке крови обнаружены ядерные эритроциты (регенерация), а также подтверждено наличие многочисленных «раздавленных» клеток (smudge cells) (рис. 4), разрушенных клеток, классификация которых в различные лейкоцитарные субпопуляции затруднена. Среди оставшихся лейкоцитов наблюдались нейтрофилы (гиперсегментированные) и моноциты.

\paragraph{На заметку}
Smudge cells — это лейкоциты с разорванным ядром или цитоплазмой, что затрудняет их классификацию автоматическими анализаторами. Если более $10\%$ этих клеток обнаруживается при дифференциальном подсчете, мазок крови следует повторить для получения надежных процентов. Основными причинами появления smudge cells являются неадекватная техника (приложение чрезмерного давления при распределении образца или неправильное распределение образца) и слабость лейкоцитов (наличие неопластических клеток или неправильное хранение образца).

\begin{figure}[htbp]
\centering
\includegraphics[width=0.6\textwidth]{placeholder.jpg}
\caption{Smudge cells (стрелки). Не следует пытаться распознать или классифицировать эти клетки при ручном подсчете. Окраска Diff-Quik, 400x.}
\end{figure}

\subsubsection*{Какие структуры видны с помощью этой окраски?}
В многочисленных эритроцитах образца наблюдались четкие грушевидные структуры, окруженные пурпурным окрашиванием (рис. 2). Эти структуры соответствуют мерозоитам \textit{Babesia canis}. Обратите внимание, что эти организмы распознаются в мазке крови только с помощью этого типа окраски. Кроме того, их присутствие непостоянно, и они часто легче идентифицируются в капиллярной крови, чем в венозной.

\subsubsection*{Учитывая клиническую картину и гематологические данные, каков диагноз?}
Клинические симптомы, среда, в которой жило животное (предрасположенность к паразитированию клещами), и гематологические данные (гемолитическая анемия с легкой нейтропенией и выраженной тромбоцитопенией) должны привести к клиническому подозрению на бабезиоз. Хотя у этого животного в крови, окрашенной по Гимзе, наблюдались мерозоиты, окончательный диагноз обычно ставится с помощью таких тестов, как ELISA или ПЦР и т.д. На рисунке 5 показаны другие гемопаразиты, которые можно наблюдать в крови собак.

\begin{figure}[htbp]
\centering
\includegraphics[width=0.8\textwidth]{placeholder.jpg}
\caption{Различные гемопаразиты собак: (a) морула \textit{Ehrlichia canis} в лимфоците; (b) прогрессирующие формы \textit{Mycoplasma haemocanis}, образующие типичную внутриэритроцитарную цепочку кокков; и (c) \textit{Hepatozoon canis} с овоидной морфологией и без окрашивания внутри цитоплазмы нейтрофила. Окраска Diff-Quik, 1000x.}
\end{figure}

\section{Случай 11. Щенок с респираторными проблемами}
\subsection*{Клиническая история и физикальное обследование}
2-месячный щенок французского бульдога поступил на прием с неврологическими симптомами (в анамнезе судороги 2 дня назад), смешанной одышкой и легким слизисто-гнойным выделением из носа (рис. 1). Животное еще не было вакцинировано или дегельминтизировано. Физикальное обследование выявило легкую спутанность сознания без клинических признаков анемии или обезвоживания.

\begin{tabular}{@{}ll@{}}
\toprule
Ректальная температура & 40°C \\
ЧСС & 180 уд/мин \\
ЧД & Выраженная одышка \\
ВНК & Менее 2 сек \\
\bottomrule
\end{tabular}

\begin{figure}[htbp]
\centering
\includegraphics[width=0.5\textwidth]{placeholder.jpg}
\caption{Пациент при поступлении.}
\end{figure}

\subsection*{Гемограмма}
\begin{tabular}{lccccc}
\toprule
Тест & Результат & Реф. диапазон & LOW & NORMAL & HIGH \\
\midrule
HCT & 34\% & 37--55 & $\blacksquare$ \\
RBC & 4.91 × 10\textsuperscript{12}/L & 5.50--8.50 & $\blacksquare$ \\
Hb & 115 g/L & 120--180 & $\blacksquare$ \\
MCV & 72 fL & 60--77 & $\blacksquare$ \\
RDW & 14\% & 12.0--15.0 & $\blacksquare$ \\
MCHC & 330 g/L & 310--340 & $\blacksquare$ \\
MCH & 23 pg & 19.5--24.5 & $\blacksquare$ \\
WBC & 30.1 × 10\textsuperscript{9}/L & 6.0--17.0 & $\blacksquare$ \\
NEU & 25.1 × 10\textsuperscript{9}/L & 3.0--12.0 & $\blacksquare$ \\
LYM & 0.39 × 10\textsuperscript{9}/L & 1--4.80 & $\blacksquare$ \\
MONO & 3.91 × 10\textsuperscript{9}/L & 0.20--1.50 & $\blacksquare$ \\
EOS & 0.7 × 10\textsuperscript{9}/L & 0--0.8 & $\blacksquare$ \\
BASO & 0 × 10\textsuperscript{9}/L & 0--0.4 & $\blacksquare$ \\
PLT & 200 × 10\textsuperscript{9}/L & 200--500 & $\blacksquare$ \\
MPV & 9.5 fL & 3.9--11.1 & $\blacksquare$ \\
\% RETIC* & 2\% & 0.5--1.5 & $\blacksquare$ \\
\% BAND* & 9.5\% & 0--3.5 & $\blacksquare$ \\
\bottomrule
\multicolumn{6}{l}{*Ручной подсчет}
\end{tabular}

\subsection*{Вопрос: Какие изменения в гемограмме?}

\subsection*{Мазки крови}
\begin{figure}[htbp]
\centering
\includegraphics[width=0.6\textwidth]{placeholder.jpg}
\caption{Мазок крови щенка. Окраска Diff-Quik, 400x.}
\end{figure}
\begin{figure}[htbp]
\centering
\includegraphics[width=0.6\textwidth]{placeholder.jpg}
\caption{Деталь мазка крови пациента. Окраска Diff-Quik, 1000x.}
\end{figure}
\begin{figure}[htbp]
\centering
\includegraphics[width=0.6\textwidth]{placeholder.jpg}
\caption{Деталь мазка крови пациента. Окраска Diff-Quik, 1000x.}
\end{figure}

\subsection*{Вопрос: Какие изменения в мазке крови?}
\subsection*{Вопрос: Учитывая клиническую картину и гематологические данные, каков диагноз?}

\subsection*{Разбор случая}
\subsubsection*{Какие изменения в гемограмме?}
Мы должны учитывать, что пациент — щенок, и гематологические референсные диапазоны отличаются для животных в возрасте до 2 месяцев (рис. 7 из случая 1). С учетом этой переменной не наблюдалось отклонений ни в красной крови (значения HCT, RBC, Hb и процента ретикулоцитов в норме для этого возраста), ни в тромбоцитах.

Наблюдался выраженный лейкоцитоз с нейтрофилией, увеличением количества палочкоядерных нейтрофилов и моноцитозом, что является типичной находкой в лейкограмме при воспалении. Также присутствовала выраженная лимфопения.

\subsubsection*{Какие изменения в мазке крови?}
В мазке крови обнаружено множество моноцитов, палочкоядерных нейтрофилов и в основном сегментоядерных нейтрофилов (рис. 2), хотя эти увеличения находились в пределах нормы для всех клеточных популяций. Что касается эритроцитов (рис. 5), в них было множество телец-включений от округлых до овальных, разного размера (от четверти до пятой части диаметра эритроцита), с розоватым окрашиванием, соответствующих тельцам-включениям чумы собак. Важно не путать эти включения с агрегатами красителя или остатками, либо с тельцами Хауэлла-Жолли (темно-синие базофильные округлые образования меньшего размера).

Те же тельца-включения чумы собак можно было увидеть и в лейкоцитах, таких как нейтрофилы (рис. 3). В этом случае важно не путать эти включения с токсической зернистостью (такой как множественные, гораздо меньшие и вариабельные по морфологии гранулы, видимые в нейтрофиле, показанном на рисунке).

\begin{figure}[htbp]
\centering
\includegraphics[width=0.6\textwidth]{placeholder.jpg}
\caption{Ацидофильные тельца-включения чумы собак (синие стрелки) против телец Хауэлла-Жолли (черные стрелки). Окраска Diff-Quik, 1000x.}
\end{figure}

\subsubsection*{Учитывая клиническую картину и гематологические данные, каков диагноз?}
Анамнез, возраст пациента и клиническое состояние должны вызвать подозрение на чуму собак. На гематологическом уровне наличие лимфопении типично для этой патологии (тромбоцитопения более вариабельна); хотя, учитывая частоту, с которой обычно появляются вторичные бактериальные инфекции (особенно при поражении органов дыхания или пищеварения), также часто наблюдаются воспалительные лейкограммы. После исключения других аномалий или артефактов с похожей картиной, наличие ацидофильных телец-включений в эритроцитах и/или лейкоцитах является диагностическим признаком чумы, хотя они обычно появляются только через 2--9 дней после заражения и не всегда совпадают с наличием клинических признаков (которые обычно появляются позже).

\section{Случай 12. Молодая собака с подкожными кровоизлияниями}
\subsection*{Клиническая история и физикальное обследование}
9-месячная собака смешанной породы поступила в клинику из-за отсутствия аппетита и потери веса, развивавшихся в течение нескольких месяцев, а также гематурии и наличия свежей крови в кале (рис. 1). Физикальное обследование выявило выраженную бледность слизистых оболочек и наличие подкожных и слизистых кровоизлияний. Место венепункции также кровоточило аномально долго.

\begin{figure}[htbp]
\centering
\includegraphics[width=0.5\textwidth]{placeholder.jpg}
\caption{Пациент в зоне госпитализации.}
\end{figure}

\subsection*{Гемограмма}
\begin{tabular}{lccccc}
\toprule
Тест & Результат & Реф. диапазон & LOW & NORMAL & HIGH \\
\midrule
HCT & 28\% & 37--55 & $\blacksquare$ \\
RBC & 4.97 × 10\textsuperscript{12}/L & 5.50--8.50 & $\blacksquare$ \\
Hb & 85 g/L & 120--180 & $\blacksquare$ \\
MCV & 52 fL & 60--77 & $\blacksquare$ \\
RDW & 18\% & 12.0--15.0 & $\blacksquare$ \\
MCHC & 300 g/L & 310--340 & $\blacksquare$ \\
MCH & 17 pg & 19.5--24.5 & $\blacksquare$ \\
WBC & 9.2 × 10\textsuperscript{9}/L & 6.0--17.0 & $\blacksquare$ \\
NEU & 7.15 × 10\textsuperscript{9}/L & 3.0--12.0 & $\blacksquare$ \\
LYM & 1.77 × 10\textsuperscript{9}/L & 1--4.80 & $\blacksquare$ \\
MONO & 0.21 × 10\textsuperscript{9}/L & 0.20--1.50 & $\blacksquare$ \\
EOS & 0.12 × 10\textsuperscript{9}/L & 0--0.80 & $\blacksquare$ \\
BASO & 0 × 10\textsuperscript{9}/L & 0--0.4 & $\blacksquare$ \\
PLT & 600 × 10\textsuperscript{9}/L & 200--500 & $\blacksquare$ \\
MPV & 12.5 fL & 3.90--11.1 & $\blacksquare$ \\
\% RETIC* & 0\% & 0.5--1.5 & $\blacksquare$ \\
\% BAND* & 2.7\% & 0--3.5 & $\blacksquare$ \\
\bottomrule
\multicolumn{6}{l}{*Ручной подсчет}
\end{tabular}

\subsection*{Вопрос: Какие изменения в гемограмме?}

\subsection*{Мазки крови}
\begin{figure}[htbp]
\centering
\includegraphics[width=0.6\textwidth]{placeholder.jpg}
\caption{Мазок крови пациента. Окраска Diff-Quik, 400x.}
\end{figure}
\begin{figure}[htbp]
\centering
\includegraphics[width=0.6\textwidth]{placeholder.jpg}
\caption{Деталь мазка крови пациента. Окраска Diff-Quik, 1000x.}
\end{figure}

\subsection*{Вопрос: Какие изменения в мазке крови?}

\subsection*{Гемостатическое исследование}
\begin{tabular}{lcc}
\toprule
Тест & Результат & Референсный диапазон \\
\midrule
Прогромбиновое время & 9 секунд & 6--12 \\
Активированное частичное тромбопластиновое время & 15 секунд & 12--22 \\
Время кровотечения слизистой оболочки щеки & 8.5 минут & < 4 \\
D-димеры & 100 ng/mL & < 250 \\
\bottomrule
\end{tabular}

\subsection*{Вопрос: Учитывая гемостатические и гематологические данные, каков диагноз?}

\subsection*{Разбор случая}
\subsubsection*{Какие изменения в гемограмме?}
У пациента наблюдалась заметная нерегенераторная анемия (снижение HCT, RBC, Hb и процента ретикулоцитов), причем поразительно, что общее количество эритроцитов снизилось меньше, чем другие показатели. Была заметная вариабельность размеров эритроцитов (RDW), сопровождающаяся уменьшением общего размера эритроцитов (низкий MCV) и эритроцитами с заметным гипохромным характером (снижение MCHC). Этот паттерн соответствовал гипохромной микроцитарной анемии, типичной для симптомов дефицита железа.

Лейкоцитарная популяция была без изменений, и наблюдалась выраженная тромбоцитопения с увеличенным средним размером тромбоцитов (MPV), типичный ответ при реактивном тромбоцитозе после продолжающегося кровотечения (случай 3).

\subsubsection*{Какие изменения в мазке крови?}
В мазке крови мы наблюдали, что большинство эритроцитов были маленькими (микроцитоз), и преобладающая гипохромия была гораздо легче обнаружима (обратите внимание, как во многих случаях центральный ореол занимал более половины диаметра эритроцита, оттесняя гемоглобин к небольшому периферическому ободку; рис. 2).

Наблюдались многочисленные макротромбоциты, по размеру сходные с эритроцитами или крупнее их (рис. 4). Наличие макротромбоцитов обычно указывает на выраженную регенерацию тромбоцитов на уровне костного мозга (после исключения типичной клинической картины у кавалер-кинг-чарльз-спаниелей).

\begin{figure}[htbp]
\centering
\includegraphics[width=0.6\textwidth]{placeholder.jpg}
\caption{Макротромбоциты (черные стрелки) по сравнению с нормальными тромбоцитами (синие стрелки). Окраска Diff-Quik, 1000x.}
\end{figure}

\subsubsection*{Учитывая гемостатические и гематологические данные, каков диагноз?}
Тесты гемостаза показали изменения только в первичном гемостазе (с участием тромбоцитов и фактора фон Виллебранда). Поскольку количество тромбоцитов было нормальным, этому пациенту был поставлен диагноз тромбоцитопатии или болезни фон Виллебранда, и для окончательного диагноза потребовались другие тесты (агрегометрия тромбоцитов, определение концентрации фактора фон Виллебранда и полимеризации и т.д.). Эти типы гемостатических нарушений вызывают только удлинение времени кровотечения из слизистой оболочки щеки, в то время как другие коагуляционные тесты остаются нормальными (табл. 2 из случая 6).

Значительная железодефицитная анемия, имевшаяся у пациента, скорее всего, была вызвана хроническим кровотечением из-за тромбоцитопатии/болезни фон Виллебранда (табл. 4 из случая 1).

\section{Случай 13. Пожилая сука с полидипсией и полиурией}
\subsection*{Клиническая история и физикальное обследование}
16-летняя нестерилизованная сука смешанной породы поступила в клинику с полиурией и постоянной полидипсией в анамнезе (в течение как минимум трех месяцев), очень заметным повышением аппетита и выраженной прибавкой в весе (рис. 1). Владелец также сообщил об апатии и непереносимости физических нагрузок. Пациентка была правильно вакцинирована и дегельминтизирована, в анамнезе — атопический дерматит с постоянным лечением кортикостероидами (дозу владелец не помнил). Слизистые оболочки были нормальными, без снижения тургора кожи.

\begin{tabular}{@{}ll@{}}
\toprule
Ректальная температура & 39.6°C \\
ЧСС & 80 уд/мин \\
ЧД & 20 дых/мин \\
ВНК & Менее 2 сек \\
\bottomrule
\end{tabular}

\begin{figure}[htbp]
\centering
\includegraphics[width=0.5\textwidth]{placeholder.jpg}
\caption{Пациентка при поступлении.}
\end{figure}

\subsection*{Гемограмма}
\begin{tabular}{lccccc}
\toprule
Тест & Результат & Реф. диапазон & LOW & NORMAL & HIGH \\
\midrule
HCT & 42\% & 37--55 & $\blacksquare$ \\
RBC & 6.01 × 10\textsuperscript{12}/L & 5.50--8.50 & $\blacksquare$ \\
Hb & 140 g/L & 120--180 & $\blacksquare$ \\
MCV & 70 fL & 60--77 & $\blacksquare$ \\
RDW & 15\% & 12.0--15.0 & $\blacksquare$ \\
MCHC & 330 g/L & 310--340 & $\blacksquare$ \\
MCH & 23 pg & 19.5--24.5 & $\blacksquare$ \\
WBC & 24.2 × 10\textsuperscript{9}/L & 6.0--17.0 & $\blacksquare$ \\
NEU & 20.2 × 10\textsuperscript{9}/L & 3.0--12.0 & $\blacksquare$ \\
LYM & 0.81 × 10\textsuperscript{9}/L & 1--4.80 & $\blacksquare$ \\
MONO & 3.20 × 10\textsuperscript{9}/L & 0.20--1.50 & $\blacksquare$ \\
EOS & 0 × 10\textsuperscript{9}/L & 0--0.80 & $\blacksquare$ \\
BASO & 0 × 10\textsuperscript{9}/L & 0--0.4 & $\blacksquare$ \\
PLT & 321 × 10\textsuperscript{9}/L & 200--500 & $\blacksquare$ \\
MPV & 8.91 fL & 3.9--11.1 & $\blacksquare$ \\
\% RETIC* & 0.9\% & 0.5--1.5 & $\blacksquare$ \\
\% BAND* & 0.6\% & 0--3.5 & $\blacksquare$ \\
\bottomrule
\multicolumn{6}{l}{*Ручной подсчет}
\end{tabular}

\subsection*{Вопрос: Какие изменения в гемограмме?}

\subsection*{Мазки крови}
\begin{figure}[htbp]
\centering
\includegraphics[width=0.6\textwidth]{placeholder.jpg}
\caption{Деталь мазка крови пациентки. Окраска Diff-Quik, 1000x.}
\end{figure}
\begin{figure}[htbp]
\centering
\includegraphics[width=0.6\textwidth]{placeholder.jpg}
\caption{Деталь мазка крови пациентки. Окраска Diff-Quik, 1000x.}
\end{figure}
\begin{figure}[htbp]
\centering
\includegraphics[width=0.6\textwidth]{placeholder.jpg}
\caption{Мазок крови животного. Окраска Diff-Quik, 400x.}
\end{figure}
\begin{figure}[htbp]
\centering
\includegraphics[width=0.6\textwidth]{placeholder.jpg}
\caption{Деталь мазка крови. Окраска Diff-Quik, 1000x.}
\end{figure}

\subsection*{Вопрос: Какие изменения наблюдались в мазке крови?}
\subsection*{Вопрос: Каково значение структур, отмеченных стрелкой?}
\subsection*{Вопрос: Учитывая анамнез и гематологические данные, каков предположительный диагноз?}

\subsection*{Разбор случая}
\subsubsection*{Какие изменения в гемограмме?}
У пациентки не было отклонений в эритрограмме или тромбоцитарных показателях. Наблюдался выраженный лейкоцитоз со зрелой нейтрофилией (без увеличения палочкоядерных нейтрофилов), моноцитозом, лимфопенией и эозинопенией. Этот тип лейкограммного ответа типичен для стресса (табл. 1 из случая 7) и обычно ассоциируется с повышенным уровнем эндогенных или экзогенных кортикостероидов (табл. 1).

\subsubsection*{Какие изменения в мазке крови?}
В мазке крови обнаружены многочисленные эритроциты с морфологией клеток-мишеней, известные как кодоциты или лептоциты (рис. 5). Эти клетки, которые в норме видны только у собак, обычно появляются в контексте гиперхолестеринемии, например, при эндокринопатиях или заболеваниях печени. Реже эта морфология может также приобретаться в результате дефицита железа (микроцитарная) или при наличии полихроматофильных эритроцитов.

\begin{table}[htbp]
\centering
\caption{Соображения при интерпретации стрессовой лейкограммы у мелких животных.}
\begin{tabular}{p{14cm}}
\toprule
Соображения относительно кортикостероидов и результатов стрессовой лейкограммы \\
\midrule
• Они вызваны клиническими картинами, связанными со значительной болью или стрессом, повышающими уровень эндогенных кортикостероидов. \\
• Введение экзогенных кортикостероидов: \\
  \quad • пик ответа = через 8 часов после введения (варьирует в зависимости от препарата и пути введения). \\
  \quad • Изменения исчезают в течение 24 часов после однократного введения (но могут сохраняться до нескольких недель после длительного лечения). \\
• Лейкоцитоз обычно < 25,000/мкл. \\
• Моноцитоз более вариабелен у кошек. \\
\bottomrule
\end{tabular}
\end{table}

\begin{figure}[htbp]
\centering
\includegraphics[width=0.4\textwidth]{placeholder.jpg}
\caption{Клетки-мишени с типичным центральным скоплением гемоглобина, окруженным выраженным беловатым ореолом. Окраска Diff-Quik, 1000x.}
\end{figure}

Что касается лейкоцитов, большинство нейтрофилов были зрелыми (не палочкоядерными), причем наиболее распространены были нейтрофилы с более чем 5 дольками (сдвиг вправо, указывающий на старение клеток \textit{in vivo}) (рис. 3, 4 и 5). Моноциты, присутствующие в мазке, не показали никаких морфологических изменений.

\subsubsection*{Каково значение структур, отмеченных стрелкой?}
Эти структуры представляют собой тельца Барра и соответствуют компактированному хроматину второй Х-хромосомы у самок. Это тельце обычно имеет форму барабанной палочки и не имеет патологического значения, указывая лишь на то, что пациент женского пола.

\subsubsection*{Учитывая анамнез и гематологические данные, каков предположительный диагноз?}
Наличие выраженного профиля стрессовой лейкограммы, наряду с симптомами (полифагия, полиурия и полидипсия) и анамнезом, совместимо с диагнозом гиперадренокортицизма. Появление многочисленных клеток-мишеней указывало на весьма вероятное заболевание печени (скорее всего, поражение печени, индуцированное кортикостероидами). Учитывая предшествующую историю приема кортикостероидов, гиперадренокортицизм у этой пациентки был, скорее всего, ятрогенным.

\section{Случай 14. Очень слабый новорожденный щенок без сосательного рефлекса}
\subsection*{Клиническая история и физикальное обследование}
5-дневный щенок смешанной породы, который потерял сосательный рефлекс, был вялым и постоянно скулил (рис. 1). Животное еще не было вакцинировано или дегельминтизировано. Остальные щенки в помете и мать были в хорошем состоянии. При физикальном обследовании пациент проявлял малую активность, но клинических признаков анемии или обезвоживания не было.

\begin{figure}[htbp]
\centering
\includegraphics[width=0.5\textwidth]{placeholder.jpg}
\caption{Пациент при поступлении.}
\end{figure}

\subsection*{Гемограмма}
\begin{tabular}{lccccc}
\toprule
Тест & Результат & Реф. диапазон & LOW & NORMAL & HIGH \\
\midrule
HCT & 45\% & 37--55 & $\blacksquare$ \\
RBC & 5.18 × 10\textsuperscript{12}/L & 5.50--8.50 & $\blacksquare$ \\
Hb & 155 g/L & 120--180 & $\blacksquare$ \\
MCV & 92 fL & 60--77 & $\blacksquare$ \\
RDW & 19\% & 12.0--15.0 & $\blacksquare$ \\
MCHC & 340 g/L & 310--340 & $\blacksquare$ \\
MCH & 30 pg & 19.5--24.5 & $\blacksquare$ \\
WBC & 2.5 × 10\textsuperscript{9}/L & 6.0--17.0 & $\blacksquare$ \\
NEU & 0.9 × 10\textsuperscript{9}/L & 3.0--12.0 & $\blacksquare$ \\
LYM & 1.34 × 10\textsuperscript{9}/L & 1--4.80 & $\blacksquare$ \\
MONO & 0.25 × 10\textsuperscript{9}/L & 0.20--1.50 & $\blacksquare$ \\
EOS & 0.15 × 10\textsuperscript{9}/L & 0--0.80 & $\blacksquare$ \\
BASO & 0 × 10\textsuperscript{9}/L & 0--0.4 & $\blacksquare$ \\
PLT & 200 × 10\textsuperscript{9}/L & 200--500 & $\blacksquare$ \\
MPV & 9.9 fL & 3.9--11.1 & $\blacksquare$ \\
\% RETIC* & 3\% & 0.5--1.5 & $\blacksquare$ \\
\% BAND* & 22\% & 0--3.5 & $\blacksquare$ \\
\bottomrule
\multicolumn{6}{l}{*Ручной подсчет}
\end{tabular}

\subsection*{Вопрос: Какие изменения в гемограмме?}

\subsection*{Мазки крови}
\begin{figure}[htbp]
\centering
\includegraphics[width=0.6\textwidth]{placeholder.jpg}
\caption{Мазок крови щенка. Окраска Diff-Quik, 1000x.}
\end{figure}
\begin{figure}[htbp]
\centering
\includegraphics[width=0.6\textwidth]{placeholder.jpg}
\caption{Деталь мазка крови пациента. Окраска Diff-Quik, 1000x.}
\end{figure}

\subsection*{Вопрос: Какие изменения в мазке крови?}
\subsection*{Вопрос: Учитывая клиническую картину и гематологические данные, каков диагноз? Какое лечение следует начать немедленно?}

\subsection*{Разбор случая}
\textbf{Неонатальная септицемия}
\subsubsection*{Какие изменения в гемограмме?}
Учитывая возраст новорожденного (случай 1), ни один эритроидный параметр не был изменен, при этом повышенные MCV, RDW и процент ретикулоцитов, а также низкие значения RBC являются нормальными для животного в возрасте менее одной недели.

У этого пациента наблюдалась выраженная лейкопения, вызванная тяжелой нейтропенией с дегенеративным левым сдвигом, что является находкой, типичной для чрезмерно воспалительных лейкограмм (табл. 1 из случая 9). Остальные лейкоцитарные и тромбоцитарные популяции не имели количественных изменений.

\subsubsection*{Какие изменения в мазке крови?}
В мазке крови обнаружены многочисленные преломляющие свет тельца внутри эритроцитов (рис. 4). Это артефактная находка, вызванная неправильной подготовкой мазка (когда мазку не дают достаточно высохнуть и слишком быстро окрашивают). Очень важно не путать эти преломляющие тельца с эритропаразитами, такими как \textit{Babesia} spp. (помните, что мерозоиты \textit{Babesia} spp. видны только при окраске по Гимзе, а не Diff-Quik).

На уровне лейкоцитов (рис. 3 и 5) наблюдались выраженные признаки токсичности, такие как наличие телец Дёле, цитоплазматическая базофилия и токсическая зернистость (случай 9).

\begin{figure}[htbp]
\centering
\includegraphics[width=0.6\textwidth]{placeholder.jpg}
\caption{Преломляющие внутриэритроцитарные тельца-включения из-за неадекватного высушивания образца (черные стрелки). Окраска Diff-Quik, 1000x.}
\end{figure}
\begin{figure}[htbp]
\centering
\includegraphics[width=0.6\textwidth]{placeholder.jpg}
\caption{Токсическая зернистость, присутствующая в палочкоядерном нейтрофиле (центральная клетка) пациента. Окраска Diff-Quik, 1000x.}
\end{figure}

\subsubsection*{Учитывая клиническую картину и гематологические данные, каков диагноз? Какое лечение?}
Учитывая возраст пациента, симптомы (типичные для синдрома угасающего щенка) и наличие нейтропении с признаками токсичности, у этого новорожденного предположительный диагноз — неонатальная септицемия. Окончательный диагноз потребовал бы посева крови и использования других методов для поиска первичного очага инфекции. Что касается лечения, учитывая нейтропению, следует начать антибактериальную терапию широкого спектра действия для обеспечения промежуточного покрытия.

\paragraph{На заметку}
Учитывая высокий риск развития оппортунистических инфекций, антибактериальную терапию широкого спектра действия следует начинать у каждого пациента с концентрацией нейтрофилов ниже 1,000 клеток/мкл.

\section{Случай 15. Находки в анализе крови суки}
\subsection*{Клиническая история и физикальное обследование}
12-летняя нестерилизованная сука смешанной породы поступила на ежегодный осмотр.

\begin{tabular}{@{}ll@{}}
\toprule
Ректальная температура & 38.1°C \\
ЧСС & 80 уд/мин \\
ЧД & 24 дых/мин \\
ВНК & Менее 2 сек \\
\bottomrule
\end{tabular}

\subsection*{Гемограмма}
\begin{tabular}{lccccc}
\toprule
Тест & Результат & Реф. диапазон & LOW & NORMAL & HIGH \\
\midrule
HCT & 41\% & 37--55 & $\blacksquare$ \\
RBC & 6.36 × 10\textsuperscript{12}/L & 5.50--8.50 & $\blacksquare$ \\
Hb & 135 g/L & 120--180 & $\blacksquare$ \\
MCV & 64 fL & 60--77 & $\blacksquare$ \\
RDW & 12\% & 12.0--15.0 & $\blacksquare$ \\
MCHC & 330 g/L & 310--340 & $\blacksquare$ \\
MCH & 21 pg & 19.5--24.5 & $\blacksquare$ \\
WBC & 17.5 × 10\textsuperscript{9}/L & 6.0--17.0 & $\blacksquare$ \\
NEU & 9.4 × 10\textsuperscript{9}/L & 3.0--12.0 & $\blacksquare$ \\
LYM & 5.41 × 10\textsuperscript{9}/L & 1--4.80 & $\blacksquare$ \\
MONO & 2.30 × 10\textsuperscript{9}/L & 0.20--1.50 & $\blacksquare$ \\
EOS & 0.38 × 10\textsuperscript{9}/L & 0--0.80 & $\blacksquare$ \\
BASO & 0 × 10\textsuperscript{9}/L & 0--0.4 & $\blacksquare$ \\
PLT & 225 × 10\textsuperscript{9}/L & 200--500 & $\blacksquare$ \\
MPV & 7.46 fL & 3.9--11.1 & $\blacksquare$ \\
\% RETIC* & 1 \% & 0.5--1.5 & $\blacksquare$ \\
\% BAND* & 1.2 \% & 0--3.5 & $\blacksquare$ \\
\bottomrule
\multicolumn{6}{l}{*Ручной подсчет}
\end{tabular}

\subsection*{Вопрос: Какие изменения в гемограмме?}

\begin{figure}[htbp]
\centering
\includegraphics[width=0.8\textwidth]{placeholder.jpg}
\caption{Гистограммы крови пациентки, классифицирующие элементы по размеру (горизонтальная ось) и количеству (вертикальная ось); (a) графическое представление лейкоцитов; (b) эритроцитов.}
\end{figure}

\subsection*{Мазки крови}
\begin{figure}[htbp]
\centering
\includegraphics[width=0.6\textwidth]{placeholder.jpg}
\caption{Микроскопическое изображение мазка крови. Окраска Diff-Quick. 400x.}
\end{figure}

\subsection*{Вопрос: Какие клетки видны в мазке крови? Каков предположительный диагноз? Какой другой диагностический тест помог бы подтвердить диагноз?}

\subsection*{Разбор случая}
\subsubsection*{Какие изменения в гемограмме?}
Гистограмма красной крови была нормальной для собак (рис. 1b; симметричная, почти идеальная гауссова кривая; сравните с регенераторной анемией на рис. 1b Случая 5). На гистограмме лейкоцитов видно отсутствие пиков субпопуляций и аномальное удлинение гранулоцитов (рис. 3). Наличие этих изменений на гистограмме лейкоцитов обычно совпадает с появлением крупных атипичных циркулирующих клеток. Сравните это с нормальной гистограммой лейкоцитов собак на рис. 1a из Случая 5.

\begin{figure}[htbp]
\centering
\includegraphics[width=0.8\textwidth]{placeholder.jpg}
\caption{Гистограмма лейкоцитов пациентки. Видно отсутствие характерных пиков каждой субпопуляции и заметное удлинение линии гранулоцитов вправо (область красным; наличие крупных циркулирующих клеток).}
\end{figure}

\subsubsection*{Какие клетки видны в мазке крови?}
Клетки, представленные на этих изображениях мазка (рис. 2), соответствуют атипичным лимфоцитам или лимфобластам. Эти клетки можно распознать по их округлой морфологии, большому размеру, скудному конденсированному хроматину и малому количеству голубоватой цитоплазмы. Любой лимфоцит, размер которого превышает в 2--2.5 раза размер эритроцитов, является атипичным у собак. Эти клетки необходимо дифференцировать от моноцитов (с переплетенным хроматином и почкообразной морфологией ядра) и они могут быть самыми крупными лейкоцитами, присутствующими в мазке крови (сравнение размеров и внешнего вида см. на рис. 4).

\subsubsection*{Каков предположительный диагноз?}
Лабораторные данные этого случая показали предположительные признаки лимфоидной лейкемии с наличием лимфобластов на гематическом уровне, а также лимфоцитоз.

\subsubsection*{Какой другой тест помог бы подтвердить диагноз?}
Учитывая, что большинство лимфом-лейкемических состояний у гериатрических собак соответствует мультицентрической лимфоме, было бы показано исследование костного мозга или пункция лимфатического узла при обнаружении лимфаденопатий. В данном случае пункция подколенного узла подтвердила первоначальное подозрение на лимфопролиферативную клиническую картину (рис. 5).

\begin{figure}[htbp]
\centering
\includegraphics[width=0.8\textwidth]{placeholder.jpg}
\caption{Сравнение размеров между эритроцитами, нейтрофилами и типичными лимфоцитами (a, b) и атипичными лимфоцитами или лимфобластами (c).}
\end{figure}
\begin{figure}[htbp]
\centering
\includegraphics[width=0.6\textwidth]{placeholder.jpg}
\caption{Пункция лимфатического узла, пораженного лимфомой, показала наличие мономорфной популяции средне-крупных лимфоцитов с переплетенным хроматином, многочисленных лимфогландулярных телец и многочисленных атипичных митозов (стрелки).}
\end{figure}

\paragraph{На заметку}
Лимфому и лимфоидную лейкемию можно считать двумя крайностями одного процесса, учитывая, что в большинстве случаев лимфомы присутствуют циркулирующие лимфобласты, и что практически все лимфоидные лейкемии имеют признаки паренхиматозной инфильтрации хотя бы одного органа.

\chapter{Биохимические клинические случаи}
\label{chap:biochem}

\section{Случай 16. Собака с кровавой диареей}
\subsection*{Клиническая история и физикальное обследование}
4-месячный щенок золотистого ретривера (рис. 1) с анамнезом депрессии, анорексии, непрерывной рвоты и диареи с переваренной и свежей кровью в течение недели (рис. 2). Животное еще не было вакцинировано, но было дегельминтизировано месяц назад. Физикальное обследование выявило некоторую спутанность сознания, бледные и заметно сухие слизистые оболочки, а также заметное удлинение времени расправления кожной складки (клиническая оценка: $10\%$ обезвоживания).

\begin{figure}[htbp]
\centering
\includegraphics[width=0.4\textwidth]{placeholder.jpg}
\caption{Пациент.}
\end{figure}
\begin{figure}[htbp]
\centering
\includegraphics[width=0.4\textwidth]{placeholder.jpg}
\caption{Изображение диареи животного.}
\end{figure}

\subsection*{Биохимия при поступлении}
\begin{tabular}{lccc}
\toprule
Тест & Результат & Референсный диапазон \\
\midrule
HCT, \% & 51 & 37--54 \\
Общий белок, g/dL & 6.1 & 5.4--8.2 \\
Альбумин, g/dL & 2.6 & 2.5--4.4 \\
Глобулины, g/dL & 3.5 & 2.3--5.2 \\
Альбумин-глобулиновое соотношение & 0.7 & 0.48--1.91 \\
Общий билирубин, mg/dL & 0.1 & 0.1--0.6 \\
Креатинин, mg/dL & 1.4 & 0.3--1.3 \\
BUN, mg/dL & 98 & 7--25 \\
Глюкоза, mg/dL & 88 & 60--110 \\
Триглицериды, mg/dL & 71 & 20--112 \\
Общий холестерин, mg/dL & 129 & 125--270 \\
Лактат, mmol/L & 2.1 & 0.60--2.90 \\
Амилаза, IU/L & 1,352 & 200--1,200 \\
Щелочная фосфатаза, IU/L & 195 & 20--150 \\
Аланинаминотрансфераза (ALT, GPT), IU/L & 133 & 10--118 \\
Аспартатаминотрансфераза (AST, GOT), IU/L & 88 & 14--45 \\
Гамма-глутамилтрансфераза ($\gamma$GT), IU/L & 5 & 0--7 \\
\bottomrule
\end{tabular}

\subsection*{Вопрос: Какие изменения в биохимии этого животного?}
\subsection*{Вопрос: После регидратации пациента, какие параметры изменятся и как?}

\subsection*{Абдоминальная рентгенограмма}
\begin{figure}[htbp]
\centering
\includegraphics[width=0.6\textwidth]{placeholder.jpg}
\caption{Латеро-латеральная рентгенограмма пациента. Ультразвуковое исследование не показало структурных аномалий на уровне брюшной полости.}
\end{figure}

\subsection*{Вопрос: Какие изменения видны на рентгенограмме?}
\subsection*{Вопрос: Учитывая клинические данные, анамнез и биохимические находки, какой наиболее вероятный возбудитель? Какие тесты следует провести для подтверждения или исключения?}

\subsection*{Разбор случая}
\subsubsection*{Какие изменения в биохимии?}
Биохимия выявила лишь слегка повышенный уровень креатинина, более выраженное повышение мочевины и умеренное повышение амилазы, ALT, AST и AP. Повышенные мочевина и креатинин могли быть вызваны обезвоживанием (преренальная почечная недостаточность), в то время как небольшое повышение трансаминаз обычно наблюдается при гастроэнтерических состояниях (из-за вторичного повреждения печени вследствие всасывания токсинов, среди прочих причин). Хотя остальные аналиты находились в пределах нормы, мы должны учитывать, что у этого животного было клиническое обезвоживание, поэтому уровни HCT и общего белка (TP) были артефактно завышены.

Хотя наличие повышенных HCT и TP часто встречается у обезвоженных животных, очень важно оценивать их гидратационный статус с помощью клинических тестов (табл. 1), поскольку не все обезвоженные пациенты демонстрируют вышеупомянутый паттерн (табл. 2). Пациент в этом клиническом случае был примером обезвоженного животного, лабораторная диагностика которого была бы очень сложной, если бы заранее не были проведены клинические тесты, такие как оценка тургора кожи, учитывая, что и HCT, и TP казались в пределах нормы.

\begin{table}[htbp]
\centering
\caption{Приблизительный расчет процента клинического обезвоживания пациента на основе клинических признаков.}
\begin{tabular}{cl}
\toprule
\% Обезвоживания & Клинические признаки (приблизительные) \\
\midrule
< 5\% & Норма (клиническое подозрение). \\
5\% & Сухие слизистые оболочки и легкое ↓ эластичности кожи. \\
7\% & Сухие слизистые оболочки и умеренное ↓ эластичности кожи. \\
10\% & Сухие слизистые оболочки, выраженное ↓ эластичности кожи, запавшие глаза и легкий ↑ ЧСС. \\
12\% & Очень сухие слизистые оболочки, полное ↓ эластичности кожи, глубоко запавшие глаза, ↑ ЧСС, ↑ ВНК и слабый периферический пульс. \\
15\% & Гиповолемический шок \\
\bottomrule
\multicolumn{2}{l}{*ЧСС = частота сердечных сокращений, ВНК = время наполнения капилляров}
\end{tabular}
\end{table}

\begin{table}[htbp]
\centering
\caption{Паттерны клинических лабораторных результатов у обезвоженных животных.}
\begin{tabular}{@{}c|c|c|c@{}}
\toprule
\multirow{2}{*}{$\uparrow$ HCT} & $\uparrow$ TP & Обезвоживание \\
& N TP & Обезвоживание + гипопротеинемия; норма (например, у грейхаундов); полицитемия \\
& $\downarrow$ TP & Обезвоживание + тяжелая гипопротеинемия \\
\midrule
\multirow{2}{*}{N HCT} & $\uparrow$ TP & Гиперглобулинемия; обезвоживание + анемия \\
& N TP & Норма; сверхострое кровотечение; обезвоживание + анемия + гипопротеинемия \\
& $\downarrow$ TP & Гипопротеинемия; острое кровотечение \\
\midrule
\multirow{2}{*}{$\downarrow$ HCT} & $\uparrow$ TP & Обезвоживание + тяжелая анемия; анемия + гиперглобулинемия \\
& N TP & Анемия (не всегда) \\
& $\downarrow$ TP & Анемия + гипопротеинемия; гипергидратация \\
\bottomrule
\multicolumn{4}{l}{TP = общий белок, HCT = гематокрит, N = нормальный уровень общих белков.}
\end{tabular}
\end{table}

\subsubsection*{После регидратации пациента, какие параметры изменятся и как?}
После регидратации биохимия животного показана ниже. После регидратации уровни HCT и TP снизятся, что приведет к диагнозу умеренной анемии с гипопротеинемией, вызванной гипоальбуминемией у этого животного. Уровни креатинина и мочевины также должны нормализоваться после регидратации (когда восстановится почечный объем). В этом случае мочевина у пациента осталась повышенной (даже когда животное начало мочиться нормально). Эта находка могла быть объяснена наличием желудочно-кишечного кровотечения, которое вызывает повышение уровня мочевины независимо от функции почек.

В таких случаях мы должны измерить уровень альбумина (глобулины обычно рассчитываются путем вычитания альбумина из TP), чтобы оценить, какая из двух фракций ответственна за гипопротеинемию.

\begin{tabular}{lcc}
\toprule
Тест & Результат & Референсный диапазон \\
\midrule
HCT, \% & 33 & 37--54 \\
Общий белок, g/dL & 4.9 & 5.4--8.2 \\
Альбумин, g/dL & 1.7 & 2.5--4.4 \\
Глобулины, g/dL & 3.2 & 2.3--5.2 \\
Альбумин-глобулиновое соотношение & 0.53 & 0.48--1.91 \\
Креатинин, mg/dL & 1.1 & 0.5--1.5 \\
BUN, mg/dL & 68 & 7--25 \\
\bottomrule
\end{tabular}

\paragraph{На заметку}
Основными причинами гипопротеинемии являются кровотечения, потери в третье пространство, гломерулопатии, гастроэнтериты, тяжелая мальабсорбция или недоедание, экссудативный дерматит, иммуносупрессия/дефицит, тяжелая печеночная недостаточность или гипергидратация.

Альбумин-глобулиновое соотношение (A:G) помогает нам изучить, как варьируют эти подтипы белков по отношению друг к другу и теряет ли пациент оба типа белка одинаково или один тип больше другого (табл. 3). В этом клиническом случае мы обнаружили, что наиболее вероятной причиной потери белка был геморрагический гастроэнтерит.

\begin{table}[htbp]
\centering
\caption{Дифференциация причин гипопротеинемии в соответствии с альбумин-глобулиновым соотношением.}
\begin{tabular}{cc}
\toprule
Альбумин-глобулиновое соотношение & Причина \\
\midrule
\uparrow & Гипоглобулинемия (иммунодефицит или иммуносупрессия). \\
Нормальное & Гипоальбуминемия + гипоглобулинемия (гастроэнтерит, дерматит, кровотечение или гипергидратация). \\
\downarrow & Селективная гипоальбуминемия (гломерулопатия, печеночная недостаточность, кишечная мальабсорбция или тяжелое недоедание). \\
\bottomrule
\end{tabular}
\end{table}

\subsubsection*{Какие изменения видны на рентгенограмме?}
Рентгенограмма не обнаружила данных, которые помогли бы продвинуться в диагнозе. Важно не гипердиагностировать наличие нерентгеноконтрастных инородных тел или других обструкций, основываясь только на рентгенографических проекциях. У этого пациента любое такое состояние было исключено с помощью УЗИ.

\subsubsection*{Учитывая клинические данные, анамнез и биохимию, какой наиболее вероятный возбудитель? Какие тесты?}
Диагноз у этого пациента — геморрагический гастроэнтерит, наиболее вероятным возбудителем которого является парвовирус собак. Хотя существуют и другие агенты (коронавирус, чума, \textit{Salmonella}, \textit{Campylobacter}, \textit{Giardia} или \textit{Clostridium} и т.д.), которые могут вызывать сходную клиническую картину и анамнез, учитывая тяжесть парвовируса собак, этот этиологический агент следует исключить в первую очередь.

Учитывая, что у этого пациента уже была неделя клинических симптомов и, следовательно, он находился на продвинутой стадии инфекции (учитывая, что симптомы обычно появляются через 7--8 дней после заражения), наиболее рекомендованным действием был бы ПЦР-тест или, в меньшей степени, серийные серологические исследования.

\begin{table}[htbp]
\centering
\caption{Факторы, которые следует учитывать в отношении различных диагностических тестов на парвовирус собак.}
\begin{tabular}{p{4cm}p{10cm}}
\toprule
Тест & Соображения \\
\midrule
Быстрый фекальный тест (типа SNAP) & • Наиболее полезен на ранних клинических фазах (+ выделение вируса). \\
& • Сбои при очень ранней или хронической инфекции (< 4 или > 14 дней после заражения). \\
ПЦР & Обычно фекальная. Более чувствительна и специфична. Различие между вакцинацией и инфекцией. \\
Серология & Проблема в дифференцировке вакцинация/инфекция (повторные тесты лучше). \\
\bottomrule
\end{tabular}
\end{table}

\section{Случай 17. Гериатрическая сука с хромотой}
\subsection*{Клиническая история и физикальное обследование}
15-летняя сука немецкой овчарки с хроническим анамнезом (два месяца) депрессии, потери аппетита, полиурии/полидипсии, появления подкожных кровотечений и гематохезии, а также перемежающейся хромоты на передние конечности (рис. 1). Животное было правильно вакцинировано и дегельминтизировано, без предшествующей истории заболеваний. Физикальное обследование выявило спутанность сознания и бледные слизистые оболочки с петехиями, хотя гидратация была адекватной.

\begin{tabular}{@{}ll@{}}
\toprule
Ректальная температура & 38.1°C \\
ЧСС & 50 уд/мин \\
ЧД & 28 дых/мин \\
ВНК & Менее 2 сек \\
\bottomrule
\end{tabular}

\begin{figure}[htbp]
\centering
\includegraphics[width=0.5\textwidth]{placeholder.jpg}
\caption{Пациентка во время госпитализации.}
\end{figure}

\subsection*{Биохимия}
\begin{tabular}{lcc}
\toprule
Тест & Результат & Референсный диапазон \\
\midrule
HCT, \% & 31 & 37--54 \\
Общий белок, g/dL & 10.2 & 5.4--8.2 \\
Альбумин, g/dL & 2.8 & 2.5--4.4 \\
Глобулины, g/dL & 7.3 & 2.3--5.2 \\
Альбумин-глобулиновое соотношение & 0.38 & 0.48--1.91 \\
Общий билирубин, mg/dL & 0.1 & 0.1--0.6 \\
Креатинин, mg/dL & 2.6 & 0.3--1.3 \\
BUN, mg/dL & 79 & 7--25 \\
Глюкоза, mg/dL & 102 & 60--110 \\
Триглицериды, mg/dL & 32 & 20--112 \\
Общий холестерин, mg/dL & 159 & 125--270 \\
Амилаза, IU/L & 652 & 200--1,200 \\
Щелочная фосфатаза, IU/L & 395 & 20--150 \\
Аланинаминотрансфераза (ALT, GPT), IU/L & 40 & 10--118 \\
Аспартатаминотрансфераза (AST, GOT), IU/L & 29 & 14--45 \\
Гамма-глутамилтрансфераза ($\gamma$GT), IU/L & 5 & 0--7 \\
Общий кальций, mg/dL & 15.2 & 8.6--11.8 \\
\bottomrule
\end{tabular}

\subsection*{Вопрос: Какие изменения в биохимии?}

\subsection*{Гемостатическое исследование}
\begin{tabular}{lcc}
\toprule
Тест & Результат & Референсный диапазон \\
\midrule
Прогромбиновое время & 34 секунды & 6--12 \\
Активированное частичное тромбопластиновое время & 51 секунда & 12--22 \\
Время кровотечения слизистой оболочки рта & 6.5 минут & < 4 \\
D-димеры & 100 ng/ml & < 250 \\
Тромбоциты & 350 k/$\mu$L & 200--500 \\
\bottomrule
\end{tabular}

\subsection*{Мочевая тест-полоска}
\begin{figure}[htbp]
\centering
\includegraphics[width=0.5\textwidth]{placeholder.jpg}
\caption{Моча пациентки была проверена с помощью мочевой тест-полоски для оценки наличия и концентрации белков в моче.}
\end{figure}

\begin{figure}[htbp]
\centering
\includegraphics[width=0.8\textwidth]{placeholder.jpg}
\caption{Латеро-латеральная рентгенограмма левой передней конечности.}
\end{figure}
\begin{figure}[htbp]
\centering
\includegraphics[width=0.8\textwidth]{placeholder.jpg}
\caption{Протеинограмма сыворотки пациентки.}
\end{figure}

\subsection*{Вопрос: Учитывая полученные результаты, каков предположительный диагноз?}
\subsection*{Вопрос: Какой тест будет проведен для подтверждения диагноза?}

\subsection*{Разбор случая}
\subsubsection*{Какие изменения в биохимии?}
Биохимия этого животного выявила умеренную анемию (умеренно сниженный HCT), причем гораздо более поразительной была гиперпротеинемия, вызванная гиперглобулинемией, нормальными уровнями альбумина и сниженным альбумин-глобулиновым соотношением (A:G).

У животных с гиперглобулинемией мы должны различать моноклональные и поликлональные состояния, выполняя сывороточные протеинограммы. Протеинограмма этой пациентки (рис. 4) показала характерный моноклональный пик с более узким основанием, чем у альбумина, поэтому ее гиперглобулинемия была более совместима с опухолевой клинической картиной (рис. 5).

\begin{figure}[htbp]
\centering
\includegraphics[width=0.8\textwidth]{placeholder.jpg}
\caption{Различные паттерны протеинограмм собак: (a) норма; (b) повышенные $\alpha$ и $\beta$ глобулины, обычно ассоциированные с инфекционными или воспалительными состояниями, с повышением белков острой фазы; (c) поликлональная гаммапатия с более широким основанием, чем у альбумина, которая обычно появляется в результате антигенной стимуляции; и (d) моноклональная гаммапатия (глобулиновый пик с основанием уже, чем у альбумина), вызванная опухолью.}
\end{figure}

\paragraph{На заметку}
Основными причинами гиперпротеинемии являются обезвоживание (повышение альбумина и глобулинов) и гиперглобулинемия, связанная с антигенной стимуляцией (инфекции и аутоиммунные состояния) или новообразованиями (множественная миелома и лимфомы).

Важно отметить, что в некоторых случаях протеинограмма не столь информативна и может привести к диагностическим ошибкам. Поэтому этот тест всегда следует оценивать в сочетании с другими анализами.

\begin{figure}[htbp]
\centering
\includegraphics[width=0.6\textwidth]{placeholder.jpg}
\caption{Деталь рентгенограммы пациентки, показывающая многочисленные остеолитические очаги в костном мозге (стрелки), придающие ему вид «изъеденного молью».}
\end{figure}

У пациентки также была азотемия (повышены креатинин и мочевина) с очень выраженной протеинурией, выявленной при анализе мочи с помощью тест-полоски (рис. 2). В случаях с моноклональными гаммапатиями (гиперглобулинемиями) может появляться поражение почек в результате гипервязкости плазмы (среди прочих причин), и большое количество глобулинов может просачиваться в мочу (протеинурия Бенс-Джонса).

Повышенный уровень щелочной фосфатазы мог быть результатом остеолиза, вызванного опухолью. В костном мозге многих костей этой пациентки (рис. 6) наблюдались множественные рентгеноконтрастные очаги (вид «скученности»), что является находкой, совместимой с резорбцией кости из-за роста определенных типов новообразований (особенно множественной миеломы) на уровне костного мозга.

\paragraph{На заметку}
Существует множество изоформ щелочной фосфатазы (AP), которые могут быть ответственны за повышенный уровень. Наиболее частыми причинами повышения AP являются холестаз или повреждение желчных протоков, повреждение печени, лечение кортикостероидами (индуцированная изоформа) и выраженное ремоделирование костей (тяжелое повреждение костей или рост у гигантских пород). Некоторые более редкие изоформы также происходят из плаценты, почек или кишечника.

Наконец, у пациентки была выраженная гиперкальциемия, которая, вероятно, была опухолевого происхождения (злокачественная гиперкальциемия) или могла быть связана с наличием резорбции костей. Основные причины гиперкальциемии указаны в таблице 1.

\begin{table}[htbp]
\centering
\caption{Основные причины гиперкальциемии у мелких животных.}
\begin{tabular}{p{7cm}p{7cm}}
\toprule
Наиболее частые & Наименее частые \\
\midrule
Злокачественная гиперкальциемия (лимфомы, множественные миеломы или карцинома анального мешка и т.д.). & Гипоадренокортицизм (болезнь Аддисона). \\
Почечная недостаточность (чаще встречается гипокальциемия). & Отравление витамином D (некоторые родентициды). \\
Первичный гиперпаратиреоз (чаще аденомы). & Гранулематозные заболевания или тяжелые поражения костей и т.д. \\
\bottomrule
\end{tabular}
\end{table}

\subsubsection*{Учитывая результаты, каков предположительный диагноз?}
Наличие моноклональной гаммапатии с протеинурией, а также участков остеолиза, где сосредоточено большое количество костного мозга, указывало на множественную миелому. Гемостатическая недостаточность у этой пациентки может быть объяснена вмешательством глобулинов, продуцируемых опухолью, во весь каскад коагуляции.

\subsubsection*{Какой тест будет проведен для подтверждения?}
Диагноз может быть подтвержден проведением исследования костного мозга для проверки наличия неопластических плазматических клеток (рис. 7). Протеинурию Бенс-Джонса также можно проверить с помощью электрофореза белков мочи.

\begin{figure}[htbp]
\centering
\includegraphics[width=0.6\textwidth]{placeholder.jpg}
\caption{Цитология костного мозга, полученная при пункционной биопсии головки плечевой кости, показала, что до 60\% присутствующих клеток составляли плазматические клетки (что совместимо с диагнозом множественной миеломы). Окраска Diff-Quik, 400x.}
\end{figure}

\section{Случай 18. Собака с абдоминальной болью}
\subsection*{Клиническая история и физикальное обследование}
5-летний кастрированный кобель лабрадора-ретривера (рис. 1) поступил в клинику после двух дней депрессии, потери аппетита и продуктивной рвоты каждые несколько часов. Животное обычно питалось домашней пищей и имело небольшой избыточный вес (4/5). Физикальное обследование выявило спутанность сознания и легкое обезвоживание с задержкой кожной складки и сухостью слизистых оболочек. Животное реагировало на пальпацию обоих подреберьев болью и демонстрировало защитное напряжение живота.

\begin{tabular}{@{}ll@{}}
\toprule
Ректальная температура & 38.0°C \\
ЧСС & 150 уд/мин \\
ЧД & 20 дых/мин \\
ВНК & Менее 2 сек \\
\bottomrule
\end{tabular}

\begin{figure}[htbp]
\centering
\includegraphics[width=0.5\textwidth]{placeholder.jpg}
\caption{Пациент на приеме с умеренной депрессией и легким обезвоживанием.}
\end{figure}

\subsection*{Гемограмма}
\begin{tabular}{lccccc}
\toprule
Тест & Результат & Реф. диапазон & LOW & NORMAL & HIGH \\
\midrule
HCT & 33\% & 37--55 & $\blacksquare$ \\
RBC & 4.36 × 10\textsuperscript{12}/L & 5.50--8.50 & $\blacksquare$ \\
Hb & 105 g/L & 120--180 & $\blacksquare$ \\
MCV & 65 fL & 60--77 & $\blacksquare$ \\
RDW & 12\% & 12.0--15.0 & $\blacksquare$ \\
MCHC & 320 g/L & 310--340 & $\blacksquare$ \\
MCH & 24 pg & 19.5--24.5 & $\blacksquare$ \\
WBC & 19.5 × 10\textsuperscript{9}/L & 6.0--17.0 & $\blacksquare$ \\
NEU & 13.9 × 10\textsuperscript{9}/L & 3.0--12.0 & $\blacksquare$ \\
LYM & 4.05 × 10\textsuperscript{9}/L & 1--4.80 & $\blacksquare$ \\
MONO & 0.97 × 10\textsuperscript{9}/L & 0.20--1.50 & $\blacksquare$ \\
EOS & 0.58 × 10\textsuperscript{9}/L & 0--0.80 & $\blacksquare$ \\
BASO & 0 × 10\textsuperscript{9}/L & 0--0.4 & $\blacksquare$ \\
PLT & 84 × 10\textsuperscript{9}/L & 200--500 & $\blacksquare$ \\
MPV & 7.46 fL & 3.90--11.1 & $\blacksquare$ \\
\% RETIC* & 1\% & 0.5--1.5 & $\blacksquare$ \\
\% BAND* & 4\% & 0--3.5 & $\blacksquare$ \\
\bottomrule
\multicolumn{6}{l}{*Ручной подсчет}
\end{tabular}

\begin{tabular}{lcc}
\toprule
Тест & Результат & Референсный диапазон \\
\midrule
Прогромбиновое время & 7 секунд & 6--12 \\
Активированное частичное тромбопластиновое время & 18 секунд & 12--22 \\
\bottomrule
\end{tabular}

\subsection*{Вопрос: Какие изменения в гемограмме?}

\subsection*{Биохимия}
\begin{tabular}{lcc}
\toprule
Тест & Результат & Референсный диапазон \\
\midrule
Общий белок, g/dL & 5.5 & 5.4--8.2 \\
Альбумин, g/dL & 3 & 2.5--4.4 \\
Фибриноген, mg/dL & 410 & 100--245 \\
Общий билирубин, mg/dL & 1.1 & 0.1--0.6 \\
Креатинин, mg/dL & 1.4 & 0.3--1.3 \\
BUN, mg/dL & 29 & 7--25 \\
Глюкоза, mg/dL & 210 & 60--110 \\
Триглицериды, mg/dL & 399 & 20--112 \\
Общий холестерин, mg/dL & 498 & 125--270 \\
Лактат, mmol/L & 2.1 & 0.60--2.90 \\
Амилаза, IU/L & 1,252 & 200--1,200 \\
Креатинфосфокиназа (CK, CPK), IU/L & 55 & 20--200 \\
Щелочная фосфатаза, IU/L & 215 & 20--120 \\
Аланинаминотрансфераза (ALT, GPT), IU/L & 143 & 10--118 \\
Аспартатаминотрансфераза (AST, GOT), IU/L & 98 & 14--45 \\
Гамма-глутамилтрансфераза ($\gamma$GT), IU/L & 6 & 0--7 \\
Лактатдегидрогеназа (LDH), IU/L & 294 & 25--220 \\
Общий кальций, mg/dL & 6.2 & 8.6--11.8 \\
Фосфор, mg/dL & 5.1 & 2.9--6.6 \\
Калий, mmol/L & 3.7 & 3.7--5.8 \\
Натрий, mmol/L & 158 & 138--160 \\
Хлор, mmol/L & 114 & 106--120 \\
\bottomrule
\end{tabular}

\subsection*{Вопрос: Какие изменения в биохимии?}
\subsection*{Вопрос: Каков предположительный диагноз?}
\subsection*{Вопрос: Какой лабораторный тест подтвердит предположительный диагноз?}
\subsection*{Вопрос: Какой другой диагностический тест помог бы подтвердить его?}

\subsection*{Разбор случая}
\subsubsection*{Какие изменения в гемограмме?}
У пациента нормохромная нормоцитарная анемия со снижением гематокрита, гемоглобина и эритроцитов, но MCV и MCHC в пределах нормы. Этот тип анемии часто встречается у пациентов с хроническими воспалительными или инфекционными состояниями. Анализ лейкоцитов показал лейкоцитоз с нейтрофилией и регенераторным левым сдвигом (увеличение количества палочкоядерных нейтрофилов). Что касается тромбоцитов, обнаружена тромбоцитопения. Следует проверить наличие тромбоцитарных агрегатов в мазке крови (рис. 2) или, лучше, выполнить ручной подсчет тромбоцитов, чтобы исключить возможность псевдотромбоцитопении у животного.

\paragraph{На заметку}
Одной из наиболее частых причин низких значений тромбоцитов является псевдотромбоцитопения, вызванная образованием тромбоцитарных агрегатов при заборе крови. Эти агрегаты могут включать даже лейкоциты, что снижает их количество (лейкоцитарно-тромбоцитарные агрегаты). Их образование можно уменьшить, собирая образец в пробирки с цитратом натрия.

\begin{figure}[htbp]
\centering
\includegraphics[width=0.4\textwidth]{placeholder.jpg}
\caption{Лейкоцитарно-тромбоцитарный агрегат при увеличении 40x.}
\end{figure}

\subsubsection*{Какие изменения в биохимии?}
У животного были следующие изменения:
\begin{itemize}
\item Выраженная гипокальциемия.
\item Повышенный уровень фибриногена: белок острой фазы, являющийся неспецифическим индикатором заболевания.
\item Азотемия (повышение мочевины и креатинина), указывающая на развивающееся повреждение почек, а также повышенные маркеры печеночного повреждения, такие как общий билирубин и ферменты гепатоцеллюлярного повреждения (ALT, AST и AP).
\item Гиперамилаземия у этого животного могла указывать на раннее повреждение поджелудочной железы, хотя это не является однозначным.
\item Гипергликемия могла подтвердить вышеупомянутое повреждение поджелудочной железы, хотя другие причинные состояния должны быть исключены.
\item Повышение триглицеридов и холестерина в сыворотке (гиперлипидемия) также могло повлиять на лабораторные измерения.
\item Повышение LDH могло быть вызвано повреждением почек (учитывая, что CK не была повышена, мышечное повреждение было маловероятным), но это не может быть подтверждено без определения изоформ этого фермента.
\end{itemize}

\paragraph{На заметку}
Чтобы избежать ложных диагнозов гипокальциемии из-за эффекта гипоальбуминемии, значение кальция всегда следует корректировать следующим образом: \\[2mm]
\textbf{СКОРРЕКТИРОВАННЫЙ ОБЩИЙ Ca (mg/dL) = 3.5 - АЛЬБУМИН (g/dL) + ИЗМЕРЕННЫЙ ОБЩИЙ Ca (mg/dL).}

Основные причины ГИПОКАЛЬЦИЕМИИ у мелких животных:

\begin{tabular}{p{5cm}p{9cm}}
\toprule
ПРИЧИНА & ПОДТВЕРЖДЕНИЕ ДИАГНОЗА \\
\midrule
Гипоальбуминемия & Используйте скорректированную формулу кальция. \\
Терминальная стадия почечной недостаточности & + Азотемия и гиперфосфатемия (обычно). \\
Острый панкреатит & Клиническая картина, амилаза, повреждение печени или гипергликемия и т.д. \\
Вторичный гиперпаратиреоз & Признаки почечной недостаточности. \\
Гипопаратиреоз & + Гиперфосфатемия БЕЗ азотемии. \\
Диабетический кетоацидоз & + Гипергликемия и метаболический ацидоз. \\
Отравление этиленгликолем & Клиническая картина, метаболический ацидоз или кристаллы оксалата кальция. \\
\bottomrule
\end{tabular}

Гиперлипидемия (рис. 3) может быть физиологической (постпрандиальной) или патологической (обычно эндокринной). Чтобы дифференцировать эти причины, пациента можно подвергнуть голоданию, провести тест с охлаждением или ультрацентрифугированием образца.

Наличие гиперлипидемии в образце может вызвать следующие лабораторные артефакты:
\begin{itemize}
\item Очень выраженный гемолиз.
\item Трудности при считывании уровня общего белка (TP) рефрактометрией.
\item Гипергликемия.
\item Гиперкальциемия и гиперфосфатемия.
\item Повышенный билирубин.
\item Снижение альбумина TP (по данным колориметрии), а также снижение уровня натрия и калия.
\end{itemize}

\begin{figure}[htbp]
\centering
\includegraphics[width=0.5\textwidth]{placeholder.jpg}
\caption{Липемичная плазма беловатого цвета у собаки-пациента.}
\end{figure}

\subsubsection*{Каков предположительный диагноз?}
Симптомы животного, наряду с изменениями, наблюдаемыми в его гемограмме и биохимии, должны заставить нас заподозрить клиническую картину острого панкреатита.

\subsubsection*{Какой лабораторный тест подтвердит предположительный диагноз?}
Этому животному показан тест на иммунореактивность панкреатической липазы (PLI).

\subsubsection*{Какой другой диагностический тест помог бы подтвердить диагноз?}
У животных с острым панкреатитом абдоминальные рентгенограммы и УЗИ могут помочь в диагностике процесса (рис. 4).

\begin{table}[htbp]
\centering
\caption{Диагностическое сравнение различных лабораторных тестов для оценки поджелудочной железы.}
\begin{tabular}{p{4cm}p{5cm}p{5cm}}
\toprule
ТЕСТ & ХАРАКТЕРИСТИКИ & НЕДОСТАТКИ \\
\midrule
Амилаза & • Повышение (7--10×) при клинической картине острого панкреатита у собак. & • Малоценна у кошек. \\
& • Пик через 12--48 часов. & • Также повышается (< 4×) при сахарном диабете, почечной недостаточности, гастроэнтерических симптомах, гепатобилиарных заболеваниях и новообразованиях и т.д. \\
& • Остается повышенной в течение 8--14 дней. & \\
Липаза & • Повышение (3--4×) при панкреатите. & • Малоценна у кошек. \\
& • Пик через 24--48 часов. & • Также повышается (< 3×) при перитоните, почечной недостаточности, гастроэнтеритах и гепатобилиарных заболеваниях. \\
& • Остается повышенной в течение 8--14 дней. & • Терапия кортикостероидами также повышает уровень липазы (5×). \\
& & • Панкреатические или гепатобилиарные новообразования могут повышать липазу до 90×. \\
Трипсиноподобная иммунореактивность (TLI) & • Требует 12-часового голодания. & • Высокое количество ложноотрицательных результатов (30--35\%). \\
& • Уровни прямо пропорциональны массе поджелудочной железы. & • Проблема для диагностики хронических или подхронических клинических картин. \\
& & • TLI также повышается в результате почечной недостаточности. \\
Иммунореактивность панкреатической липазы (PLI) & • Наиболее специфичный тест для диагностики панкреатита у собак и кошек. & • Этот тест имеет умеренную специфичность и чувствительность у кошек. \\
& • Измерение радиоиммунным анализом в специализированных лабораториях или использование коррелированных тестов SNAP®. & \\
& • PLI не изменяется при терапии кортикостероидами. & \\
\bottomrule
\end{tabular}
\end{table}

\begin{figure}[htbp]
\centering
\includegraphics[width=0.6\textwidth]{placeholder.jpg}
\caption{Типичное ультразвуковое изображение на уровне поджелудочной железы-двенадцатиперстной кишки у пациентов с острым панкреатитом. Наблюдалась эхогенная неоднородность в области поджелудочной железы и «гофрированная» стенка на уровне двенадцатиперстной кишки.}
\end{figure}

\section{Случай 19. Очень угнетенная и спутанная сука}
\subsection*{Клиническая история и физикальное обследование}
Сука йоркширского терьера в возрасте 2,5 лет (рис. 1) поступила с несколькими эпизодами судорог (генерализованные припадки) в последние месяцы. Она также стала привередливой в еде с прогрессирующей потерей аппетита и умеренной полиурией и полидипсией. Животное было правильно вакцинировано и дегельминтизировано. Владельцев поразили ее выраженная апатия, отсутствие внимания и способности к обучению, а также плохой рост по сравнению с остальными щенками из того же помета. При физикальном обследовании отмечены выраженная депрессия, сонливость и летаргия. Слизистые оболочки казались нормального цвета, и состояние гидратации пациентки было адекватным.

\begin{figure}[htbp]
\centering
\includegraphics[width=0.5\textwidth]{placeholder.jpg}
\caption{Пациентка при поступлении.}
\end{figure}

\subsection*{Биохимия}
\begin{tabular}{lcc}
\toprule
Тест & Результат & Референсный диапазон \\
\midrule
HCT, \% & 30 & 37--54 \\
Общий белок, g/dL & 5.3 & 5.4--8.2 \\
Альбумин, g/dL & 1.6 & 2.5--4.4 \\
Глобулины, g/dL & 3.7 & 2.3--5.2 \\
Альбумин-глобулиновое соотношение & 0.43 & 0.48--1.91 \\
Общий билирубин, mg/dL & 0.2 & 0.1--0.6 \\
Креатинин, mg/dL & 0.6 & 0.3--1.3 \\
BUN, mg/dL & 4 & 7--25 \\
Глюкоза, mg/dL & 32 & 60--110 \\
Триглицериды, mg/dL & 45 & 20--112 \\
Общий холестерин, mg/dL & 99 & 125--270 \\
Амилаза, IU/L & 590 & 200--1,200 \\
Щелочная фосфатаза, IU/L & 45 & 20--150 \\
Аланинаминотрансфераза (ALT, GPT), IU/L & 20 & 10--118 \\
Аспартатаминотрансфераза (AST, GOT), IU/L & 19 & 14--45 \\
Гамма-глутамилтрансфераза ($\gamma$GT), IU/L & 6 & 0--7 \\
\bottomrule
\end{tabular}

\subsection*{Вопрос: Какие изменения в биохимии?}

\subsection*{Мазок крови}
\begin{figure}[htbp]
\centering
\includegraphics[width=0.6\textwidth]{placeholder.jpg}
\caption{Мазок крови пациентки. Окраска Diff-Quik, 400x.}
\end{figure}

\subsection*{Вопрос: Какие изменения в мазке крови?}

\subsection*{Осадок мочи (после катетеризации)}
\begin{figure}[htbp]
\centering
\includegraphics[width=0.6\textwidth]{placeholder.jpg}
\caption{Неокрашенный осадок мочи после центрифугирования. Получен после катетеризации через уретру. 100x.}
\end{figure}

\subsection*{Вопрос: Какие изменения видны в осадке мочи?}
\subsection*{Вопрос: Учитывая полученные результаты, каков предположительный диагноз?}
\subsection*{Вопрос: Какой лабораторный тест будет выполнен для подтверждения?}

\subsection*{Разбор случая}
\textbf{Портосистемный шунт}

\subsubsection*{Какие изменения в биохимии?}
Были обнаружены умеренная анемия и гипопротеинемия за счет гипоальбуминемии. У этой пациентки альбумин-глобулиновое соотношение (A:G) использовалось для проверки наличия селективной потери/снижения альбумина, что совместимо с картинами гломерулопатии, печеночной недостаточности, кишечной мальабсорбции или тяжелого недоедания (табл. 3 из случая 16).

Кроме того, у животного было снижение концентрации холестерина и мочевины, причем последнее особенно является признаком печеночной недостаточности.

У пациентки также была значительная гипогликемия, которая, вероятно, была причиной ее спутанности сознания. Основные причины гипогликемии перечислены в таблице 1.

Наконец, хотя все параметры указывали на недостаточность или отказ функции печени (табл. 2), не было выраженного повреждения гепатоцитов или холестаза, поскольку уровни $\gamma$GT и трансаминаз не были повышены.

\paragraph{На заметку}
Мочевина синтезируется в печени из аммиака, образующегося в толстой кишке при бактериальном переваривании аминокислот, и выводится в основном почками. Основными причинами повышения мочевины являются снижение скорости почечной элиминации (из-за повреждения почек или уменьшения почечного объема) или повышение уровня белка в толстой кишке (обычно из-за кишечного кровотечения). Снижение уровня мочевины обычно вызвано прогрессирующей печеночной недостаточностью.

\begin{table}[htbp]
\centering
\caption{Основные причины гипогликемии у мелких животных.}
\begin{tabular}{p{6cm}p{8cm}}
\toprule
Артефактная & Задержка отделения сыворотки/плазмы от клеток в пробирке. \\
Сбой цикла глюкозы & Печеночная недостаточность, тяжелое недоедание или неонатальные случаи. \\
Избыток инсулина & Инсулиномы или ятрогенные причины. \\
Влияние других гормонов & Гипоадренокортицизм (болезнь Аддисона) или гипофизарный нанизм. \\
Другие причины & Сепсис, паранеопластический синдром или экстремальный лейкоцитоз и т.д. \\
\bottomrule
\end{tabular}
\end{table}

\begin{table}[htbp]
\centering
\caption{Примеры различных лабораторных диагностических тестов на уровне печени.}
\begin{tabular}{p{3cm}p{4cm}p{7cm}}
\toprule
Тип маркера & Примеры & Примечания \\
\midrule
Повреждение гепатоцитов (утечка) & ALT, AST или AP и т.д. & Повышение из-за разрушения или повреждения гепатоцитов (будь то обратимое или нет). Эти тесты не предоставляют информации об уровне функции печени. \\
Холестаз & AP и $\gamma$GT & Повреждение или обструкция желчных протоков. \\
Функция печени & Билирубин, желчные кислоты, аммиак, альбумин, глюкоза или мочевина и т.д. & Они указывают на функцию печени, но не на повреждение гепатоцитов. \\
\bottomrule
\multicolumn{3}{l}{ALT = аланинаминотрансфераза, AST = аспартатаминотрансфераза, AP = щелочная фосфатаза, $\gamma$GT = гамма-глутамилтрансфераза.}
\end{tabular}
\end{table}

\subsubsection*{Какие изменения в мазке крови?}
Наблюдался умеренный дефицит железа, при этом эритроциты имели очень заметные и расширенные центральные ореолы. Также были обнаружены редкие микроциты с гораздо меньшим размером, чем остальные эритроциты (и пропорционально присутствующим нейтрофилам) (рис. 4), а также акантоциты с нерегулярными и асимметричными выступами.

\subsubsection*{Какие изменения видны в осадке мочи?}
В осадке на дне образца наблюдалось много эритроцитов (гематурия, в данном случае скорее всего связанная с катетеризацией мочевыводящих путей), а также множество кристаллов биурата аммония. Эти кристаллы были в основном от коричневатого до оранжевого цвета и напоминали по форме прорастающий картофель (овальные или округлые с нерегулярными выступами). Этот тип кристаллов очень типичен для клинической картины портосистемного шунта, а также для некоторых непеченочных заболеваний (в результате дефицита ферментов в цикле аммиака) у далматинцев и бульдогов. Наличие этих кристаллов не всегда сопровождается образованием камней в мочевом пузыре.

\begin{figure}[htbp]
\centering
\includegraphics[width=0.6\textwidth]{placeholder.jpg}
\caption{Эритроциты с выраженными расширенными ореолами (умеренная гипохромия), а также редкие микроциты (черные стрелки) и акантоциты (синие стрелки).}
\end{figure}

\subsubsection*{Учитывая результаты, каков предположительный диагноз?}
Анамнез дефицита роста и неврологические отклонения (из-за гипогликемии и печеночной энцефалопатии, вызванной аммонием), а также все биохимические и гематологические результаты и наличие кристаллов биурата аммония в моче соответствовали диагнозу портосистемного шунта. Этот диагноз может быть подтвержден с помощью УЗИ (иногда трудно достижимо), контрастной рентгенографии (портограмма) или более продвинутых методов, таких как сцинтиграфия, компьютерная томография или магнитно-резонансная томография, а также с помощью лабораторных тестов.

\subsubsection*{Какой лабораторный тест будет выполнен для подтверждения?}
Как определение желчных кислот, так и измерение аммиака могут помочь подтвердить диагноз. Учитывая летучесть аммиака и связанные с этим трудности в получении достоверного результата, а также сходную чувствительность двух методов, гораздо более целесообразно измерить уровни желчных кислот. Кроме того, для получения гораздо более точного результата показана оценка желчных кислот как натощак, так и через два часа после еды.

\section{Случай 20. Кот с полидипсией и полиурией}
\subsection*{Клиническая история и физикальное обследование}
3-летний некастрированный кот (рис. 1) с анамнезом потери аппетита и веса, полидипсии и полиурии в течение двух месяцев. За три дня до обращения животное перестало мочиться и стало очень вялым. Пациент не был вакцинирован или дегельминтизирован. При физикальном обследовании отмечены сухие и слегка бледные слизистые оболочки, а также умеренное снижение тургора кожи. Образец мочи был собран с помощью цистоцентеза, и была определена плотность мочи (рис. 2). После проведения общего биохимического анализа пациент был госпитализирован и находился под наблюдением (табл. 1).

\begin{figure}[htbp]
\centering
\includegraphics[width=0.5\textwidth]{placeholder.jpg}
\caption{Пациент при поступлении.}
\end{figure}

\begin{tabular}{@{}ll@{}}
\toprule
Ректальная температура & 37.5°C \\
ЧСС & 90 уд/мин \\
ЧД & 52 дых/мин \\
ВНК & Менее 2 сек \\
\bottomrule
\end{tabular}

\begin{figure}[htbp]
\centering
\includegraphics[width=0.4\textwidth]{placeholder.jpg}
\caption{Плотность мочи образца, полученного цистоцентезом и измеренного рефрактометрией.}
\end{figure}

\subsection*{Биохимия}
\begin{tabular}{lcc}
\toprule
Тест & Результат & Референсный диапазон \\
\midrule
HCT, \% & 22 & 24--45 \\
Общий белок, g/dL & 10.2 & 5.4--8.2 \\
Альбумин, g/dL & 4.4 & 2.2--4.4 \\
Глобулины, g/dL & 5.8 & 1.5--5.7 \\
Альбумин-глобулиновое соотношение & 0.75 & 0.38--2.93 \\
Общий билирубин, mg/dL & 0.1 & 0.1--0.6 \\
Глюкоза, mg/dL & 102 & 70--150 \\
Триглицериды, mg/dL & 25 & 10--100 \\
Общий холестерин, mg/dL & 133 & 90--205 \\
Амилаза, IU/L & 1,487 & 300--1,500 \\
Щелочная фосфатаза, IU/L & 58 & 10--90 \\
Аланинаминотрансфераза (ALT, GPT), IU/L & 98 & 20--100 \\
Аспартатаминотрансфераза (AST, GOT), IU/L & 34 & 12--43 \\
Гамма-глутамилтрансфераза ($\gamma$GT), IU/L & 1 & 0--2 \\
\bottomrule
\end{tabular}

\subsection*{Эволюция во время госпитализации}
\begin{table}[htbp]
\centering
\caption{Наблюдение за пациентом в течение недели госпитализации.}
\begin{tabular}{lccccc}
\toprule
& Реф. & День 0 & День 1 & День 2 & День 7 \\
\midrule
Дегидратация (\%) & 0 & 7 & 5 & 0 & 0 \\
Объем мочи за 24 ч (mL) & 200 & 51 & 25 & 250 & 550 \\
BUN (mg/dL) & 10--30 & 195 & 137 & 105 & 72 \\
Креатинин (mg/dL) & 0.3--1.6 & 5.2 & 4.1 & 3.9 & 3.8 \\
BUN--креатининовое соотношение & < 20 & 37.5 & 33.4 & 26.9 & 18.9 \\
Плотность мочи & > 1,020 & 1,002 & 1,005 & 1,006 & 1,002 \\
Примечания & & • Начата в/в инфузия. & • Нормальный анализ мочи. & • В/в инфузия прекращена. & • Нормальное систолическое АД. \\
& & & • Нет протеинурии. & • HCT = 17 & \\
& & & & • TP = 7.6 & \\
\bottomrule
\multicolumn{6}{l}{В/в: внутривенно; TP: общий белок.}
\end{tabular}
\end{table}

\subsection*{Вопрос: Какие изменения в биохимии?}

\subsection*{УЗИ почек}
\begin{figure}[htbp]
\centering
\includegraphics[width=0.6\textwidth]{placeholder.jpg}
\caption{УЗИ почек пациента.}
\end{figure}

\subsection*{Вопрос: Учитывая результаты, каков предположительный диагноз?}
\subsection*{Вопрос: На какой стадии заболевания находится пациент?}

\subsection*{Разбор случая}
\textbf{Острая почечная недостаточность на фоне хронической почечной болезни}

\subsubsection*{Какие изменения в биохимии?}
В начальной биохимии наблюдалась анемия (низкий HCT, даже с клиническим обезвоживанием), а также повышенные TP, альбумин и глобулины (из-за обезвоживания). Также было обнаружено повышение амилазы, которое можно объяснить почечной недостаточностью выведения.

Анемия была подтверждена при последующем наблюдении за животным (табл. 1). На 0-й день наблюдения были обнаружены типичные признаки острой преренальной почечной недостаточности (из-за обезвоживания) с наличием олигоанурии, азотемии (повышение мочевины и креатинина) и повышенного соотношения мочевина-креатинин (табл. 2).

\paragraph{На заметку}
Соотношение мочевина-креатинин следует использовать только для различения преренальной и ренальной недостаточности, всегда помня, что существует перекрытие между этими двумя состояниями и что оба эти параметра зависят от других факторов (мышечная масса, функция печени, наличие желудочно-кишечного кровотечения и т.д.).

Однако, несмотря на обезвоживание, моча была гипостенуричной, в то время как у обезвоженных животных обычно наблюдается повышенная плотность мочи (табл. 3). Эта находка должна вызвать подозрение на лежащее в основе повреждение почек.

\paragraph{На заметку}
У некоторых кошек с преренальной почечной недостаточностью концентрация мочи не происходит. В этих случаях самый простой способ отличить вышеуказанное от чисто почечной недостаточности — наблюдать за ответом после восстановления объема крови (при преренальной недостаточности азотемия должна разрешиться).

\begin{table}[htbp]
\centering
\caption{Основные характеристики различных подтипов почечной недостаточности у мелких животных.}
\begin{tabular}{lccccc}
\toprule
Тип ПН & Азотемия & Соотношение мочевина-креатинин & Плотность мочи & Другие \\
\midrule
Преренальная & ДА & > 30 & \uparrow\uparrow\uparrow & Олигоанурия \\
Ренальная & ДА & < 30 & Гипо- или изостенурия & Полиурия или олигоанурия; вариабельно \\
Постренальная & ДА & Вариабельно & Вариабельно & Дизурия, странгурия или активный осадок и т.д. \\
\bottomrule
\end{tabular}
\end{table}

\begin{table}[htbp]
\centering
\caption{Классификация плотности мочи у мелких животных.}
\begin{tabular}{lcc}
\toprule
Плотность мочи & Собака & Кошка \\
\midrule
Норма & 1,020--1,045 & 1,020--1,050 \\
При 5\% обезвоживании & 1,040--1,075 & 1,045--1,088 \\
Изостенурия & \multicolumn{2}{c}{1,007--1,013} \\
Гипостенурия & \multicolumn{2}{c}{< 1,007} \\
\bottomrule
\end{tabular}
\end{table}

Во время госпитализации обезвоживание разрешилось, уровни мочевины снизились (поскольку продукция мочи восстановилась), а уровни креатинина снижались гораздо медленнее. К концу госпитализации у пациента все еще была азотемия, и он не мог концентрировать мочу (типичные признаки чисто почечной недостаточности).

\subsubsection*{Учитывая результаты, каков предположительный диагноз?}
Клинические данные, первоначальные лабораторные тесты и частичный ответ на регидратацию, представленный этим котом, были совместимы с острой почечной недостаточностью из-за обезвоживания. Кроме того, учитывая, что почечная недостаточность сохранялась после разрешения проблемы с объемом крови, а также ультразвуковая картина почек, показывающая выраженную гиперэхогенность с потерей кортико-медуллярной дифференцировки и уменьшенным размером (рис. 3), у этого пациента был фон хронической почечной недостаточности.

\subsubsection*{На какой стадии заболевания находится пациент?}
Согласно классификации Международного общества по изучению почек (IRIS), на момент поступления (день 0) у этого пациента была IV степень острой почечной недостаточности (табл. 4). После разрешения обезвоживания и острой почечной недостаточности (день 7) хроническая почечная недостаточность пациента была классифицирована как стадия 3 (табл. 5).

\paragraph{На заметку}
Основными причинами гипо/изо/стенурии являются почечная недостаточность (сопутствующая азотемия), несахарный диабет (центральный или нефрогенный), форсированный диурез (инфузионная терапия), психогенная полидипсия, промывание спинного мозга (печеночная недостаточность или болезнь Аддисона) или наличие выраженной глюкозурии (сахарный диабет).

\begin{table}[htbp]
\centering
\caption{Классификация острой почечной недостаточности (ОПН) на основе сывороточного уровня креатинина согласно классификации Международного общества по изучению почек (IRIS).}
\begin{tabular}{cc}
\toprule
Степень ОПН & Уровень креатинина (mg/dL) \\
\midrule
I & < 1.6 \\
II & 1.7--2.5 \\
III & 2.6--5.0 \\
IV & 5.0--10 \\
V & > 10 \\
\bottomrule
\multicolumn{2}{l}{Существуют подстепени ОПН в зависимости от наличия олигоанурии или необходимости трансплантации почки. Пациенты без азотемии классифицируются на основе эволюции уровней мочевины и креатинина и/или начала олигоанурии.}
\end{tabular}
\end{table}

\begin{table}[htbp]
\centering
\caption{Классификация хронической почечной недостаточности (ХПН) на основе сывороточного уровня креатинина согласно классификации Международного общества по изучению почек (IRIS).}
\begin{tabular}{ccc}
\toprule
Стадия ХПН & \multicolumn{2}{c}{Уровень креатинина (mg/dL)} \\
& Собака & Кошка \\
\midrule
1 & < 1.4 & < 1.6 \\
2 & 1.4--2.0 & 1.6--2.8 \\
3 & 2.1--5.0 & 2.9--5.0 \\
4 & \multicolumn{2}{c}{> 5} \\
\bottomrule
\multicolumn{3}{l}{Существуют подстадии в зависимости от наличия протеинурии (определяемой на основе UPC) и/или артериальной гипертензии.}
\multicolumn{3}{l}{UPC: соотношение белок-креатинин в моче.}
\end{tabular}
\end{table}

\section{Случай 21. Кот с дизурией}
\subsection*{Клиническая история и физикальное обследование}
6-летний кастрированный кот с недавним анамнезом потери аппетита и веса, а также изменениями привычек мочеиспускания (рис. 1). Животное в течение нескольких месяцев мочилось очень часто и малыми порциями, иногда выделяя красноватую мочу, временами вокализируя при попытке мочеиспускания и избегая лотка в последние недели. Он также проявлял умеренную полидипсию и не ел в течение двух дней. Физикальное обследование выявило слизистые оболочки нормального цвета и адекватный статус гидратации. Пальпировался большой мочевой пузырь, и животное жаловалось при манипуляциях. Во время забора пробы животное сопротивлялось и заметно напряглось, что привело к выраженному гемолизу в полученном образце плазмы (рис. 2).

\begin{figure}[htbp]
\centering
\includegraphics[width=0.5\textwidth]{placeholder.jpg}
\caption{Пациент при поступлении.}
\end{figure}

\begin{tabular}{@{}ll@{}}
\toprule
Ректальная температура & 37.6°C \\
ЧСС & 90 уд/мин \\
ЧД & 56 дых/мин \\
ВНК & Менее 2 сек \\
\bottomrule
\end{tabular}

\begin{figure}[htbp]
\centering
\includegraphics[width=0.5\textwidth]{placeholder.jpg}
\caption{Гепариновая пробирка пациента (слева) по сравнению с образцом другого кота (справа) после центрифугирования.}
\end{figure}

\subsection*{Биохимия}
\begin{tabular}{lcc}
\toprule
Тест & Результат & Референсный диапазон \\
\midrule
HCT, \% & 31 & 24--45 \\
Общий белок, g/dL & 9.2 & 5.4--8.2 \\
Альбумин, g/dL & 2.6 & 2.2--4.4 \\
Глобулины, g/dL & 6.6 & 1.5--5.7 \\
Альбумин-глобулиновое соотношение & 0.39 & 0.38--2.93 \\
Общий билирубин, mg/dL & 0.0 & 0.1--0.6 \\
Креатинин, mg/dL & 0.9 & 0.3--1.6 \\
BUN, mg/dL & 27 & 10--30 \\
Глюкоза, mg/dL & 199 & 70--150 \\
Триглицериды, mg/dL & 295 & 10--100 \\
Общий холестерин, mg/dL & 101 & 90--205 \\
Амилаза, IU/L & 990 & 300--1,500 \\
Щелочная фосфатаза, IU/L & 0 & 10--90 \\
Аланинаминотрансфераза (ALT, GPT), IU/L & 65 & 20--100 \\
Аспартатаминотрансфераза (AST, GOT), IU/L & 455 & 12--43 \\
Гамма-глутамилтрансфераза ($\gamma$GT), IU/L & 0 & 0--2 \\
\bottomrule
\end{tabular}

\subsection*{Вопрос: Какие изменения в биохимии?}
\subsection*{Вопрос: Какие биохимические параметры могут быть недостоверны в этом образце?}

\subsection*{Анализ мочи (цистоцентез)}
\begin{figure}[htbp]
\centering
\includegraphics[width=0.5\textwidth]{placeholder.jpg}
\caption{Моча, содержащаяся в мочевом пузыре.}
\end{figure}

\begin{tabular}{l}
\caption{Результаты анализа мочи.} \\
\toprule
Цвет и мутность \& Оранжевая, прозрачная \\
Плотность мочи (рефрактометрия) \& 1,050 \\
Глюкоза (тест-полоска) \& $\approx$ 250 mg/dL \\
Белки (тест-полоска) \& +++ ($\approx$ 300 mg/dL) \\
pH (тест-полоска) \& 8 \\
Бактерии (посев) \& \textit{Escherichia coli} \\
\bottomrule
\end{tabular}

\subsection*{Последующее наблюдение}
После адекватного лечения дизурии животное вернулось на контрольный осмотр и больше не проявляло признаков дизурии или странгурии, но все еще демонстрировало полиурию и полидипсию. В лаборатории была подтверждена стойкая гипергликемия и глюкозурия (образцы натощак), хотя другие отклонения исчезли.

\begin{figure}[htbp]
\centering
\includegraphics[width=0.6\textwidth]{placeholder.jpg}
\caption{Осадок мочи пациента, 400x.}
\end{figure}

\subsection*{Вопрос: Какие изменения в анализе мочи?}
\subsection*{Вопрос: Учитывая результаты, каков предположительный диагноз?}
\subsection*{Вопрос: Какой лабораторный тест будет проведен для подтверждения?}

\subsection*{Разбор случая}
\textbf{Сахарный диабет и инфекция мочевыводящих путей}

\subsubsection*{Какие изменения в биохимии?}
Лабораторные тесты выявили кажущуюся гиперпротеинемию (из-за гиперглобулинемии), а также гипергликемию, гипертриглицеридемию и очень заметное повышение AST. Примечательно, что значения общего билирубина, щелочной фосфатазы и $\gamma$GT были равны 0. Основные причины гипергликемии перечислены в таблице 2.

\subsubsection*{Какие биохимические параметры могут быть недостоверны в этом образце?}
Учитывая выраженный гемолиз, наблюдаемый в образце плазмы (рис. 2), большинство изменений, обнаруженных в анализе, можно объяснить влиянием свободного гемоглобина на лабораторные методики. В таблице 3 показаны биохимические параметры, которые могут быть изменены гемолизом у мелких животных.

Гемолиз не мог объяснить повышенный уровень глюкозы у этого пациента. Поскольку животное было натощак, гипергликемия могла быть вызвана реакцией бегства во время забора пробы. Чтобы проверить/исключить этот вариант, рекомендуется повторить измерение через 30 минут, тщательно собрав образец без стресса. Для оценки глюкозы мы также должны убедиться, что животное правильно голодало (минимум 12 часов).

В этом клиническом случае после установки катетера и забора образца (после адекватного и тщательного голодания) все биохимические параметры были в пределах нормы, за исключением глюкозы (202 mg/dL).

\begin{table}[htbp]
\centering
\caption{Основные причины гипергликемии у мелких животных.}
\begin{tabular}{p{6cm}p{8cm}}
\toprule
Артефактные & Отсутствие голодания, страх или реакция бегства у кошек. \\
Фармакологические & Кортикостероиды, глюкозные сыворотки, медроксипрогестерон, кетамин или ксилазин и т.д. \\
Недостаточность/функциональность инсулина & Сахарный диабет. \\
Влияние других гормонов & Гиперадренокортицизм, акромегалия или гипертиреоз. \\
Другие причины & Сепсис, панкреатит или карциномы поджелудочной железы и т.д. \\
\bottomrule
\end{tabular}
\end{table}

\begin{table}[htbp]
\centering
\caption{Влияние разных степеней гемолиза на биохимические параметры у мелких животных. + Эквивалентно $100mg / dL$ свободного гемоглобина плазмы, $+++$ эквивалентно $>2,000mg / dL$. Влияние гемолиза может варьировать в зависимости от различного лабораторного оборудования.}
\begin{tabular}{cc}
\toprule
Степень гемолиза & \uparrow Артефактное \& \downarrow Артефактное \\
\midrule
+ & Fe \& Билирубин. \\
++ & Fe, AST, CK и LDH. \& Билирубин. \\
+++ & Fe, AST, CK, LDH, P и K. \& Билирубин и $\gamma$GT. \\
++++ & Fe, AST, CK, LDH, P, K, TP и триглицериды. \& Билирубин, $\gamma$GT и щелочная фосфатаза. \\
\bottomrule
\multicolumn{3}{l}{Fe = железо, AST = аспартатаминотрансфераза, CK = креатинкиназа, LDH = лактатдегидрогеназа, P = фосфор, K = калий, TP = общий белок.}
\end{tabular}
\end{table}

\subsubsection*{Какие изменения в анализе мочи?}
В осадке мочи была выраженная гематурия, лейкоцитурия и наличие кристаллов струвита (рис. 6). Как гематурия, так и лейкоцитурия считаются, когда насчитывается более пяти этих клеток в поле $400\times$. Учитывая, что образец был собран цистоцентезом и, следовательно, загрязнение при зондировании можно исключить, эта находка всегда будет аномальной. Наличие кристаллов струвита само по себе не обязательно совпадает с появлением камней, состоящих из этого минерала, и часто наблюдается у животных с щелочной мочой (тем более, если также имеется инфекция мочевыводящих путей). Камни струвита не были обнаружены рентгенографически у этого пациента.

Кроме того, плотность мочи была высокой, что можно объяснить наличием выраженной глюкозурии и протеинурии. Среди прочих причин, любая причина выраженной гипергликемии (тем более, если она сохраняется) может привести к глюкозурии, поскольку почечный порог фильтрации превышен. В таблице 4 перечислены основные подтипы протеинурии у мелких животных. В данном случае наиболее вероятна была постренальная протеинурия (из-за наличия активного осадка).

pH мочи был заметно щелочной, и была бактериурия (информация, полученная из посева).

\begin{table}[htbp]
\centering
\caption{Подтипы протеинурии у мелких животных.}
\begin{tabular}{p{5cm}p{9cm}}
\toprule
Преренальная протеинурия & Вызвана выраженным гемолизом или рабдомиолизом. Реже связана с судорогами или лихорадкой. \\
Ренальная протеинурия & Единственный тип протеинурии, обычно связанный с гипопротеинемией. Связан с гломерулонефритом и амилоидозом и т.д. \\
Постренальная протеинурия & Наиболее частая причина, в результате инфекций мочевыводящих путей (ИМП), цистита, пиелонефрита и новообразований и т.д. Связана с активным мочевым осадком (с заметной клеточностью). \\
\bottomrule
\end{tabular}
\end{table}

\begin{figure}[htbp]
\centering
\includegraphics[width=0.6\textwidth]{placeholder.jpg}
\caption{Вид эритроцитов (черные стрелки) и лейкоцитов (синие стрелки) в неокрашенном образце осадка мочи. Справа показаны несколько крупных кристаллов струвита с их типичной формой «крышки гроба», 400x.}
\end{figure}

\subsubsection*{Учитывая результаты, каков предположительный диагноз?}
Анализ мочи выявил инфекцию мочевыводящих путей. Этот процесс мог объяснить многие первоначальные симптомы у пациента. У этого животного дизурия разрешилась после антибактериального лечения без необходимости изменять минеральный состав мочи. Однако наличие глюкозурии натощак и гипергликемии должно вызвать подозрение на иное основное заболевание. С учетом последующего наблюдения у этого животного предположительный диагноз — сахарный диабет, поскольку была продемонстрирована гипергликемия натощак (и не связанная со стрессом) вместе с глюкозурией. Следует отметить, что инфекции мочевыводящих путей редки у кошек, хотя наличие глюкозурии может способствовать их появлению.

\subsubsection*{Какой лабораторный тест будет проведен для подтверждения?}
Клиническое подозрение может быть подтверждено в лаборатории путем определения уровня фруктозамина. Этот параметр информирует нас о наличии гипергликемии, сохраняющейся с течением времени, таким образом избегая эффектов, которые стресс может вызвать при единичных определениях.

\section{Случай 22. Сука с генерализованной лимфаденомегалией}
\subsection*{Клиническая история и физикальное обследование}
3-летняя сука немецкого дога (рис. 1) с прогрессирующим клиническим анамнезом (два месяца) рвоты, потери аппетита, полиурии/полидипсии, асцита и парапареза задних конечностей. Животное было правильно вакцинировано и дегельминтизировано, без предшествующей истории заболеваний. Физикальное обследование выявило генерализованное увеличение лимфатических узлов (подколенных, подмышечных, подлопаточных и поднижнечелюстных и т.д.), слегка бледные слизистые оболочки, адекватный статус гидратации и гиперестезию в сочетании с дефицитом периферического пульса на обеих задних конечностях. Перитонеальный выпот был классифицирован как чистый транссудат.

\begin{figure}[htbp]
\centering
\includegraphics[width=0.5\textwidth]{placeholder.jpg}
\caption{Пациентка.}
\end{figure}

\subsection*{Биохимия}
\begin{tabular}{lcc}
\toprule
Тест & Результат & Референсный диапазон \\
\midrule
HCT, \% & 34 & 37--54 \\
Общий белок, g/dL & 9.6 & 5.4--8.2 \\
Альбумин, g/dL & 1.9 & 2.5--4.4 \\
Глобулины, g/dL & 7.7 & 2.3--5.2 \\
Альбумин-глобулиновое соотношение & 0.2 & 0.48--1.91 \\
Общий билирубин, mg/dL & 0.2 & 0.1--0.6 \\
Креатинин, mg/dL & 3.1 & 0.3--1.3 \\
BUN, mg/dL & 32 & 7--25 \\
Соотношение BUN--креатинин & 10 & < 30 \\
Глюкоза, mg/dL & 81 & 60--110 \\
Триглицериды, mg/dL & 115 & 20--112 \\
Общий холестерин, mg/dL & 451 & 125--270 \\
Щелочная фосфатаза, IU/L & 89 & 20--150 \\
Аланинаминотрансфераза (ALT, GPT), IU/L & 189 & 10--118 \\
Аспартатаминотрансфераза (AST, GOT), IU/L & 74 & 14--45 \\
Гамма-глутамилтрансфераза ($\gamma$GT), IU/L & 5 & 0--7 \\
\bottomrule
\end{tabular}

\subsection*{Вопрос: Какие изменения в биохимии?}

\subsection*{Гемостатическое исследование}
\begin{tabular}{lcc}
\toprule
Тест & Результат & Референсный диапазон \\
\midrule
Прогромбиновое время & 12 секунд & 6--12 \\
Активированное частичное тромбопластиновое время & 25 секунд & 12--22 \\
Время кровотечения слизистой оболочки рта & 4 мин & < 4 \\
D-димеры & 400 ng/mL & < 250 \\
Тромбоциты & 150 k/uL & 200--500 \\
Антитромбин III & 20\% & 65--145 \\
\bottomrule
\end{tabular}

\subsection*{Вопрос: Какие изменения в гемостазе?}

\subsection*{Анализ мочи (цистоцентез)}
\begin{figure}[htbp]
\centering
\includegraphics[width=0.6\textwidth]{placeholder.jpg}
\caption{Осадок мочи пациентки, 100x.}
\end{figure}

\begin{tabular}{l}
\caption{Результаты анализа мочи.} \\
\toprule
Цвет и мутность \& Желтоватая и прозрачная. \\
Плотность мочи (рефрактометрия) \& 1,010 \\
Глюкоза (тест-полоска) \& - \\
Белки (тест-полоска) \& ++++ ($\geq$ 2,000 mg/dL) \\
UPC \& 5 \\
pH (тест-полоска) \& 7 \\
Бактерии (посев) \& Отрицательно \\
Кристаллы \& Не обнаружены \\
\bottomrule
\multicolumn{2}{l}{UPC: соотношение белок-креатинин в моче.}
\end{tabular}

\subsection*{Вопрос: Какие изменения в анализе мочи?}

\begin{figure}[htbp]
\centering
\includegraphics[width=0.8\textwidth]{placeholder.jpg}
\caption{Протеинограмма пациентки.}
\end{figure}

\subsection*{Лимфоузловая цитология}
\begin{figure}[htbp]
\centering
\includegraphics[width=0.6\textwidth]{placeholder.jpg}
\caption{Цитология подколенного лимфоузла, полученная тонкоигольной аспирационной пункцией. Окраска Diff-Quick, 400x.}
\end{figure}
\begin{figure}[htbp]
\centering
\includegraphics[width=0.6\textwidth]{placeholder.jpg}
\caption{Цитология подколенного лимфоузла, полученная тонкоигольной аспирационной пункцией. Окраска Diff-Quick, 1000x.}
\end{figure}

\subsection*{Вопрос: Какие изменения в протеинограмме пациентки?}
\subsection*{Вопрос: Учитывая полученные результаты, каков предположительный диагноз?}

\subsection*{Разбор случая}
\textbf{Нефротический синдром и лейшманиозная тромбоэмболия}

\subsubsection*{Какие изменения в биохимии?}
Была умеренная анемия (снижение HCT), тяжелая гиперпротеинемия за счет гиперглобулинемии наряду с гипоальбуминемией и снижение альбумин-глобулинового соотношения (селективная потеря альбумина). Помните, что основными причинами селективной гипоальбуминемии являются гломерулопатии, печеночная недостаточность, кишечная мальабсорбция и тяжелое недоедание (табл. 3 из случая 16).

Также наблюдались ренальная азотемия (без обезвоживания и с соотношением мочевина-креатинин, типичным для повреждения почечной паренхимы), а также печеночное повреждение (повышение трансаминаз), гипертриглицеридемия и гиперхолестеринемия. Примечание: макроскопическая липемия в образце не обнаружена, и связанных помех в измерениях не было.

\paragraph{На заметку}
Уровень сывороточного креатинина зависит от мышечной массы пациента (и, следовательно, от его размера, мышечного статуса, возраста и породы). Чтобы получить более объективное значение гломерулярной фильтрации в конкретных измерениях (особенно при выраженной потере веса или у щенков, гериатрических животных, гигантских или миниатюрных пород), рекомендуется определять уровень симметричного диметиларгинина (SDMA).

\subsubsection*{Какие изменения в гемостазе?}
Наиболее специфическим изменением было снижение уровня антитромбина III наряду с небольшим увеличением активированного частичного тромбопластинового времени, тромбоцитопенией и повышением D-димеров (все это индикаторы гиперфибринолиза). Антитромбин III является одним из основных антикоагулянтных белков, поэтому его потеря приводит к гиперкоагуляции и склонности к появлению тромбоэмболий в крупных сосудах (аорта или бедренная артерия и т.д.) с последующим повышением фибринолиза и изменениями в других тестах коагуляции потребления.

\subsubsection*{Какие изменения в анализе мочи?}
В осадке мочи этой пациентки обнаружены гематурия, лейкоцитурия и цилиндрурия (зернистые цилиндры) (рис. 2), а также тяжелая протеинурия. Учитывая, что другие причины протеинурии у этого животного можно было исключить (табл. 4 из случая 21), ее происхождение было классифицировано как чисто почечное. Значение протеинурии, скорректированное по креатинину в моче (UPC), было очень высоким у этой пациентки (табл. 1) и соответствовало картине нефропатии с потерей белка. Эта выраженная протеинурия ответственна за большую часть гипоальбуминемии у пациентки, хотя печеночная недостаточность также могла сыграть свою роль.

\paragraph{На заметку}
Цилиндры в осадке мочи — это агрегаты с тубулярной морфологией, которые обычно являются индикаторами заметного повреждения почек. Эти структуры быстро разрушаются, поэтому важно оценивать осадок быстро после получения образца мочи. Существуют разные типы цилиндров в зависимости от их состава (зернистые, клеточные, восковидные и т.д.), и их следует интерпретировать в сочетании с другими показателями повреждения почек.

\begin{table}[htbp]
\centering
\caption{Классификация протеинурии в соответствии с соотношением белок-креатинин в моче (UPC) у мелких животных.}
\begin{tabular}{lcc}
\toprule
Классификация & \multicolumn{2}{c}{UPC} \\
& Собака & Кошка \\
\midrule
Нет протеинурии & \multicolumn{2}{c}{< 0.2} \\
Пограничная & 0.2--0.5 & 0.2--0.4 \\
Протеинурия & > 0.5 & > 0.4 \\
\bottomrule
\multicolumn{3}{l}{Сначала исключите пре- и постренальные причины протеинурии.}
\end{tabular}
\end{table}

\subsubsection*{Какие изменения в протеинограмме?}
При сравнении протеинограммы пациентки с референсными значениями для собак (рис. 5 из случая 17) наблюдалась поликлональная гаммапатия. Это совместимо с инфекционной клинической картиной, вызывающей значительную продукцию антител, а также повышенные $\alpha$ и $\beta$ глобулины и даже выраженный пик $\beta - 2$ глобулинов.

Учитывая предыдущие результаты и типы белков, которые мигрируют в каждую область протеинограммы (табл. 2), повышение $\alpha$ и $\beta$ глобулинов может быть объяснено повышенными липопротеинами (HDL, LDL и VLDL), вызванными нефротическим синдромом. В свою очередь, отчетливый пик $\beta - 2$ глобулинов, скорее всего, был вызван ошибочным использованием плазмы, а не сыворотки для тестирования (протеинограммы должны выполняться и интерпретироваться с использованием образцов сыворотки, которые не содержат фибриноген, поскольку коагуляции не произошло).

\subsubsection*{Учитывая результаты, каков предположительный диагноз?}
Результаты выявили наличие нефротического синдрома с заметным повреждением почек и потерей антитромбина III (симптомы, совместимые с тромбоэмболией аорты). Эта нефропатия также вызывает повышение уровня липопротеинов (следовательно, также повышение триглицеридов и холестерина), что является частой находкой при нефротическом синдроме, отчасти из-за повышенной продукции липопротеинов печенью для баланса осмотического давления, а также из-за потери белка (что способствует липолизу). Наконец, также наблюдалась выраженная поликлональная гаммапатия и умеренное повреждение печени. Все эти эффекты появились у этой пациентки из-за повреждения, вызванного инфекцией \textit{Leishmania} spp., которая была диагностирована с помощью цитологии лимфоузла. Обратите внимание на большое количество плазматических клеток и макрофагов, присутствующих на рис. 4, а также на наличие многочисленных амастигот в цитоплазме макрофага в центре рис. 5.

\begin{table}[htbp]
\centering
\caption{Основные белки, которые мигрируют в каждой области протеинограммы собак, и значение их повышения. Примечание: различие между областями 1 (наиболее ранние миграции) и 2 (более поздние миграции) обычно специфично для лаборатории. APP: белок острой фазы, белок, повышение которого указывает на наличие острого неспецифического воспалительного или инфекционного состояния.}
\begin{tabular}{p{3cm}p{5cm}p{6cm}}
\toprule
Области протеинограммы & Присутствующие белки & Значение повышения \\
\midrule
\multirow{3}{*}{$\alpha$-1 глобулины} & Антитрипсин, PFA & \\
& HDL & Гиперхолестеринемия (нефротический синдром или эндокринопатия). \\
& Гаптоглобин и церулоплазмин. PFA & \\
\multirow{2}{*}{$\alpha$-2 глобулины} & VLDL & Гипертриглицеридемия и холестеринемия (нефротический синдром или эндокринопатия). \\
& Трансферрин и др. Сомнительно & \\
\multirow{4}{*}{$\beta$-1 глобулины} & Фибриноген & Ошибка образца (редко для APP). \\
& C-реактивный белок, PFA & \\
& LDL & Гиперхолестеринемия (нефротический синдром или эндокринопатия). \\
& Ig A или Ig M. & Заболевание печени, инфекционные процессы или новообразования. \\
$\gamma$ глобулины & Ig G & Инфекционные процессы или новообразования. \\
\bottomrule
\multicolumn{3}{l}{LDL: липопротеины низкой плотности; VLDL: липопротеины очень низкой плотности; HDL: липопротеины высокой плотности; Ig: иммуноглобулин.}
\end{tabular}
\end{table}

\section{Случай 23. Сука с желтушными слизистыми оболочками}
\subsection*{Клиническая история и физикальное обследование}
7-летняя сука смешанной породы с анамнезом желтухи (рис. 1), кровотечения из слизистых оболочек, вялости, анорексии и рвоты в течение двух месяцев. Животное было правильно вакцинировано и дегельминтизировано. Владельцы подозревали пищевое отравление. Направляющий ветеринар назначил лечение гепатопротекторами, строгую диету, способствующую работе печени, и мониторинг трансаминаз в течение одного месяца (табл. 1). Физикальное обследование не выявило признаков обезвоживания.

\begin{tabular}{@{}ll@{}}
\toprule
Ректальная температура & 37.9°C \\
ЧСС & 52 уд/мин \\
ЧД & 32 дых/мин \\
ВНК & Менее 2 сек \\
\bottomrule
\end{tabular}

\begin{figure}[htbp]
\centering
\includegraphics[width=0.5\textwidth]{placeholder.jpg}
\caption{Пациентка во время госпитализации.}
\end{figure}

\subsection*{Биохимия}
\begin{tabular}{lcc}
\toprule
Тест & Результат & Референсный диапазон \\
\midrule
HCT, \% & 41 & 37--54 \\
Общий белок, g/dL & 6.4 & 5.4--8.2 \\
Альбумин, g/dL & 2.7 & 2.5--4.4 \\
Глобулины, g/dL & 3.7 & 2.3--5.2 \\
Альбумин-глобулиновое соотношение & 0.7 & 0.48--1.91 \\
Общий билирубин, mg/dL & 2.3 & 0.1--0.6 \\
Креатинин, mg/dL & 0.6 & 0.3--1.3 \\
BUN, mg/dL & 8 & 7--25 \\
Глюкоза, mg/dL & 81 & 60--110 \\
Триглицериды, mg/dL & 44 & 20--112 \\
Общий холестерин, mg/dL & 139 & 125--270 \\
Щелочная фосфатаза, IU/L & 238 & 20--150 \\
Аланинаминотрансфераза (ALT, GPT), IU/L & 929 & 10--118 \\
Аспартатаминотрансфераза (AST, GOT), IU/L & 788 & 14--45 \\
Гамма-глутамилтрансфераза ($\gamma$GT), IU/L & 6 & 0--7 \\
\bottomrule
\end{tabular}

\subsection*{Вопрос: Какие изменения в биохимии?}

\begin{tabular}{lccccc}
\caption{Эволюция трансаминаз, измеренных у пациентки после того, как направивший ветеринар назначил гепатопротекторное лечение и диету, способствующую работе печени. Животное поступило в нашу клинику на 30-й день этого лечения, и представленные здесь данные биохимии соответствуют этому дню.} \\
\toprule
& Реф. диапазон & День 7 & День 14 & День 21 \\
\midrule
AP (IU/L) & 20--120 & 175 & 212 & 221 \\
ALT (IU/L) & 10--118 & 535 & 725 & 899 \\
AST (IU/L) & 14--45 & 49 & 96 & 126 \\
\bottomrule
\multicolumn{5}{l}{AP: щелочная фосфатаза; ALT: аланинаминотрансфераза; AST: аспартатаминотрансфераза.}
\end{tabular}

\subsection*{Вопрос: Судя по эволюции этих маркеров, считаете ли вы, что причина патологии заключалась в диете пациентки?}

\subsection*{УЗИ печени}
\begin{figure}[htbp]
\centering
\includegraphics[width=0.6\textwidth]{placeholder.jpg}
\caption{УЗИ печени пациентки (день 30, при поступлении).}
\end{figure}
\begin{figure}[htbp]
\centering
\includegraphics[width=0.6\textwidth]{placeholder.jpg}
\caption{Цитология под контролем УЗИ, полученная тонкоигольной пункцией узелка, указанного на рис. 2. Окраска Diff-Quik, 400x.}
\end{figure}

\subsection*{Вопрос: Учитывая результаты, каков предположительный диагноз?}

\subsection*{Разбор случая}
\subsubsection*{Какие изменения в биохимии?}
У пациентки было выраженное повышение билирубина и маркеров повреждения гепатоцитов (щелочная фосфатаза, ALT и AST). Повышения $\gamma$GT не наблюдалось, и не было изменений в показателях функции печени (табл. 2 из случая 19). Основные причины гипербилирубинемии перечислены в таблице 2.

\paragraph{На заметку}
Лучшими маркерами холестаза являются щелочная фосфатаза и $\gamma$GT. Первая также повышается при повреждении гепатоцитов (утечка ферментов). Если необходимо выбрать только один показатель холестаза, щелочная фосфатаза более информативна у собак, в то время как $\gamma$GT является лучшим выбором для кошек.

\begin{table}[htbp]
\centering
\caption{Дифференциал основных причин гипербилирубинемии у мелких животных.}
\begin{tabular}{p{4cm}p{6cm}p{4cm}}
\toprule
Тип гипербилирубинемии & Причины & Подтверждение диагноза \\
\midrule
Гемолитического происхождения & Аутоиммунная гемолитическая анемия, бабезиоз, тельца Гейнца и т.д. & Гемограмма или мазок крови и т.д. \\
Печеночного происхождения & Токсические, лекарственные, новообразования, хронический гепатит и т.д. & ALT, AST, AP и маркеры функции печени и т.д. \\
Холестатического происхождения & Холангит, билиарные новообразования, холелитиаз, мукоцеле, панкреатит и т.д. & AP, $\gamma$GT и УЗИ и т.д. \\
\bottomrule
\end{tabular}
\end{table}

\subsubsection*{Судя по эволюции маркеров, заключалась ли причина в диете?}
Учитывая динамику маркеров повреждения печени (рис. 4) и тот факт, что все уровни трансаминаз продолжали расти, а их значения не нормализовались, казалось маловероятным, что меры, принятые предыдущим ветеринаром, устранили причину повреждения печени. Обратите внимание, что хронический активный гепатит, вызванный отравленной диетой, — редкая возможность, но это происходит через процесс самоподдержания.

\begin{figure}[htbp]
\centering
\includegraphics[width=0.6\textwidth]{placeholder.jpg}
\caption{Типичная эволюция трансаминаз после одного острого эпизода повреждения печени.}
\end{figure}

\subsubsection*{Учитывая результаты, каков предположительный диагноз?}
Гипербилирубинемия и желтуха у этой пациентки, по-видимому, были связаны с выраженным повреждением печени и разрушением гепатоцитов. Цитология показала наличие многочисленных лимфобластов (в фоне, окружающем скопления гепатоцитов), как в узлах с аномальной эхогенностью, так и в окружающей печени (рис. 3). Окончательным диагнозом этой пациентки была лимфома печени. Примечательно, что не все опухоли печени вызывают повышение уровня ферментов (известное как «утечка»), и это еще менее вероятно в случаях хорошо капсулированных опухолей или при их прогрессирующем и медленном росте.

\section{Случай 24. Экстренный случай у собаки после наезда}
\subsection*{Клиническая история и физикальное обследование}
2-летняя собака породы пойнтер поступила экстренно после наезда 12 часов назад. Животное было дезориентировано и вялое, с беловатыми слизистыми оболочками.

\begin{tabular}{@{}ll@{}}
\toprule
Ректальная температура & 37.1°C \\
ЧСС & 170 уд/мин \\
ЧД & 48 дых/мин \\
ВНК & Более 2 секунд \\
\bottomrule
\end{tabular}

\subsection*{Биохимия}
\begin{tabular}{lcc}
\toprule
Тест & Результат & Референсный диапазон \\
\midrule
HCT, \% & 38 & 37--54 \\
Общий белок, g/dL & 3.9 & 5.4--8.2 \\
Альбумин, g/dL & 1.7 & 2.5--4.4 \\
Глобулины, g/dL & 2.2 & 2.3--5.2 \\
Альбумин-глобулиновое соотношение & 0.77 & 0.48--1.91 \\
Общий билирубин, mg/dL & 0.1 & 0.1--0.6 \\
Креатинин, mg/dL & 1.1 & 0.3--1.3 \\
BUN, mg/dL & 10 & 7--25 \\
Глюкоза, mg/dL & 97 & 60--110 \\
Креатинфосфокиназа (CK, CPK), IU/L & 4,500 & 20--200 \\
Щелочная фосфатаза, IU/L & 101 & 20--150 \\
Аланинаминотрансфераза (ALT, GPT), IU/L & 22 & 10--118 \\
Аспартатаминотрансфераза (AST, GOT), IU/L & 199 & 14--45 \\
Гамма-глутамилтрансфераза ($\gamma$GT), IU/L & 4 & 0--7 \\
Лактат, mmol/L & 3.5 & 0.60--2.90 \\
\bottomrule
\end{tabular}

\subsection*{Вопрос: Какие изменения в биохимии?}

\subsection*{Анализ мочи (цистоцентез)}
\begin{figure}[htbp]
\centering
\includegraphics[width=0.5\textwidth]{placeholder.jpg}
\caption{Моча была получена с помощью цистоцентеза под контролем УЗИ.}
\end{figure}

\begin{tabular}{l}
\caption{Результаты анализа мочи.} \\
\toprule
Цвет и мутность \& Коричневато-красноватая и прозрачная. \\
Плотность мочи (рефрактометрия) \& 1,035 \\
Белки (тест-полоска) \& +++ ($\leq$ 300 mg/dL) \\
Кровь (тест-полоска) \& +++ ($\leq$ 200 эритроцитов/мкл) \\
Эритроциты (осадок) \& Не обнаружены. \\
Другие клетки (осадок) \& Не обнаружены. \\
\bottomrule
\end{tabular}

\subsection*{Вопрос: Какие изменения в анализе мочи?}
\subsection*{Вопрос: Учитывая результаты, каков предположительный диагноз? Какой орган подвергается неминуемому риску значительного повреждения?}

\subsection*{Разбор случая}
\textbf{Рабдомиолиз и гиповолемический шок вследствие травмы}

\subsubsection*{Какие изменения в биохимии?}
Была подтверждена гипопротеинемия с гипоальбуминемией и гипоглобулинемией, обе, скорее всего, вызваны тяжелым острым кровотечением, хотя сокращение селезенки предотвратило снижение уровня HCT (табл. 4 из случая 2). Также было обнаружено тяжелое повышение маркеров мышечного повреждения (CK и AST). Повышение AST у этого пациента было вызвано травматической миопатией, а повреждение печени было исключено, поскольку уровень ALT был в пределах нормы (табл. 2).

Наконец, было заметное повышение уровня лактата, скорее всего, связанное с клинической картиной гиповолемического шока из-за острой кровопотери (табл. 3).

\paragraph{На заметку}
Основными причинами гиперлактатемии являются гиповолемия (тяжелое кровотечение и/или обезвоживание), гипоперфузия (прогрессирующая сердечная недостаточность или дистрибутивный шок и т.д.), гиповентиляция, очень высокий уровень физической нагрузки, судороги и, реже, заболевания печени или почечная недостаточность.

\begin{table}[htbp]
\centering
\caption{Непеченочные причины повышения основных трансаминаз и их интерпретация.}
\begin{tabular}{p{3cm}p{5cm}p{5cm}}
\toprule
Маркер & Непеченочные причины, которые могут вызвать \uparrow & Подтверждение диагноза \\
\midrule
ALT & Энтерит, панкреатит, перитонит, сахарный диабет или гипотиреоз и т.д. & • Более специфичный маркер печени. \\
& & • Одновременное повышение вместе с AST. \\
AST & Гемолиз, миопатия, энтерит, панкреатит, перитонит, сахарный диабет или гипотиреоз и т.д. & • Исключить гемолиз \textit{in vitro/in vivo}. \\
& & • Если миопатия \uparrow, CK также будет \uparrow. \\
& & • Если заболевание печени \uparrow, ALT будет \uparrow. \\
AP & Холестаз, глюкокортикоиды, остеопатии, панкреатит, сахарный диабет или гипотиреоз и т.д. & • Изучение других маркеров холестаза. \\
& & • Изучение изоформ. \\
\bottomrule
\end{tabular}
\end{table}

\begin{table}[htbp]
\centering
\caption{Градация гиповолемического шока у собак в зависимости от клинической картины.}
\begin{tabular}{lcccccc}
\toprule
Степень шока & ЧСС & ВНК & Периферический пульс & Цвет слизистых оболочек & Психический статус & Конечности \\
\midrule
Умеренный & 120--150 & Норма & Сильный & Розовый & Нормальный & Нормальные \\
Декомпенсированный & 150--170 & $\approx$ 2 сек. & Слабый & Бледный & Угнетенный & Нормальные--холодные \\
Тяжелый & 170--220 & > 2 сек. & Очень слабый & Очень бледный & Очень угнетенный & Холодные \\
\bottomrule
\multicolumn{7}{l}{Также учитывайте изменения артериального давления и диуреза и т.д.}
\end{tabular}
\end{table}

\subsubsection*{Какие изменения в анализе мочи?}
Окраска мочи указывала на выраженную пигментурию с дифференциальным диагнозом гематурия, гемоглобинурия или миоглобинурия. Учитывая биохимические данные и данные осадка, наиболее вероятной причиной у этого пациента была миоглобинурия (табл. 4). Наблюдаемая протеинурия была преренальной (табл. 4 из случая 21). Наконец, всегда важно проверять наличие эритроцитов и лейкоцитов в осадке, поскольку маркер «кровь» на тест-полосках мочи обнаруживает только гем-группу и, следовательно, не различает гематурию или гемоглобинурию из-за перекрестных реакций с миоглобинурией.

\begin{table}[htbp]
\centering
\caption{Дифференциал наиболее частых причин пигментурии у мелких животных. Обратите внимание, что эритроциты могут лизироваться (\textit{in vivo} или \textit{in vitro}) в моче, в результате чего некоторые гематурии теряют свои типичные характеристики.}
\begin{tabular}{p{3cm}p{2cm}p{2cm}p{2cm}p{2cm}p{2cm}}
\toprule
Причина пигментурии & Моча & Супернатант после центрифугирования & Осадок после центрифугирования & Эритроциты в осадке & Другие находки \\
\midrule
Гематурия & Красноватая & Прозрачный & Да & Да & - \\
Гемоглобинурия & Красноватая & Красноватый & Нет & Нет & Признаки гемолитической анемии \\
Миоглобинурия & Красная/ коричневая & Красноватый & Нет & Нет & Признаки миопатии (T CK) \\
\bottomrule
\end{tabular}
\end{table}

\subsubsection*{Учитывая результаты, каков предположительный диагноз? Какой орган в опасности?}
У животного были признаки тяжелого острого кровотечения, приведшего к гиповолемическому шоку, а также тяжелое мышечное повреждение (травматический рабдомиолиз), который привел к миоглобинурии.

Следует контролировать функцию почек, поскольку как гиповолемия, так и миоглобинурия являются особенно важными причинами повреждения этого органа.

\section{Случай 25. Собака с зловонной диареей}
\subsection*{Клиническая история и физикальное обследование}
6-летняя немецкая овчарка (рис. 1) с хроническим анамнезом (два месяца) большого объема зловонной диареи (рис. 2), а также спорадической рвоты. У животного был повышенный аппетит, но оно испытывало постепенную потерю веса. Пациент был вакцинирован и дегельминтизирован как внутренне, так и внешне. Физикальное обследование не выявило неврологических отклонений или признаков анемии, но было обнаружено $7\%$ обезвоживание.

\begin{figure}[htbp]
\centering
\includegraphics[width=0.5\textwidth]{placeholder.jpg}
\caption{Пациент во время госпитализации.}
\end{figure}

\begin{tabular}{@{}ll@{}}
\toprule
Ректальная температура & 38.2°C \\
ЧСС & 54 уд/мин \\
ЧД & 36 дых/мин \\
ВНК & Менее 2 сек \\
\bottomrule
\end{tabular}

\begin{figure}[htbp]
\centering
\includegraphics[width=0.5\textwidth]{placeholder.jpg}
\caption{Диарея, производимая животным.}
\end{figure}

\subsection*{Биохимия}
\begin{tabular}{lcc}
\toprule
Тест & Результат & Референсный диапазон \\
\midrule
HCT, \% & 55 & 37--54 \\
Общий белок, g/dL & 8.3 & 5.4--8.2 \\
Альбумин, g/dL & 4.7 & 2.5--4.4 \\
Глобулины, g/dL & 3.6 & 2.3--5.2 \\
Альбумин-глобулиновое соотношение & 1.3 & 0.48--1.91 \\
Общий билирубин, mg/dL & 0.2 & 0.1--0.6 \\
Креатинин, mg/dL & 1.2 & 0.3--1.3 \\
BUN, mg/dL & 12 & 7--25 \\
Глюкоза, mg/dL & 82 & 60--110 \\
Триглицериды, mg/dL & 101 & 20--112 \\
Общий холестерин, mg/dL & 219 & 125--275 \\
Амилаза, IU/L & 629 & 200--1,200 \\
Щелочная фосфатаза, IU/L & 25 & 20--150 \\
Аланинаминотрансфераза (ALT, GTP), IU/L & 20 & 10--118 \\
Аспартатаминотрансфераза (AST, GOT), IU/L & 28 & 14--45 \\
Гамма-глутамилтрансфераза ($\gamma$GT), IU/L & 3 & 0--7 \\
\bottomrule
\end{tabular}

\subsection*{Специфические биохимические тесты}
\begin{tabular}{lcc}
\toprule
Тест & Результат & Референсный диапазон \\
\midrule
Липаза (IU/L) & 125 & < 300 \\
TLI (ng/mL) & 0.9 & 9.6--31.2 \\
Кобаламин (ng/L) & 40 & 270--875 \\
Фолат ($\mu$g/dL) & 81 & 6.5--18.5 \\
cPLI (SNAP® +/-) & \multicolumn{2}{c}{Результаты в пределах референсного диапазона.} \\
Калий (mmol/L) & 3.4 & 3.7--5.8 \\
Натрий (mmol/L) & 134 & 138--160 \\
Хлор (mmol/L) & 114 & 106--120 \\
\bottomrule
\end{tabular}

\subsection*{Вопрос: Какие изменения в биохимических тестах?}
\subsection*{Вопрос: Учитывая результаты, каков предположительный диагноз?}
\subsection*{Вопрос: Какое лечение будет начато у этого пациента?}

\subsection*{Разбор случая}
\subsubsection*{Какие изменения в биохимических тестах?}
Учитывая повышение гематокрита, альбумина и общего белка, у этого животного в биохимии наблюдалось умеренное обезвоживание, что совпадало с данными клинической оценки.

Специфические биохимические тесты выявили снижение трипсиноподобной иммунореактивности (TLI), очень специфичную находку, связанную с экзокринной недостаточностью поджелудочной железы (ЭНПЖ). Кроме того, было исключено наличие панкреатита, поскольку амилаза, липаза и cPLI были в пределах нормы (см. биохимию из случая 18), и наблюдались гипокобаламинемия и повышенный уровень фолата в сыворотке. Изменения этих двух последних параметров также часто наблюдаются в случаях ЭНПЖ из-за неспособности животного абсорбировать кобаламин (что зависит от секреции внутреннего фактора поджелудочной железой) и вторичного бактериального роста (синдром избыточного бактериального роста в тонком кишечнике, SIBO), который повышает уровень фолата. Однако другие энтеральные процессы, поражающие определенные отделы кишечника, в которых всасываются эти витамины, или вызывающие SIBO, также могут изменять эти значения (табл. 1).

В ионограмме также наблюдались гипокалиемия и гипонатриемия, частые находки у пациентов с хронической диареей. Основные причины электролитных нарушений у мелких животных перечислены в таблице 2.

\begin{table}[htbp]
\centering
\caption{Изменения некоторых маркеров энтерально-панкреатических заболеваний в зависимости от типа и локализации патологии у мелких животных.}
\begin{tabular}{lccccc}
\toprule
& ЭНПЖ & Проксимальный энтерит & Дистальный энтерит & Диффузный энтерит & SIBO \\
\midrule
TLI & ↓ & Норма & Норма & Норма & Норма \\
Фолат & ↑ & ↓ & Норма & ↓ & ↑ \\
Кобаламин & ↓ & Норма & ↓ & ↓ & ↓ \\
\bottomrule
\multicolumn{6}{l}{TLI: трипсиноподобная иммунореактивность; SIBO: синдром избыточного бактериального роста в тонком кишечнике.}
\end{tabular}
\end{table}

\begin{table}[htbp]
\centering
\caption{Основные причины электролитных нарушений у мелких животных.}
\begin{tabular}{p{4cm}p{4cm}p{4cm}}
\toprule
& \multicolumn{2}{c}{Причины, которые ↑} & \multicolumn{2}{c}{Причины, которые ↓} \\
\midrule
Калий & Олигоанурия, разрыв мочевого пузыря, уроабдомен или болезнь Аддисона. & Рвота, диарея, полиурия и инсулинотерапия. \\
Натрий & Рвота, диарея, несахарный диабет или терминальная стадия почечной недостаточности (ESRD). & Рвота, диарея, терминальная стадия почечной недостаточности (ESRD), сердечная недостаточность (HF) или болезнь Аддисона. \\
Хлор & Диарея или почечная ацидурия. & Рвота или диуретики и т.д. \\
\bottomrule
\end{tabular}
\end{table}

\subsubsection*{Учитывая результаты, каков предположительный диагноз?}
Учитывая клиническую историю и признаки (обратите внимание на стеаторею, присутствующую на рис. 2) и лабораторные данные этого пациента, наиболее вероятным диагнозом была экзокринная недостаточность поджелудочной железы и легкое обезвоживание.

\subsubsection*{Какое лечение будет начато у этого пациента?}
В дополнение к инфузионной терапии с добавлением калия следует начать лечение панкреатическими ферментами для облегчения всасывания питательных веществ, а также добавки кобаламина и антибактериальную терапию для разрешения SIBO.

\chapter{Газометрические клинические случаи}
\label{chap:gas}

\section{Случай 26. Собака с хронической диареей}
\subsection*{Клиническая история и физикальное обследование}
5-летняя бельгийская овчарка с хроническим анамнезом (три недели) депрессии, потери аппетита и выраженной диареи, которой уже был поставлен диагноз (частичная обструкция инородным телом в дистальном отделе тонкого кишечника) и запланировано хирургическое вмешательство. Физикальное обследование выявило умеренное обезвоживание $(5\%)$.

\begin{tabular}{@{}ll@{}}
\toprule
Ректальная температура & 38.2°C \\
ЧСС & 50 уд/мин \\
ЧД & 32 дых/мин \\
ВНК & Менее 2 сек \\
\bottomrule
\end{tabular}

\subsection*{Венозная газометрия}
\begin{tabular}{lccccc}
\toprule
Тест & Результат & Реф. диапазон & LOW & NORMAL & HIGH \\
\midrule
pH & 7.28 & 7.30--7.48 & $\blacksquare$ \\
PCO2 & 25 mmHg & 27 -- 47 & $\blacksquare$ \\
PO2 & 33 mmHg & 24--48 & $\blacksquare$ \\
HCO3- & 16 mmol/L & 18--26 & $\blacksquare$ \\
TCO2 & 18 mmol/L & 19--28 & $\blacksquare$ \\
\bottomrule
\end{tabular}

\subsection*{Вопрос: Какие изменения в газометрии?}

\subsection*{Электролитная панель}
\begin{tabular}{lcc}
\toprule
Тест & Результат & Референсный диапазон \\
\midrule
Калий, mmol/L & 3.0 & 3.4--4.7 \\
Натрий, mmol/L & 136 & 142--152 \\
Хлор, mmol/L & 114 & 110--124 \\
Анионный интервал, mmol/L & 9 & 8--16 \\
\bottomrule
\end{tabular}

\subsection*{Вопрос: Какие изменения в электролитной панели?}
\subsection*{Вопрос: Учитывая результаты, какое кислотно-щелочное нарушение у этого животного? Какая жидкость считается идеальной для восстановления гидратации в этом случае?}

\subsection*{Разбор случая}
\textbf{Метаболический ацидоз с компенсаторным респираторным алкалозом}

\subsubsection*{Какие изменения в газометрии?}
Наблюдалось снижение pH (ацидемия, тенденция к ацидозу), наряду со снижением бикарбоната (тенденция к ацидозу в метаболическом компартменте) и снижением $\mathrm{PCO}_{2}$ (тенденция к алкалозу в респираторном компартменте). На рис. 1 показана интерпретация трех основных газометрических параметров: $\mathrm{pH}$, бикарбонат (метаболический компартмент) и $\mathrm{CO}_{2}$ или $\mathrm{PCO}_{2}$ парциальное давление (респираторный компартмент), с указанием того, когда изменение каждого из них вызывает тенденцию к ацидозу или алкалозу.

Поскольку изменение pH у этого пациента было вызвано снижением бикарбоната, можно сделать вывод, что у животного был первичный метаболический ацидоз. Если также наблюдается падение $\mathrm{PCO}_{2}$ (тенденция, противоположная pH), это будет компенсаторный респираторный алкалоз (рис. 1 и 2; табл. 1).

\paragraph{На заметку}
Ацидоз и алкалоз классифицируются как респираторные или метаболические в зависимости от компартмента, вызывающего изменение pH. В некоторых случаях (не всегда) противоположный компартмент может пытаться компенсировать pH (компенсаторный ответ), демонстрируя таким образом тенденцию, противоположную уровню pH. Даже когда этот компенсаторный ответ появляется, он обычно не может нормализовать pH. По определению, гиперкомпенсации не существует.

\begin{table}[htbp]
\centering
\caption{Возможные компенсаторные ответы при кислотно-щелочных проблемах, а также время, которое обычно требуется для их проявления, и степень ожидаемой компенсации (последняя ориентировочна, но плохо определена у кошек). Как время, так и степень компенсации сильно варьируют, когда первичное нарушение является респираторным (в зависимости от степени хронизации процесса).}
\begin{tabular}{p{3.5cm}p{3cm}p{3cm}p{3.5cm}}
\toprule
Первичная клиническая картина & Компенсаторная клиническая картина & Время, обычно необходимое для проявления компенсации & Степень компенсации* \\
\midrule
МЕТАБОЛИЧЕСКИЙ АЦИДОЗ & РЕСПИРАТОРНЫЙ АЛКАЛОЗ & 24 часа & ↓ PCO2 (0.7 mmHg на mmol/L ↓ бикарбоната). \\
МЕТАБОЛИЧЕСКИЙ АЛКАЛОЗ & РЕСПИРАТОРНЫЙ АЦИДОЗ & 24 часа & ↑ PCO2 (0.7 mmHg на mmol/L ↑ бикарбоната). \\
РЕСПИРАТОРНЫЙ АЦИДОЗ & МЕТАБОЛИЧЕСКИЙ АЛКАЛОЗ & До 30 дней & ↑ бикарбоната (0.15--0.55 mmol/L на mmHg ↑ в PCO2). \\
РЕСПИРАТОРНЫЙ АЛКАЛОЗ & МЕТАБОЛИЧЕСКИЙ АЦИДОЗ & До 30 дней & ↓ бикарбоната (0.25--0.55 mmol/L на mmHg ↓ в PCO2). \\
\bottomrule
\end{tabular}
\end{table}

\begin{figure}[htbp]
\centering
\includegraphics[width=0.8\textwidth]{placeholder.jpg}
\caption{Интерпретация изменений трех основных параметров кислотно-щелочного состояния. Обратите внимание, что изменения бикарбоната и $\mathrm{PCO}_{2}$ в одном направлении вызывают противоречивые тенденции.}
\end{figure}
\begin{figure}[htbp]
\centering
\includegraphics[width=0.8\textwidth]{placeholder.jpg}
\caption{Алгоритм интерпретации газометрического анализа. Начните с оценки изменений pH, а затем определите, какие изменения существуют в $\mathrm{PCO}_{2}$ и/или бикарбонате. На этом рисунке показаны только первичные изменения, а не возможные компенсаторные изменения.}
\end{figure}

Общий $\mathrm{CO}_{2}$ $\mathrm{(TCO_{2})}$ также был снижен. Мы должны учитывать, что этот параметр (предоставляемый некоторыми анализаторами) получается путем объединения концентраций $\mathrm{PCO}_{2}$ и бикарбоната. Поскольку бикарбонат по весу гораздо более распространен, чем $\mathrm{PCO}_{2}$, это значение обычно близко к концентрации бикарбоната.

Наконец, $\mathrm{PO}_{2}$ находился в пределах нормального референсного диапазона. Однако, учитывая, что это был венозный образец, следует проявлять осторожность при интерпретации этого значения.

\paragraph{На заметку}
Венозные газометрии не полезны для изучения оксигенации крови, поэтому оценка парциального давления $0_{2}$ или $\mathrm{PO}_{2}$ на основе этих тестов не показана. Однако эти типы образцов используются для определения кислотно-щелочных нарушений.

\subsubsection*{Какие изменения в электролитной панели?}
Были обнаружены гипокалиемия и гипонатриемия, обе, скорее всего, вызваны потерями ионов при диарее (табл. 2 из случая 25).

Анионный интервал был в пределах нормы, что позволило классифицировать метаболический ацидоз этого пациента как гиперхлоремический (табл. 2). Обратите внимание, что этот тип ацидоза называется гиперхлоремическим, даже когда хлор находится в пределах нормы (поскольку он относится к соотношению хлора и натрия и других ионов).

\begin{table}[htbp]
\centering
\caption{Использование анионного интервала в диагностике метаболического ацидоза, а также основные причины каждого подтипа метаболического ацидоза.}
\begin{tabular}{p{7cm}p{7cm}}
\toprule
\multicolumn{2}{c}{Анионный интервал (AG) — это формула, которая помогает различать две большие группы метаболического ацидоза.} \\
\multicolumn{2}{c}{Референсные значения: 7--17 у собак и 9--20 у кошек.} \\
\midrule
Метаболический ацидоз с AG в пределах нормы (также называемый ГИПЕРхлоремическим ацидозом) & Метаболический ацидоз с ↑ AG (также называемый НОРМОхлоремическим ацидозом) \\
\midrule
• Диарея. & • Лактат-ацидоз. \\
• Гипоадренокортицизм (болезнь Аддисона). & • Уремический ацидоз. \\
• Дилюционный ацидоз. & • Диабетический кетоацидоз. \\
• Введение хлорида аммония. & • Ацидоз из-за экзогенных токсинов (этиленгликоль или салицилат и т.д.). \\
• Почечный тубулярный ацидоз. & \\
\bottomrule
\end{tabular}
\end{table}

\subsubsection*{Учитывая результаты, какое кислотно-щелочное нарушение у животного? Какая жидкость идеальна?}
Метаболический ацидоз с анионным интервалом в пределах нормы типичен для пациентов с тяжелой диареей, а также в случаях гипокалиемии и гипонатриемии. Учитывая наблюдаемые изменения, регидратация у этого пациента должна основываться на ощелачивающей жидкости, которая обеспечивает как натрий, так и калий, что делает раствор Рингера лактата идеальным. Добавление бикарбоната натрия обычно предназначено для пациентов с более тяжелым снижением pH или бикарбоната ($< 7.2$ или $< 12$ мЭкв/л), хотя это решение должно приниматься клиницистом на основе эволюции и прогноза каждого клинического случая.

\section{Случай 27. Собака с неврологическими симптомами}
\subsection*{Клиническая история и физикальное обследование}
7-летний ирландский сеттер с тяжелой диареей, рвотой и легким обезвоживанием $(5\%)$. После первоначальной оценки была начата агрессивная инфузионная терапия раствором Рингера лактата с добавлением бикарбоната натрия. Через два часа у пациента начали развиваться клиническая картина судорог, тетания (рис. 1) и гипертермия $(40.2^{\circ}C)$. До начала этого лечения у животного не было никаких неврологических признаков.

\begin{tabular}{@{}ll@{}}
\toprule
Ректальная температура & 38.2°C \\
ЧСС & 54 уд/мин \\
ЧД & 36 дых/мин \\
ВНК & Менее 2 сек \\
\bottomrule
\end{tabular}

\begin{figure}[htbp]
\centering
\includegraphics[width=0.5\textwidth]{placeholder.jpg}
\caption{Пациент во время госпитализации.}
\end{figure}

\subsection*{Первоначальная венозная газометрия}
\begin{tabular}{lccccc}
\toprule
Тест & Результат & Реф. диапазон & LOW & NORMAL & HIGH \\
\midrule
pH & 7.21 & 7.30--7.48 & $\blacksquare$ \\
PCO2 & 25 mmHg & 27--47 & $\blacksquare$ \\
PO2 & 42 mmHg & 24--48 & $\blacksquare$ \\
HCO3- & 15 mmol/L & 18--26 & $\blacksquare$ \\
TCO2 & 17 mmol/L & 19--28 & $\blacksquare$ \\
\bottomrule
\end{tabular}

\subsection*{Венозные газометрии после судорог}
\begin{tabular}{lccccc}
\toprule
Тест & Результат & Реф. диапазон & LOW & NORMAL & HIGH \\
\midrule
pH & 7.4 & 7.30--7.48 & $\blacksquare$ \\
PCO2 & 32 mmHg & 27--47 & $\blacksquare$ \\
PO2 & 44 mmHg & 24--48 & $\blacksquare$ \\
HCO3- & 25 mmol/L & 18--26 & $\blacksquare$ \\
TCO2 & 27 mmol/L & 19--28 & $\blacksquare$ \\
\bottomrule
\end{tabular}

\subsection*{Вопрос: Какие изменения в газометриях?}

\subsection*{Эволюция биохимии и электролитной панели}
\begin{tabular}{lccc}
\toprule
Тест & Начальный результат & Результат после судорог & Референсный диапазон \\
\midrule
Альбумин, g/dL & 3.2 & 3.0 & 2.5--4.4 \\
Глюкоза, mg/dL & 100 & 81 & 60--110 \\
Аланинаминотрансфераза (ALT, GTP), IU/L & 16 & 17 & 10--118 \\
Калий, mmol/L & 4.4 & 4.2 & 3.4--4.7 \\
Натрий, mmol/L & 148 & 149 & 139--150 \\
Хлор, mmol/L & 119 & 114 & 110--120 \\
Анионный интервал, mmol/L & 16 & 13 & 7--17 \\
Общий кальций, mg/dL & 9.9 & 9.5 & 8.6--11.8 \\
\bottomrule
\end{tabular}

\subsection*{Вопрос: Может ли какой-либо из этих параметров объяснить неврологические симптомы? Какой аналит следует измерить для подтверждения предположительного диагноза?}
\subsection*{Вопрос: Какая профилактическая мера позволила бы избежать этой реакции на лечение?}

\subsection*{Разбор случая}
\textbf{Ионная гипокальциемия вследствие внезапного алкалоза}

\subsubsection*{Какие изменения в газометриях?}
В первоначальной газометрии обнаружены низкий pH, снижение бикарбоната (тенденция к ацидозу в метаболическом компартменте) и снижение $\mathrm{PCO_2}$ (тенденция к алкалозу в респираторном компартменте). Парциальное давление кислорода не следует оценивать (учитывая, что это венозный образец), а $\mathrm{TCO_2}$ показал тенденцию, сходную с бикарбонатом. Первоначальным диагнозом пациента будет первичный метаболический ацидоз с компенсаторным респираторным алкалозом. Поскольку анионный интервал был в пределах нормы, это состояние, скорее всего, было вызвано диареей (табл. 2 из случая 26).

Газометрический анализ после начала неврологических симптомов показал, что все значения (pH, бикарбонат и $\mathrm{PCO_2}$) быстро нормализовались, поэтому можно сделать вывод, что инфузионная терапия вызвала внезапный выраженный алкалоз.

\subsubsection*{Может ли какой-либо из этих параметров объяснить неврологические симптомы? Какой аналит измерить?}
Все собранные биохимические параметры были в пределах нормы как изначально, так и после начала судорог. Однако мы должны учитывать, что эволюция у этого пациента указывала на метаболическое происхождение неврологической клинической картины (внезапное начало после агрессивного лечения у взрослого животного без предшествующей истории судорог), поэтому причины этой симптоматики (табл. 1) следует исключать только с крайней осторожностью.

\begin{table}[htbp]
\centering
\caption{Основные метаболические причины судорог у мелких животных, а также возможные подтверждающие тесты.}
\begin{tabular}{p{6cm}p{8cm}}
\toprule
Метаболические причины судорог & Подтверждение диагноза \\
\midrule
Гипогликемия & Измерение уровня глюкозы крови натощак. \\
& Частая причина у самок в послеродовом периоде или при инсулиноме. \\
Печеночная недостаточность & Определение функции печени (с помощью функциональных тестов, а не только ферментов). \\
Ионная гипокальциемия & Необходимо измерить ионизированный кальций (не общий кальций). \\
& Подозрение на общую гипокальциемию или очень быструю алкализацию. \\
Тяжелая анемия (редко) & Выполнить общий анализ крови. Поддержка в клинике. \\
Другие менее частые причины (гиперлипопротеинемия, гипервязкость, гиперосмолярность или тяжелая уремия и т.д.) & Специфические биохимические определения, в зависимости от причины. \\
\bottomrule
\end{tabular}
\end{table}

Поскольку у этого животного произошла очень быстрая алкализация крови, судороги, вероятно, были вызваны ионной гипокальциемией, поэтому мы должны определить этот параметр. В данном клиническом случае концентрация ионизированного кальция после судорог составила 0.89 mmol/L (референсные диапазоны: 1.19--1.40 mmol/L для собак и 1.20--1.32 mmol/L у кошек).

\paragraph{На заметку}
Ионизированный кальций (активная форма) находится в равновесии с кальцием, связанным с белками (в основном альбумином), и солевыми формами. При ацидозе большая часть этого связанного с белком кальция высвобождается в ионной форме в процессе изменения, который обычно является постепенным. Если внутренняя среда очень быстро и агрессивно ощелачивается, происходит обратное (много ионизированного кальция связывается с белками, и, даже когда общее количество остается неизменным, могут появиться клинические признаки гипокальциемии из-за снижения фракции ионизированного кальция).

\subsubsection*{Какая профилактическая мера позволила бы избежать этой реакции?}
Чтобы избежать подобных проблем, мы должны постепенно изменять pH крови и избегать агрессивных методов лечения. Например, если для лечения метаболического ацидоза используется бикарбонат натрия, показано медленное введение только половины дозы. Затем кислотно-щелочной статус будет пересмотрен, чтобы изучить, следует ли вводить вторую половину, учитывая, что во многих случаях нормализация происходит без необходимости в полной дозе.

\section{Случай 28. Хроническая рвота у собаки}
\subsection*{Клиническая история и физикальное обследование}
9-летняя андалузская бодегера с предыдущим диагнозом лейшманиоза и хронической почечной недостаточности, у которой была рвота и признаки депрессии, потери аппетита и полиурия/полидипсия в течение двух месяцев. Животное прекратило лечение лейшманиоза год назад и соблюдало только почечную диету. Отмечено умеренное обезвоживание $(7\%)$.

\begin{tabular}{@{}ll@{}}
\toprule
Ректальная температура & 38.3°C \\
ЧСС & 70 уд/мин \\
ЧД & 44 дых/мин \\
ВНК & Менее 2 сек \\
\bottomrule
\end{tabular}

\subsection*{Венозная газометрия}
\begin{tabular}{lccccc}
\toprule
Тест & Результат & Реф. диапазон & LOW & NORMAL & HIGH \\
\midrule
pH & 7.51 & 7.30--7.48 & $\blacksquare$ \\
PCO2 & 52 mmHg & 27--47 & $\blacksquare$ \\
HCO3- & 33 mmol/L & 18--26 & $\blacksquare$ \\
\bottomrule
\end{tabular}

\subsection*{Вопрос: Какие изменения в газометрии?}

\subsection*{Биохимия и электролитная панель}
\begin{tabular}{lcc}
\toprule
Тест & Результат & Референсный диапазон \\
\midrule
Альбумин, g/dL & 2.2 & 2.5--4.4 \\
Общий белок, g/dL & 4.2 & 5.4--8.2 \\
BUN, mg/dL & 89 & 7--25 \\
Креатинин & 2.2 & 0.3--1.3 \\
Калий, mmol/L & 2.8 & 3.4--4.7 \\
Натрий, mmol/L & 149 & 139--150 \\
Хлор, mmol/L & 107 & 110--120 \\
Анионный интервал, mmol/L & 12 & 7--17 \\
\bottomrule
\end{tabular}

\subsection*{Вопрос: Какие изменения в биохимических тестах? Есть ли биохимический параметр, который объяснял бы кислотно-щелочные нарушения у этого пациента?}

\subsection*{Разбор случая}
\subsubsection*{Какие изменения в газометрии?}
Данные показали повышенный $\mathsf{pH}$, повышенный $\mathsf{PaCO}_{2}$ и повышенный бикарбонат (первичный респираторный ацидоз с компенсаторным метаболическим алкалозом). Однако исходное венозное PCO2 было повышено вместе с pH, что указывает на первичный метаболический алкалоз с компенсаторным респираторным ацидозом (поскольку PCO2 повышен). Согласно предоставленной газометрии: pH 7.51 (алкалемия), PCO2 52 (повышен, тенденция к ацидозу), HCO3- 33 (повышен, тенденция к алкалозу). Таким образом, первичное изменение — метаболический алкалоз (повышение бикарбоната), а повышение PCO2 является компенсаторным респираторным ацидозом. Это соответствует хронической рвоте.

\paragraph{На заметку}
Основные причины метаболического алкалоза у мелких животных:

\begin{tabular}{p{7cm}p{7cm}}
\toprule
ГИПОХЛОРЕМИЧЕСКИЙ метаболический алкалоз* (реагирует на введение Cl) & НОРМОХЛОРЕМИЧЕСКИЙ метаболический алкалоз* (не реагирует на введение Cl) \\
\midrule
• Рвота (или другие причины потери желудочного содержимого). & • Первичный гиперальдостеронизм (обычно опухоль). \\
• Использование петлевых диуретиков. & • Гиперадренокортицизм (болезнь Кушинга), хотя эффекты непостоянны. \\
• Введение щелочных продуктов (бикарбонат натрия, цитрат или, реже, гидроксид Al и т.д.). & \\
• Постгиперкапния. & \\
\bottomrule
\multicolumn{2}{l}{*Это название не обязательно подразумевает наличие гипо/нормохлоремии в анализе, и во многих случаях эта находка относительна к натриемии.}
\end{tabular}

\subsubsection*{Какие изменения в биохимических тестах? Есть ли объясняющий параметр?}
У животного была гипопротеинемия, гипоальбуминемия и азотемия (вероятно, из-за хронической почечной недостаточности). Присутствовали электролитные нарушения (гипохлоремия и гипокалиемия), типичные для пациентов с тяжелой рвотой (табл. 2 из случая 25). В этом клиническом случае мы пришли к выводу, что эта рвота была вторичной по отношению к уремии. Из этих параметров мы должны учитывать, что как белки (особенно альбумин), так и электролиты (в частности, соотношение хлор-натрий) влияют на кислотно-щелочной уровень.

Белки (в основном альбумин) являются слабыми кислотами, поэтому их снижение вызывает определенную тенденцию к алкалозу. В свою очередь, гипохлоремия и гипернатриемия (абсолютная или относительная друг к другу) приводят к тенденции к алкалозу, тогда как гиперхлоремия и гипонатриемия склонны приводить к ацидозу. Влияние этих параметров на кислотно-щелочной статус показано в таблице 2.

В этом клиническом случае гипопротеинемия (и гипоальбуминемия) и гипохлоремия (и относительная гипернатриемия) углубляли степень алкалоза (повышенная разница сильных ионов, сниженное соотношение хлор-натрий, повышенная разница натрий-хлор и низкий $A_{\mathrm{TOT}}$).

\begin{table}[htbp]
\centering
\caption{Формулы для расчета влияния белков и электролитов на кислотно-щелочные параметры.}
\begin{tabular}{p{14cm}}
\toprule
$A_{\mathrm{TOT}}$, или концентрация нелетучих буферов, определяет влияние белков на кислотно-щелочной статус. \\
$A_{\mathrm{TOT}} = 2.24 \times $ общий белок. \\
Референсные значения: 10--16 у собак и 8--17 у кошек* \\
↓ $A_{\mathrm{TOT}}$ = тенденция к алкалозу; ↑ $A_{\mathrm{TOT}}$ = тенденция к ацидозу. \\
\midrule
Разница сильных ионов (SID) оценивает влияние основных ионов на кислотно-щелочной статус. \\
SID = Na + K – Cl \\
Референсные значения: 34--40 у собак и 32--34 у кошек. \\
↓ SID = тенденция к ацидозу; ↑ SID = тенденция к алкалозу. \\
\midrule
Важность хлора и натрия в кислотно-щелочном статусе также можно оценить (даже когда оба иона находятся в пределах нормы), рассчитав следующие параметры: \\
• Соотношение хлор--натрий (референсные значения: 0.72--0.78 у собак и 0.74--0.8 у кошек). \\
• Разница натрий--хлор (референсные значения: 32--40 мЭкв/л у собак и кошек). \\
Высокое соотношение или малая разница: относительная гиперхлоремия (тенденция к ацидозу). \\
Сниженное соотношение или большая разница: относительная гипохлоремия (тенденция к алкалозу). \\
\bottomrule
\multicolumn{1}{l}{*Доступны более сложные формулы для расчета $\mathbb{A}_{\mathrm{TOT}}$, которые также учитывают концентрации альбумина, глобулинов и фосфора.}
\end{tabular}
\end{table}

\section{Случай 29. Кот с одышкой}
\subsection*{Клиническая история и физикальное обследование}
2-летний персидский кот, владельцы которого наблюдали постепенную потерю веса и апатию, и который в течение месяца стал привередливым в еде, с одышкой, тахипноэ и дыханием с открытым ртом, появившимися в последние дни. Животное никогда не было вакцинировано или дегельминтизировано. Физикальное обследование выявило слегка цианотичные слизистые оболочки и умеренное обезвоживание $(5\%)$, но была выраженная трудность при аускультации сердца и дыхательных шумов.

\begin{tabular}{@{}ll@{}}
\toprule
Ректальная температура & 37.7°C \\
ЧСС & 150 уд/мин \\
ЧД & Тяжелая одышка \\
ВНК & Менее 2 сек \\
\bottomrule
\end{tabular}

\subsection*{Артериальная газометрия}
\begin{tabular}{lccccc}
\toprule
Тест & Результат & Реф. диапазон & LOW & NORMAL & HIGH \\
\midrule
pH & 7.15 & 7.20--7.52 & $\blacksquare$ \\
PaCO2 & 59 mmHg & 22--50 & $\blacksquare$ \\
PaO2 & 80 mmHg & 90--110 & $\blacksquare$ \\
HCO3- & 27 mmol/L & 17--24 & $\blacksquare$ \\
\bottomrule
\end{tabular}

\subsection*{Вопрос: Какие изменения в газометрии?}

\subsection*{Рентгенография}
\begin{figure}[htbp]
\centering
\includegraphics[width=0.6\textwidth]{placeholder.jpg}
\caption{Латеро-латеральная рентгенограмма грудной клетки пациента.}
\end{figure}

\subsection*{Анализ плевральной жидкости}
\begin{tabular}{lc}
\toprule
Тест & Результат \\
\midrule
Цвет и внешний вид & Янтарная, слегка мутная \\
Общий белок, g/dL & 4.1 \\
Гематокрит (\%) & 2 \\
Общее количество ядерных клеток, на $\mu$L & 7,500 \\
\% Нейтрофилов & 78 \\
\bottomrule
\end{tabular}

\subsection*{Цитология плеврального выпота}
\begin{figure}[htbp]
\centering
\includegraphics[width=0.6\textwidth]{placeholder.jpg}
\caption{Цитология плеврального выпота. Окраска Diff-Quik, 40x.}
\end{figure}

\subsection*{Вопрос: Учитывая результаты, каков предположительный диагноз?}

\subsection*{Разбор случая}
\textbf{Респираторный ацидоз вследствие инфекционного перитонита кошек (FIP)}

\subsubsection*{Какие изменения в газометрии?}
Полученные данные показали снижение $\mathsf{pH}$, повышенный $\mathsf{PaCO}_{2}$ и повышенный бикарбонат (первичный респираторный ацидоз с компенсаторным метаболическим алкалозом). Парциальное артериальное давление кислорода было снижено (гипоксемия). В таблице 1 перечислены основные причины первичных респираторных кислотно-щелочных нарушений.

\begin{table}[htbp]
\centering
\caption{Основные причины респираторных кислотно-щелочных нарушений у мелких животных.}
\begin{tabular}{p{7cm}p{7cm}}
\toprule
\multicolumn{2}{c}{Первичный респираторный ацидоз обычно появляется при состояниях, связанных с гиповентиляцией; почти всегда присутствует сопутствующая гипоксемия.} \\
\midrule
• Обструкция дыхательных путей (коллапс трахеи, брахицефалический синдром, астма или аспирация и т.д.) & • Заболевания легких (пневмония, отек легких или астма и т.д.) \\
• Угнетение дыхательного центра (лекарства, анестетики, повреждение спинного или головного мозга и т.д.) & • Экстрапульмональные рестриктивные заболевания (диафрагмальные грыжи, пневмоторакс или плевральный выпот и т.д.) \\
• Нервно-мышечные заболевания (миастения гравис и т.д.) & \\
\midrule
\multicolumn{2}{c}{Первичный респираторный алкалоз возникает из-за гипервентиляции.} \\
\multicolumn{2}{c}{Респираторный алкалоз С гипоксемией} \\
\multicolumn{2}{c}{Респираторный алкалоз БЕЗ гипоксемии} \\
• Легочные патологии, нарушающие вентиляционно-перфузионное соотношение (тромбоэмболия легочной артерии, легочный фиброз или некоторые пневмонии и т.д.). & • Повреждения головного мозга в дыхательном центре (опухоли или травмы и т.д.). \\
• Сосудистые шунты в легочных сосудах. & • Стимуляция легочных рецепторов (некоторые пневмонии или тромбоэмболия легочной артерии и т.д.). \\
• Экологическая гипоксия (большие высоты). & • Экстремальная физическая нагрузка, тепловой удар или усиленное пыхтение и т.д. \\
• Тяжелая гипотензия или анемия. & \\
\bottomrule
\end{tabular}
\end{table}

\subsubsection*{Учитывая результаты, каков предположительный диагноз?}
Рентгенология выявила выраженный плевральный выпот, который был охарактеризован как воспалительный экссудат (табл. 2). Цитология показала результаты, типичные для асептического экссудата при инфекционном перитоните кошек (FIP); наблюдалось преобладание нейтрофилов (недегенерированных) и очень яркий, очень плотный белковый фон, который даже образовывал типичные «полумесяцы» (примечательно, что важно не путать точечные отложения белка с бактериями).

\begin{table}[htbp]
\centering
\caption{Основные типы полостных выпотов у мелких животных.}
\begin{tabular}{lcc}
\toprule
& Общий белок, g/dL & Общее количество ядерных клеток, на $\mu$L \\
\midrule
Транссудат & < 2.5 (обычно < 1.5) & < 1,500 \\
Модифицированный транссудат & > 2.5 & < 7,000 \\
Экссудат & > 2.5 (обычно > 3) & > 7,000 \\
\bottomrule
\end{tabular}
\end{table}

\section{Случай 30. Щенок с рвотой и диареей}
\subsection*{Клиническая история и физикальное обследование}
1-годовалая собака смешанной породы с тяжелой диареей и спорадической рвотой в течение последних двух недель. Пациент был очень угнетен и дышал с большим трудом за три дня до обращения. Животное было правильно вакцинировано и дегельминтизировано. Физикальное обследование выявило спутанность сознания и выраженную одышку смешанного типа, а также бледные и цианотичные слизистые оболочки и слабый периферический пульс.

\begin{tabular}{@{}ll@{}}
\toprule
Ректальная температура & 38.2°C \\
ЧСС & 150 уд/мин \\
ЧД & Выраженная одышка \\
ВНК & Немного более 2 секунд \\
\bottomrule
\end{tabular}

\subsection*{Артериальная газометрия}
\begin{tabular}{lccccc}
\toprule
Тест & Результат & Реф. диапазон & LOW & NORMAL & HIGH \\
\midrule
pH & 7.22 & 7.30--7.48 & $\blacksquare$ \\
PaCO2 & 55 mmHg & 27--47 & $\blacksquare$ \\
PaO2 & 69 mmHg & 85--100 & $\blacksquare$ \\
HCO3- & 14 mmol/L & 18--26 & $\blacksquare$ \\
\bottomrule
\end{tabular}

\subsection*{Вопрос: Какие изменения в газометрии?}

\subsection*{Биохимия и электролитная панель}
\begin{tabular}{lcc}
\toprule
Тест & Результат & Референсный диапазон \\
\midrule
Альбумин, g/dL & 4.5 & 2.5--4.4 \\
Общий белок, g/dL & 8.5 & 5.4--8.2 \\
Калий, mmol/L & 4 & 3.4--4.7 \\
Натрий, mmol/L & 147 & 139--150 \\
Хлор, mmol/L & 119 & 110--124 \\
Анионный интервал, mmol/L & 18 & 7--17 \\
$A_{\mathrm{TOT}}$, g/dL & 19 & 10--16 \\
SID, mmol/L & 32 & 34--40 \\
Лактат, mmol/L & 3.2 & 0.60--2.90 \\
\bottomrule
\end{tabular}

\subsection*{Вопрос: Какие изменения в биохимических тестах?}

\subsection*{Рентгенография}
\begin{figure}[htbp]
\centering
\includegraphics[width=0.6\textwidth]{placeholder.jpg}
\end{figure}

\subsection*{Вопрос: Учитывая результаты, каков предположительный диагноз?}

\subsection*{Разбор случая}
\textbf{Смешанный ацидоз вследствие диареи, гиповолемии и аспирационной пневмонии}

\subsubsection*{Какие изменения в газометрии?}
Наблюдалось падение pH (ацидемия, тенденция к ацидозу), наряду со снижением бикарбоната (тенденция к ацидозу в метаболическом компартменте) и повышением $\mathrm{PCO_2}$ (также тенденция к ацидозу в респираторном компартменте). Ни одна из этих клинических картин не была компенсаторной (все они вызывали одну и ту же тенденцию). $\mathrm{PO_2}$ был снижен (гипоксемия), что объясняло цианоз и указывало на то, что гиперкапния возникла из-за гиповентиляции.

Поскольку присутствовало более одной первичной клинической картины, у животного была смешанная клиническая картина как респираторного, так и метаболического ацидоза (табл. 1).

\paragraph{На заметку}
Мы говорим о смешанных кислотно-щелочных нарушениях, когда у пациента появляются два или более кислотно-щелочных изменения, не компенсирующих друг друга.

\begin{table}[htbp]
\centering
\caption{Как распознать смешанную кислотно-щелочную картину и некоторые примеры этих нарушений.}
\begin{tabular}{p{14cm}}
\toprule
Как распознать смешанное кислотно-щелочное нарушение \\
• PCO2 и бикарбонат изменяются в противоположных направлениях (↓/↑ или ↑/↓). \\
• pH в пределах нормы при PCO2 и/или бикарбонате вне диапазона*. \\
• Изменение pH, не соответствующее основному заболеванию пациента. \\
• В случаях метаболических нарушений оценивайте как анионный интервал, так и влияние $A_{\mathrm{TOT}}$ и SID. \\
Примеры смешанных кислотно-щелочных нарушений. \\
• Респираторный ацидоз и метаболический ацидоз. \\
• Респираторный алкалоз и метаболический алкалоз. \\
• Гиперхлоремический и нормохлоремический метаболический ацидоз. \\
• Гипохлоремический и нормохлоремический метаболический алкалоз. \\
• Метаболический ацидоз и алкалоз. \\
• Комбинации респираторной клинической картины с двумя метаболическими факторами и т.д. \\
Также учитывайте анамнез, клинические признаки и другие результаты тестов пациента, поскольку они могут указывать на то, что следует подозревать смешанные процессы (например, кот-диабетик с тяжелой рвотой после потери гликемического контроля [кетоацидоз]; pH в пределах нормы, но низкий хлор, высокий бикарбонат и высокий анионный интервал; диагноз: ацидоз и метаболический алкалоз). \\
\bottomrule
\multicolumn{1}{l}{*Помните, что компенсаторные клинические картины почти никогда не нормализуют pH.}
\end{tabular}
\end{table}

\subsubsection*{Какие изменения в биохимических тестах?}
У пациента был повышенный анионный интервал (из-за гиперлактатемии), а также низкий SID (относительная гиперхлоремия) и повышенный $\mathbf{A}_{\mathrm{TOT}}$ (гиперпротеинемия). Относительная гиперхлоремия и гипонатриемия также могут быть подтверждены расчетом соотношения хлор-натрий (0.81) или разницы между этими двумя ионами (28 mmol/L; табл. 2 из случая 28).

Таким образом, у этого пациента на момент обследования наблюдались три компонента метаболического ацидоза (гиперхлоремия, гиперпротеинемия и повышенный анионный интервал), и все они вызывали снижение бикарбоната и pH крови.

\subsubsection*{Учитывая результаты, каков предположительный диагноз?}
Рентгенографическая картина на уровне грудной клетки была совместима с клинической картиной аспирационной пневмонии (пневмонический очаг, сосредоточенный в гравитационно-зависимых областях; рис. 1), что могло объяснить появление респираторных симптомов после рвоты. Физикальное обследование (и гиперлактатемия) показало наличие гиповолемического шока из-за тяжелого обезвоживания (вероятно, из-за диареи и рвоты). Это обезвоживание также было ответственно за относительную гиперпротеинемию. Последующие тесты продемонстрировали наличие проксимальной кишечной инвагинации у этого пациента.

\vfill
\begin{center}
\includegraphics[width=0.6\textwidth]{placeholder.jpg}
\end{center}
\vspace{1cm}

\begin{center}
Гематологический анализ — один из наиболее широко используемых диагностических тестов в ветеринарных клиниках для мелких животных и входит в состав большинства рабочих протоколов. Это связано с тем, что он предоставляет широкий спектр информации об общем состоянии здоровья пациента и, следовательно, обеспечивает наиболее подходящие терапевтические меры в каждом конкретном случае. В других случаях гематологическая оценка помогает в диагностике различных заболеваний. Интерпретация гемограмм, сывороточной биохимии, анализов газов крови или мазков крови часто может представлять клиническую проблему, поэтому автор написал это руководство, чтобы проиллюстрировать различные клинические ситуации и способы их наилучшего разрешения.
\end{center}

\end{document}
